% Cada capítulo corresponde a una sesión completa del curso. Las subdivisiones
% a, b y c de los materiales se integran en un mismo capítulo.
\clearpage
\phantomsection
\addcontentsline{toc}{chapter}{\chaptertocunnumbered{Introducción}}
\markboth{Introducción}{Introducción}
\thispagestyle{plain}
\noindent{\tocdisplaytitle[0.9cm]{Introducción}\cftaftertoctitle}\par
\vspace{3em}
\section{La idea y el alcance del curso}

La teoría de esquemas es el lenguaje en el que la geometría algebraica moderna pone en correspondencia objetos geométricos y anillos conmutativos. La idea central, que guiará el curso, asocia a cada anillo $A$ su espectro primo $\Spec A$. Así, incluso $\Spec\mathbb Z$ se interpreta geométricamente y los problemas aritméticos pueden estudiarse con herramientas de geometría. Los puntos no serán únicamente soluciones sobre un cuerpo algebraicamente cerrado: serán todos los ideales primos, incluidos los puntos genéricos y los puntos que reflejan información aritmética.

La perspectiva aritmética estará presente desde el comienzo mediante ejemplos sobre $\mathbb Z$, cuerpos finitos, anillos locales y reducción módulo primos. El curso parte de los esquemas afines y su gavilla estructural; después construye esquemas generales por pegado y estudia sus morfismos y fibras. Más adelante aborda dimensión, regularidad, gavillas cuasicoherentes y coherentes, gavillas invertibles, divisores y la construcción $\Proj$. El programa concluye con curvas y ejemplos aritméticos. El programa académico distribuye estos temas en 24 sesiones y deja cohomología avanzada, planitud, geometría étale, descenso y stacks para un segundo curso.

El hilo conceptual inicial puede resumirse así: los ideales primos de $A$ son los puntos de $\Spec A$; la topología de Zariski describe sus especializaciones; las localizaciones $A_f$ y $A_{\mathfrak p}$ describen funciones en abiertos y gérmenes; la gavilla estructural reúne esa información local; y el pegado de abiertos afines produce esquemas que ya no tienen por qué ser afines.

\section{Lecturas y prerrequisitos}

Los textos principales indicados en el programa son \emph{The Rising Sea}, de Ravi Vakil, y \emph{Algebraic Geometry and Arithmetic Curves}, de Qing Liu. Como referencias complementarias aparecen Hartshorne, Görtz--Wedhorn, Mumford, Shafarevich, Atiyah--Macdonald, Eisenbud, el \emph{Stacks Project} y EGA. Para las primeras sesiones se espera familiaridad con ideales primos y maximales, localización, módulos, productos tensoriales, anillos noetherianos, dimensión de Krull, el lema de Nakayama y los teoremas de going-up y going-down; también se usa topología general básica. Las nociones de categorías, funtores, transformaciones naturales y propiedades universales se incorporan cuando hacen falta.

% \begin{nota}
% Este capítulo recoge la introducción temática del programa y de la presentación de la Sesión 0. Las indicaciones de inscripción, plataformas y evaluación pertenecen a la organización del curso y no forman parte del desarrollo matemático de los apuntes.
% \end{nota}


\chapter{Motivación y el conjunto \texorpdfstring{$\Spec A$}{Spec A}}
% Sesión 1.
\section{¿Por qué esquemas? De las variedades clásicas a los anillos}

La geometría algebraica clásica estudia conjuntos de soluciones de ecuaciones polinomiales sobre un cuerpo. Si $k$ es algebraicamente cerrado y $I\subseteq k[x_1,\ldots,x_n]$ es un ideal, los puntos de $V(I)$ corresponden a ideales maximales del anillo de coordenadas: $a=(a_1,\ldots,a_n)$ corresponde a $\mathfrak m_a=(x_1-a_1,\ldots,x_n-a_n)$.

\begin{teo}[Nullstellensatz de Hilbert (enunciado)]
\label{teo:s01-nullstellensatz-hilbert-enunciado}
La asignación $a=(a_1,\ldots,a_n)\mapsto\mathfrak m_a=(x_1-a_1,\ldots,x_n-a_n)$ induce una biyección entre $V(I)$ y los ideales maximales de $k[x_1,\ldots,x_n]/I$. Además, $f$ se anula en todo $V(I)$ si y solo si $f^m\in I$ para algún $m\ge1$.
\end{teo}

Esta correspondencia es el punto de partida, pero presenta cuatro limitaciones que motivan el paso a esquemas: sobre un cuerpo no algebraicamente cerrado algunas ecuaciones sin soluciones visibles esconden geometría (por ejemplo, $x^2+1=0$ sobre $\mathbb R$); los anillos aritméticos como $\mathbb Z$ no son cuerpos; el conjunto de puntos pierde multiplicidades e información infinitesimal; y una subvariedad irreducible, como una curva dentro de una superficie, todavía no aparece como un punto.

La propuesta de Grothendieck consiste en asociar un espacio geométrico a todo anillo conmutativo con unidad, no solo a anillos de polinomios sobre cuerpos. Sus puntos serán los ideales primos y sus funciones los elementos del anillo. La elección de primos, en vez de solo maximales, también garantiza la funtorialidad.

\section{El espectro primo como conjunto}

\begin{defi}[Espectro primo]
\label{defi:s01-espectro-primo}
Sea $A$ un anillo conmutativo con unidad. Su espectro primo es
\[
\Spec A=\set*{\mathfrak p\subseteq A \given \mathfrak p\text{ es un ideal primo}}.
\]
Escribiremos $\mathfrak p$ cuando pensemos algebraicamente y $x\in\Spec A$ cuando pensemos geométricamente; en este último caso $\mathfrak p_x$ denota el primo correspondiente.
\end{defi}

Recordemos que un ideal $\mathfrak p$ es primo si es propio y $fg\in\mathfrak p$ implica $f\in\mathfrak p$ o $g\in\mathfrak p$; equivalentemente, $A/\mathfrak p$ es un dominio entero.

\begin{lema}[Preimagen de un primo]
\label{lema:s01-preimagen-primo}
Si $\varphi:A\to B$ es un morfismo de anillos y $\mathfrak q\in\Spec B$, entonces $\varphi^{-1}(\mathfrak q)\in\Spec A$. Por tanto, $\varphi$ induce un mapa
\[
\Spec\varphi:\Spec B\longrightarrow\Spec A,\qquad \mathfrak q\longmapsto\varphi^{-1}(\mathfrak q).
\]
\end{lema}
\begin{proof}
El compuesto $A\xrightarrow{\varphi}B\to B/\mathfrak q$ tiene núcleo $\varphi^{-1}(\mathfrak q)$, así que induce una inyección $A/\varphi^{-1}(\mathfrak q)\hookrightarrow B/\mathfrak q$. Como un subanillo de un dominio es un dominio, el cociente de la izquierda es un dominio. Además, $\varphi^{-1}(\mathfrak q)$ es propio: si contuviera a $1$, entonces $1=\varphi(1)\in\mathfrak q$, contradicción. Luego es primo.
\end{proof}

\begin{obs}[Por qué primos y no maximales]
\label{obs:s01-primos-no-maximales}
Con maximales la construcción no funciona: la inclusión $\mathbb Z\hookrightarrow\mathbb Q$ lleva por preimagen el maximal $(0)\subset\mathbb Q$ al primo $(0)\subset\mathbb Z$, que no es maximal. En cambio, la preimagen de un primo siempre es primo. Los primos no maximales también serán necesarios para representar puntos genéricos.
\end{obs}

\begin{teo}[El espectro es vacío]
\label{teo:s01-espectro-vacio}
Se tiene $\Spec A=\varnothing$ si y solo si $A=0$.
\end{teo}
\begin{proof}
Si $A=0$, el ideal $(0)=A$ no es propio, así que no hay ideales primos. Supongamos ahora $A\ne0$. Considérese el conjunto $\Sigma$ de ideales propios de $A$, ordenado por inclusión. Es no vacío pues contiene a $(0)$. La unión de una cadena de ideales propios sigue siendo un ideal propio: si contuviera a $1$, alguno de los ideales de la cadena ya contendría a $1$. Por el lema de Zorn, $\Sigma$ tiene un elemento maximal, que es un ideal maximal de $A$ y, en particular, primo.
\end{proof}

\section{\texorpdfstring{Los elementos de $A$ como funciones sobre $\Spec A$}{Los elementos de A como funciones sobre Spec A}}

Para interpretar $f\in A$ como función en $\Spec A$, primero especificamos el cuerpo donde toma su valor en cada punto.

\begin{defi}[Cuerpo residual y evaluación]
\label{defi:s01-cuerpo-residual-evaluacion}
Para $\mathfrak p\in\Spec A$, definimos el cuerpo residual
\[
\kappa(\mathfrak p)=\Frac(A/\mathfrak p).
\]
Si $f\in A$, su valor $f(\mathfrak p)\in\kappa(\mathfrak p)$ es su imagen bajo $A\to A/\mathfrak p\hookrightarrow\kappa(\mathfrak p)$.
\end{defi}

\begin{prop}
\label{prop:s01-f-p-tiene-f-p}
Para $f\in A$ y $\mathfrak p\in\Spec A$, se tiene $f(\mathfrak p)=0$ si y solo si $f\in\mathfrak p$. El punto es exactamente el conjunto de funciones que se anulan en él.
\end{prop}
\begin{proof}
El valor es la imagen de $f$ en $A/\mathfrak p\hookrightarrow\kappa(\mathfrak p)$; como la segunda flecha es inyectiva, se anula exactamente cuando la clase de $f$ en $A/\mathfrak p$ es cero.
\end{proof}

El cuerpo de valores puede cambiar con el punto. Por ejemplo, en $\Spec\mathbb Z$, la función $15$ vale $1$ módulo $7$, se anula en $(3)$ y $(5)$, y en el punto genérico $(0)$ toma el valor $15\in\mathbb Q$.

En los cálculos de participación de la clase, $15$ toma el valor $4$ en $(11)$; la función $12$ vale $0$ en $(2)$ y $(3)$, y vale $2$ en $(5)$.

\begin{nota}[Advertencia: los valores no determinan la función]
Los valores puntuales no siempre determinan un elemento del anillo. En $A=k[\varepsilon]/(\varepsilon^2)$, el elemento $\varepsilon\ne0$ se anula en el único punto de $\Spec A$; véase el \cref{ejem:s01-punto-gordo}. La discrepancia se mide exactamente por el nilradical.
\end{nota}

\begin{teo}[Caracterización del nilradical]
\label{teo:s01-caracterizacion-nilradical}
Para $f\in A$ son equivalentes: (i) $f(\mathfrak p)=0$ para todo $\mathfrak p\in\Spec A$; (ii) $f$ es nilpotente. En particular,
\[
\Nil(A)=\set*{f\in A \given f^n=0\text{ para algún }n\ge1}=\bigcap_{\mathfrak p\in\Spec A}\mathfrak p.
\]
\end{teo}
\begin{proof}
Si $f^n=0$, entonces $f^n\in\mathfrak p$ para todo primo $\mathfrak p$, y la primalidad implica $f\in\mathfrak p$. Para la otra dirección probamos la contrapositiva. Supongamos que $f$ no es nilpotente y consideremos
\[
\Sigma=\set*{I\subseteq A \given I\text{ es ideal y }f^n\notin I\text{ para todo }n\ge1}.
\]
El conjunto contiene $(0)$. La unión de una cadena de elementos de $\Sigma$ vuelve a pertenecer a $\Sigma$, de modo que Zorn da un ideal maximal $\mathfrak p$ de $\Sigma$. Si $g,h\notin\mathfrak p$, la maximalidad implica que existen $m,n\ge1$ con $f^m\in\mathfrak p+(g)$ y $f^n\in\mathfrak p+(h)$. Al multiplicar estas inclusiones se obtiene $f^{m+n}\in\mathfrak p+(gh)$. Si $gh\in\mathfrak p$, resultaría $f^{m+n}\in\mathfrak p$, contradiciendo la definición de $\Sigma$. Por tanto $gh\notin\mathfrak p$, y $\mathfrak p$ es primo. Además $f\notin\mathfrak p$, así que $f$ no se anula en ese punto.
\end{proof}

\section{Galería de ejemplos}

\begin{ejem}[Cuerpos]
\label{ejem:s01-cuerpos}
Si $k$ es un cuerpo, $\Spec k=\{(0)\}$: un punto con cuerpo residual $k$. Así, $\Spec\mathbb C$ y $\Spec\mathbb F_2$ son el mismo espacio de un punto, aunque sus anillos de funciones sean distintos.
\end{ejem}

\begin{ejem}[El punto gordo]
\label{ejem:s01-punto-gordo}
Sea $A=k[\varepsilon]/(\varepsilon^2)$. Todo primo contiene a $\varepsilon$, pues $\varepsilon^2=0$, y $(\varepsilon)$ es maximal con cociente $k$. Por tanto $\Spec A=\{(\varepsilon)\}$. El punto tiene una sola posición, pero el anillo conserva el infinitesimal $\varepsilon$, que codifica información de primer orden y tangencia.
\end{ejem}

\begin{ejem}[El espectro de $\mathbb Z$]
\label{ejem:s01-espectro-z}
Los primos de $\mathbb Z$ son $(0)$ y $(p)$ para cada primo entero $p$. Entonces
\[
\Spec\mathbb Z=\{(0),(2),(3),(5),(7),\ldots\},\qquad
\kappa((0))=\mathbb Q,\quad \kappa((p))=\mathbb F_p.
\]
Se piensa como una curva con un punto genérico $(0)$ y puntos cerrados $(p)$.
\end{ejem}

\begin{prop}[La recta afín]
\label{prop:s01-recta-afin}
Si $k$ es un cuerpo,
\[
\Spec k[x]=\{(0)\}\cup\set*{(f) \given f\in k[x]\text{ irreducible mónico}}.
\]
Si $k$ es algebraicamente cerrado, los puntos cerrados son $(x-a)$, $a\in k$.
\end{prop}
\begin{proof}
El anillo $k[x]$ es un dominio de ideales principales. Todo primo no nulo es maximal y está generado por un polinomio irreducible; normalizando el generador, lo tomamos mónico. Si $k$ es algebraicamente cerrado, los irreducibles mónicos son precisamente los $x-a$.
\end{proof}

\begin{ejem}[La recta afín aritmética: $\mathbb Q$ y $\mathbb R$]
\label{ejem:s01-recta-afin-aritmetica-q-r}
Sobre cuerpos no algebraicamente cerrados aparecen puntos cerrados aritméticos. En $\Spec\mathbb Q[x]$, el punto $(x^2-2)$ tiene cuerpo residual $\mathbb Q(\sqrt2)$ y representa la pareja de raíces conjugadas $\{\sqrt2,-\sqrt2\}$. En $\Spec\mathbb R[x]$, el punto $(x^2+1)$ tiene cuerpo residual $\mathbb C$; representa una ecuación sin soluciones reales.
\end{ejem}

\begin{teo}[El plano afín]
\label{teo:s01-plano-afin}
Sea $k$ algebraicamente cerrado. Los ideales primos de $k[x,y]$ son exactamente $(0)$, los ideales $(f)$ con $f$ irreducible, y los maximales $(x-a,y-b)$ con $(a,b)\in k^2$.
\end{teo}
\begin{proof}
Los ideales $(0)$, $(f)$ para $f$ irreducible y $(x-a,y-b)$ son primos: los dos primeros porque $k[x,y]$ es un dominio de factorización única, y el cociente por $(x-a,y-b)$ es $k$. Falta mostrar que no hay otros. Sea $\mathfrak p$ un primo no nulo que no sea principal generado por un irreducible. Elijamos un elemento no nulo de $\mathfrak p$ y un factor irreducible $f$ suyo. Por primalidad $f\in\mathfrak p$. Como $\mathfrak p\ne(f)$, existe $g\in\mathfrak p\setminus(f)$; así $f$ y $g$ no tienen factor irreducible común.

Vistos en $k(x)[y]$, el máximo común divisor de $f$ y $g$ es $1$. En efecto, cualquier factor común no constante en $y$ se recuperaría, al limpiar contenido, como un factor común en $k[x,y]$ por el lema de Gauss. El anillo $k(x)[y]$ es euclidiano, así que existen $a,b\in k(x)[y]$ con $af+bg=1$. Al multiplicar por un denominador común no nulo $d\in k[x]$ obtenemos $uf+vg=d$ para ciertos $u,v\in k[x,y]$. Como $f,g\in\mathfrak p$, se sigue que $0\ne d\in\mathfrak p\cap k[x]$.

El polinomio $d$ se factoriza en factores lineales sobre $k$. Como $\mathfrak p$ es primo y $d\in\mathfrak p$, algún $x-a$ pertenece a $\mathfrak p$. En el cociente $k[x,y]/(x-a)\cong k[y]$, el ideal primo $\mathfrak p/(x-a)$ es $(0)$ o está generado por un polinomio irreducible. Puesto que $k$ es algebraicamente cerrado, en el segundo caso es $(y-b)$. El caso $(0)$ daría $\mathfrak p=(x-a)$, que es el primo principal excluido. Por tanto $\mathfrak p=(x-a,y-b)$.
\end{proof}
\begin{obs}[La jerarquía del plano afín]
\label{obs:s01-jerarquia-plano-afin}
La clasificación exhibe puntos de tres tamaños: puntos clásicos, puntos genéricos de curvas y el punto genérico del plano. Sobre cuerpos no algebraicamente cerrados, los maximales corresponden a órbitas de Galois, como en el \cref{ejem:s01-recta-afin-aritmetica-q-r}.
\end{obs}

\section{Puntos cerrados, puntos genéricos y especialización}

Para un ideal $I\subseteq A$ escribimos
\[
V(I)=\set*{\mathfrak p\in\Spec A \given I\subseteq\mathfrak p}.
\]
En la siguiente sesión se verá que estos conjuntos son los cerrados de la topología de Zariski y que la cerradura de $\mathfrak p$ es $V(\mathfrak p)$.

\begin{defi}[Especialización]
\label{defi:s01-especializacion}
Para $\mathfrak p,\mathfrak q\in\Spec A$, decimos que $\mathfrak q$ es una especialización de $\mathfrak p$, y escribimos $\mathfrak p\rightsquigarrow\mathfrak q$, si $\mathfrak p\subseteq\mathfrak q$. En este caso $\mathfrak p$ generiza a $\mathfrak q$.
\end{defi}

Un punto cerrado corresponde a un ideal maximal. El primo $\mathfrak p$ es el punto genérico del cerrado $V(\mathfrak p)$. La contención de ideales se invierte geométricamente: ideales más grandes corresponden a cerrados más pequeños. En $\Spec\mathbb Z$, $(0)\rightsquigarrow(p)$ para todo primo entero $p$.

\section{Ejercicios}

\begin{ejer}
\label{ejer:s01-espectro-producto-anillos}
Sea $A=A_1\times A_2$. Demuestra que todo primo de $A$ es de la forma $\mathfrak p_1\times A_2$ o $A_1\times\mathfrak p_2$, y concluye que $\Spec A\cong\Spec A_1\amalg\Spec A_2$. (Sugerencia: considera los idempotentes $e_1=(1,0)$ y $e_2=(0,1)$.)
\end{ejer}
\begin{ejer}
\label{ejer:s01-describe-explicitamente-indicando-puntos-cuerpos}
Describe explícitamente, indicando puntos y cuerpos residuales, los espectros de $\mathbb Q[x]/(x^2-2)$, $\mathbb R[x]/(x^2+1)$, $\mathbb C[x]/(x^2-1)$, $\mathbb C[x]/(x^2)$ y $\mathbb Z/(60)$.
\end{ejer}
\begin{ejer}
\label{ejer:s01-i-cociente-demuestra-i-inyectivo}
Sea $\varphi:A\twoheadrightarrow A/I$ el cociente. Demuestra que $\Spec(A/I)\to\Spec A$ es inyectivo con imagen $V(I)$.
\end{ejer}
\begin{ejer}
\label{ejer:s01-n-nilradical-red-n-usando}
Sea $N$ el nilradical de $A$ y $A_{\mathrm{red}}=A/N$. Usando el \cref{ejer:s01-i-cociente-demuestra-i-inyectivo} y el \cref{teo:s01-caracterizacion-nilradical}, demuestra que $\Spec A_{\mathrm{red}}\to\Spec A$ es una biyección.
\end{ejer}
\begin{ejer}
\label{ejer:s01-encuentra-todos-puntos-z-1}
Encuentra todos los puntos de $\Spec\mathbb Z[1/6]$. ¿Qué puntos de $\Spec\mathbb Z$ desaparecieron? Formula una conjetura sobre $\Spec A_f$.
\end{ejer}
\begin{ejer}
\label{ejer:s01-sobre-cuerpo-k-no-necesariamente}
Sobre un cuerpo $k$ no necesariamente algebraicamente cerrado, ¿qué enunciado sustituye a la lista de los primos de $k[x,y]$ del \cref{teo:s01-plano-afin}? Formúlalo para $k=\mathbb F_p$ y da un ejemplo de ideal maximal de $\mathbb F_2[x,y]$ que no sea $(x-a,y-b)$.
\end{ejer}

\section{Los dígitos son coeficientes de Taylor}

\begin{obs}[El punto gordo, releído]
\label{obs:s01-punto-gordo-releido}
El anillo del \cref{ejem:s01-punto-gordo}, $k[\varepsilon]/(\varepsilon^2)=k[x]/(x^2)$, es el truncamiento de Taylor de orden uno: guarda $a_0$ y $a_1$ y descarta los términos de orden superior. Que $\varepsilon\ne0$ se anule en el único punto significa que su término constante es cero; el valor puntual no ve el coeficiente infinitesimal de primer orden.
\end{obs}

\begin{obs}[Completar es la misma operación de los dos lados]
\label{obs:s01-completar-misma-operacion-dos-lados}
Guardar todos los coeficientes a la vez es completar:
\[
k[[x-a]]=\varprojlim_n k[x]/(x-a)^n,\qquad
\mathbb Z_p=\varprojlim_n\mathbb Z/p^n\mathbb Z.
\]
Los enteros $p$-ádicos son series de potencias en $p$. Esta analogía explica por qué $\Spec\mathbb Z$ se dibuja como una curva: cada punto tiene vecindades infinitesimales con la estructura de las de un punto de la recta afín.
\end{obs}

\begin{ejer}
\label{ejer:s01-calcula-primeros-cuatro-digitos-1}
Calcula los primeros cuatro dígitos de $-1$ en base $7$, es decir, su desarrollo como elemento de $\mathbb Z_7$. ¿Por qué el desarrollo no termina, si $-1$ es un entero? Compáralo con el desarrollo de $1/(1-x)$ en $k[[x]]$.
\end{ejer}

\section{Todo ideal de un producto es producto de ideales}

\begin{prop}
\label{prop:s01-k-ideal-b-existen-ideales}
Si $K$ es un ideal de $A\times B$, existen ideales $I\subseteq A$ y $J\subseteq B$ tales que $K=I\times J$.
\end{prop}
\begin{proof}
Definamos $I=\set*{a\in A \given (a,0)\in K}$ y $J=\set*{b\in B \given (0,b)\in K}$. Si $(a,b)\in K$, al multiplicar por los idempotentes $(1,0)$ y $(0,1)$ se obtiene $(a,0),(0,b)\in K$, así que $a\in I$ y $b\in J$. Esto prueba $K\subseteq I\times J$. Recíprocamente, si $a\in I$ y $b\in J$, entonces $(a,0),(0,b)\in K$ y su suma $(a,b)$ pertenece a $K$.
\end{proof}
\begin{obs}
\label{obs:s01-prueba-funciona-porque-idempotentes-1}
La prueba funciona porque los idempotentes $(1,0)$ y $(0,1)$ pertenecen al anillo y separan sus dos factores; son la razón algebraica de la descomposición.
\end{obs}

\section{Spec y productos: dónde importa que sean finitos}

\begin{teo}[La anti-equivalencia, sin restricciones]
\label{teo:s01-anti-equivalencia-sin-restricciones}
El funtor $\Spec:\mathbf{Ring}^{\mathrm{op}}\to\AffSch$ es una equivalencia. Por tanto, para cualquier familia finita o infinita $\{A_i\}_{i\in I}$, $\Spec(\prod_i A_i)$ es el coproducto de los $\Spec A_i$ en $\AffSch$.
\end{teo}

\begin{nota}[Advertencia]
La unión disjunta $\coprod_i\Spec A_i$ es el coproducto en $\Sch$. Con infinitos factores no nulos no es afín, pues no es cuasicompacta, así que no es un objeto de $\AffSch$. Lo que falla no es la anti-equivalencia, sino que la inclusión $\AffSch\hookrightarrow\Sch$ no preserva coproductos infinitos.
\end{nota}

\begin{prop}[La comparación, explícita]
\label{prop:s01-comparacion-explicita}
Sea $e_i\in\prod_jA_j$ el idempotente coordenado y $J=(e_i:i\in I)$. El morfismo canónico
\[
\Phi:\coprod_i\Spec A_i\longrightarrow\Spec\Bigl(\prod_i A_i\Bigr)
\]
es una inmersión abierta con imagen $\bigcup_iD(e_i)=\Spec(\prod_iA_i)\setminus V(J)$. Es un isomorfismo si y solo si $J=(1)$, si y solo si solo hay finitos $A_i\ne0$.
\end{prop}
\begin{proof}
Al localizar en $e_i$ se anulan todas las coordenadas salvo la $i$-ésima, así que $(\prod_jA_j)_{e_i}\cong A_i$ y $\Spec A_i$ se identifica con $D(e_i)$. Los idempotentes distintos son ortogonales, de modo que estos abiertos son disjuntos. La imagen es por tanto su unión, igual al complemento de $V(J)$. En un producto finito la suma de los idempotentes coordenados es $1$; en el producto infinito el ideal que generan consta de elementos de soporte finito y no contiene la unidad constante.
\end{proof}

\begin{obs}[Por qué la finitud]
\label{obs:s01-finitud}
La finitud expresa exactamente la igualdad $\sum_i e_i=1$: los idempotentes parten el espectro y para cubrirlo todo hace falta que sumen la unidad.
\end{obs}

\begin{obs}
\label{obs:s01-ejercicio-expuesto-clase-dos-factores}
El \cref{ejer:s01-espectro-producto-anillos} trata el caso de dos factores, donde $e_1+e_2=1$ y el morfismo sí es un isomorfismo; la \cref{prop:s01-comparacion-explicita} precisa el caso de familias arbitrarias.
\end{obs}

\section{Los dos puntos de \texorpdfstring{$\Spec \mathbb{C}[x]/(x^2-1)$}{Spec C[x]/(x2 - 1)}}

Por el teorema chino del resto, $\mathbb C[x]/(x^2-1)\cong\mathbb C\times\mathbb C$, de modo que hay dos puntos cerrados, ambos con cuerpo residual $\mathbb C$.

\begin{prop}[Los dos puntos no son intrínsecamente distinguibles]
\label{prop:s01-dos-puntos-no-son-intrinsecamente}
Existe un automorfismo de $\mathbb C$-álgebras que intercambia los dos puntos de $\Spec\mathbb C[x]/(x^2-1)$.
\end{prop}
\begin{proof}
El automorfismo inducido por $x\mapsto -x$ intercambia los ideales $(x-1)$ y $(x+1)$.
\end{proof}

\begin{obs}[Lo que sí quería decir]
\label{obs:s01-queria-decir}
El elemento $x$ vale $1$ en uno de los puntos y $-1$ en el otro. El conjunto de dos puntos, por sí solo, no recuerda si provino de $\mathbb C\times\mathbb C$, $\mathbb F_2\times\mathbb F_2$ u otro producto de dos cuerpos; la información está en las funciones y en dónde se evalúan.
\end{obs}


\chapter{La topología de Zariski}
% Esta sesión comienza con la topología de Zariski; la prueba del plano afín
% El ejemplo sobre primos del plano quedó incorporado en la Sesión 1.

\section{La topología de Zariski}

Ahora dotamos $\Spec A$ de una topología cuyos cerrados reflejen conjuntos de soluciones de ecuaciones.

\begin{defi}
\label{defi:s02-s-definimos-v-s-p}
Para $S\subseteq A$, definimos
\[
V(S)=\set*{\mathfrak p\in\Spec A \given S\subseteq\mathfrak p}.
\]
Como un primo contiene a $S$ si y solo si contiene al ideal que $S$ genera, $V(S)=V((S))$; escribiremos $V(I)$ para ideales $I$.
\end{defi}

\begin{obs}
\label{obs:s02-tiene-v-s-v-s}
Se tiene $V(S)=V((S))$, donde $(S)$ es el ideal generado; basta por ello considerar $V(I)$ para ideales. Además, si $I\subseteq J$, entonces $V(J)\subseteq V(I)$: más ecuaciones dan menos soluciones.
\end{obs}

\begin{prop}
\label{prop:s02-ideales-i-j-cualquier-familia}
Para ideales $I,J\subseteq A$ y cualquier familia $\{I_\alpha\}$,
\begin{enumerate}
\item $\displaystyle\bigcap_\alpha V(I_\alpha)=V\left(\sum_\alpha I_\alpha\right)$;
\item $V(I)\cup V(J)=V(I\cap J)=V(IJ)$.
\end{enumerate}
\end{prop}
\begin{proof}
La primera igualdad se sigue de que un ideal primo contiene la suma si y solo si contiene cada $I_\alpha$. Para la segunda, $IJ\subseteq I\cap J\subseteq I,J$ y la contravariancia dan $V(I)\cup V(J)\subseteq V(I\cap J)\subseteq V(IJ)$. Para la inclusión restante, supongamos $\mathfrak p\supseteq IJ$ y $\mathfrak p\not\supseteq I$. Elijamos $f\in I\setminus\mathfrak p$. Para todo $g\in J$, $fg\in IJ\subseteq\mathfrak p$; por primalidad, $g\in\mathfrak p$. Luego $J\subseteq\mathfrak p$ y $\mathfrak p\in V(J)$.
\end{proof}

En particular, $V((0))=\Spec A$ y $V(A)=\varnothing$. Las igualdades anteriores muestran que los $V(I)$ son los cerrados de una topología, la topología de Zariski.

\begin{ejem}[Los cerrados de $\Spec\mathbb Z$]
\label{ejem:s02-cerrados-z}
Por ejemplo, todo ideal de $\mathbb Z$ es principal y, para $n\ne0$, $V(n)=\set*{(p) \given p\mid n}$ es finito. Recíprocamente, cada conjunto finito de puntos cerrados se obtiene de esta forma, tomando el producto de los primos correspondientes. Así, los cerrados propios de $\Spec\mathbb Z$ son exactamente los conjuntos finitos de primos; los abiertos no vacíos son cofinitos y contienen $(0)$. Dos abiertos no vacíos se intersectan, de modo que este espacio no es Hausdorff.
\end{ejem}

\section{Cerradura de puntos, puntos genéricos y especialización}

\begin{prop}
\label{prop:s02-todo-p-cerradura-punto-p}
Para todo $\mathfrak p\in\Spec A$, la cerradura del punto es $\overline{\{\mathfrak p\}}=V(\mathfrak p)$. En particular, $\mathfrak p$ es cerrado si y solo si es maximal.
\end{prop}
\begin{proof}
Los cerrados que contienen a $\mathfrak p$ tienen la forma $V(I)$ con $I\subseteq\mathfrak p$. Por contravariancia, todos contienen a $V(\mathfrak p)$, que es entonces el cerrado mínimo que contiene al punto.
\end{proof}

\begin{defi}
\label{defi:s02-decimos-q-especializacion-p-q}
Decimos que $\mathfrak q$ es una especialización de $\mathfrak p$ si $\mathfrak q\in\overline{\{\mathfrak p\}}$, equivalentemente, si $\mathfrak p\subseteq\mathfrak q$. En este caso $\mathfrak p$ generiza a $\mathfrak q$ y es el punto genérico de $V(\mathfrak p)$.
\end{defi}

\begin{ejem}
\label{ejem:s02-z-0-p-todo-primo}
En $\Spec\mathbb Z$, $(0)\rightsquigarrow(p)$ para todo primo $p$ y $\overline{\{(0)\}}=\Spec\mathbb Z$. En $\Spec k[x,y]$, con $k$ algebraicamente cerrado, $(0)\rightsquigarrow(y-x^2)\rightsquigarrow(x-1,y-1)$; la inclusión de ideales registra la incidencia del plano, la parábola y el punto clásico.
\end{ejem}

\begin{obs}
\label{obs:s02-primos-no-maximales-registran-subvariedades}
Los primos no maximales registran subvariedades irreducibles mediante sus puntos genéricos; precisaremos este hecho en la Sesión 9.
\end{obs}

\section{Los abiertos distinguidos}

\begin{defi}
\label{defi:s02-f-abierto-distinguido-asociado-f}
Para $f\in A$, el abierto distinguido asociado a $f$ es
\[
D(f)=\Spec A\setminus V(f)=\set*{\mathfrak p \given f\notin\mathfrak p}.
\]
\end{defi}

\begin{prop}
\label{prop:s02-cumplen-siguientes-propiedades-d-f}
Se cumplen las siguientes propiedades:
\begin{enumerate}
\item $D(f)\cap D(g)=D(fg)$;
\item los $D(f)$ forman una base de la topología de Zariski;
\item $D(f)=\varnothing$ si y solo si $f$ es nilpotente, y $D(f)=\Spec A$ si y solo si $f$ es unidad.
\end{enumerate}
\end{prop}
\begin{proof}
La primera igualdad se desprende de la primalidad: $fg\notin\mathfrak p$ equivale a que $f,g\notin\mathfrak p$. Para la segunda, si $U=\Spec A\setminus V(I)$ es abierto, entonces $U=\bigcup_{f\in I}D(f)$. Para la tercera, $D(f)=\varnothing$ equivale a que $f$ pertenezca a todos los primos; por el \cref{teo:s01-caracterizacion-nilradical}, esto equivale a que $f$ sea nilpotente. Además, $D(f)=\Spec A$ significa que $f$ no está en ningún primo, en particular en ningún maximal; esto equivale a que $f$ sea unidad.
\end{proof}

\begin{obs}[Una promesa cumplida y otra sembrada]
\label{obs:s02-promesa-cumplida-otra-sembrada}
En la demostración del \cref{teo:s01-caracterizacion-nilradical} se construye, si $f$ no es nilpotente, un primo que evita todas sus potencias; ese primo pertenece a $D(f)$, así que $D(f)\ne\varnothing$. En el \cref{teo:s03-correspondencia-localizaciones} veremos que $D(f)$ es el espectro de un anillo localizado.
\end{obs}

\section{Cuasicompacidad}

\begin{prop}
\label{prop:s02-espacio-cuasicompacto-toda-cubierta-abierta}
El espacio $\Spec A$ es cuasicompacto: toda cubierta abierta admite una subcubierta finita.
\end{prop}
\begin{proof}
Basta refinar una cubierta abierta a una cubierta por básicos. Supongamos $\Spec A=\bigcup_{\alpha}D(f_\alpha)$. Entonces ningún ideal primo contiene el ideal generado por todos los $f_\alpha$, por lo que este ideal es $A$ (si fuera propio estaría contenido en un maximal). Así, $1$ es combinación lineal de una cantidad finita de los $f_\alpha$, digamos $1=\sum_{j=1}^n a_jf_{\alpha_j}$. Los mismos índices cubren: ningún primo puede contener simultáneamente esos $f_{\alpha_j}$. Por tanto $\Spec A=\bigcup_{j=1}^nD(f_{\alpha_j})$.
\end{proof}

\begin{obs}
\label{obs:s02-no-hipotesis-finitud-sobre-cuasicompacidad}
No hay hipótesis de finitud sobre $A$: la cuasicompacidad es un fenómeno algebraico que resulta de expresar $1$ con una combinación finita.
\end{obs}

\begin{obs}[Una tensión que vale la pena notar]
\label{obs:s02-tension-vale-pena-notar}
Una unión disjunta infinita de espacios no vacíos no es cuasicompacta, pero el espectro de un producto infinito de anillos sí lo es. Los abiertos disjuntos $D(e_i)$ de los idempotentes coordenados no cubren el espectro en general, porque el ideal que generan no contiene a $1$.
\end{obs}

\begin{ejer}[Sembrador: se cosecha en la Sesión 6]
\label{ejer:s02-sembrador-cosecha-sesion-6}
Sea $\{A_i\}_{i\in\mathbb N}$ una familia de anillos no nulos y $e_i\in\prod_jA_j$ el idempotente coordenado. Comprueba que $(\prod_jA_j)[1/e_i]\cong A_i$, que los $D(e_i)$ son disjuntos y que $\coprod_i\Spec A_i$ se inyecta en $\Spec(\prod_iA_i)$. Describe $J=(e_i:i\in\mathbb N)$ en términos de soportes, demuestra $J\ne(1)$ y explica en qué punto se usa que la familia sea infinita; compáralo con el caso finito.
\end{ejer}

\section{Ejercicios}

\begin{ejer}
\label{ejer:s02-demuestra-v-i-v-i}
Demuestra que $V(I)=V(\sqrt I)$ y que $V(I)\subseteq V(J)$ si y solo si $\sqrt J\subseteq\sqrt I$. Puedes usar $\sqrt I=\bigcap_{\mathfrak p\supseteq I}\mathfrak p$, que se deduce del \cref{teo:s01-caracterizacion-nilradical} aplicado a $A/I$.
\end{ejer}
\begin{ejer}
\label{ejer:s02-describe-todos-cerrados-z-k}
Describe todos los cerrados de $\Spec\mathbb Z$ y de $\Spec k[x]$, y determina qué subconjuntos son densos.
\end{ejer}
\begin{ejer}
\label{ejer:s02-demuestra-d-f-d-g}
Demuestra que $D(f)\subseteq D(g)$ si y solo si $f\in\sqrt{(g)}$.
\end{ejer}
\begin{ejer}
\label{ejer:s02-b-morfismo-anillos-demuestra-1}
Sea $\varphi:A\to B$ un morfismo de anillos. Demuestra que $(\Spec\varphi)^{-1}(V(I))=V(\varphi(I)B)$ y concluye que $\Spec\varphi$ es continua.
\end{ejer}
\begin{ejer}
\label{ejer:s02-espacio-no-vacio-irreducible-no}
Un espacio no vacío es irreducible si no es unión de dos cerrados propios. Demuestra que $V(\mathfrak p)$ es irreducible para todo primo $\mathfrak p$. ¿Es cierto el recíproco: todo cerrado irreducible es $V(\mathfrak p)$ para un único primo?
\end{ejer}



\chapter{Localización y funciones regulares}
\section{Dos promesas pendientes}

La Sesión 1 dejó abierto el cálculo de $\Spec\mathbb Z[1/6]$ y la Sesión 2 anunció que $D(f)$ debería ser el espectro de un anillo. La localización construye ese anillo. También introduce $A_{\mathfrak p}$, que describe el entorno algebraico de un punto.

\section{La localización de un anillo}

\begin{defi}[Localización]
\label{defi:s03-localizacion}
Un subconjunto $S\subseteq A$ es multiplicativo si $1\in S$ y es cerrado bajo productos. La localización $S^{-1}A$ consta de fracciones formales $a/s$, sujetas a
\[
\frac{a_1}{s_1}=\frac{a_2}{s_2}\quad\Longleftrightarrow\quad
\exists t\in S\text{ tal que }t(s_2a_1-s_1a_2)=0,
\]
con suma y producto como para fracciones ordinarias. El morfismo canónico es $\varphi:A\to S^{-1}A$, $a\mapsto a/1$.
\end{defi}

\begin{ejem}[Los dos sabores]
\label{ejem:s03-dos-sabores}
Las elecciones principales son $S=\{1,f,f^2,\ldots\}$, que da $A_f$, y $S=A\setminus\mathfrak p$, que da $A_{\mathfrak p}$. Si $A$ es dominio y $S=A\setminus\{0\}$, se obtiene $\Frac(A)$. Hay que distinguir $A_{\mathfrak p}$, que invierte lo que no pertenece a $\mathfrak p$, de $A_f$, que invierte las potencias de $f$.
\end{ejem}

\begin{obs}[El caso no conmutativo]
\label{obs:s03-caso-no-conmutativo}
La construcción usa que $A$ es conmutativo. En un anillo no conmutativo, para construir fracciones se requiere una condición de Ore sobre $S$; no toda multiplicativa admite una localización con las propiedades esperadas.
\end{obs}

\begin{prop}[Propiedad universal de la localización]
\label{prop:s03-propiedad-universal-localizacion}
El morfismo $\varphi:A\to S^{-1}A$ es inicial entre los morfismos desde $A$ que envían todos los elementos de $S$ a unidades. Es decir, si $\psi:A\to B$ y cada $\psi(s)$ es unidad, existe un único morfismo $\widetilde\psi:S^{-1}A\to B$ tal que $\widetilde\psi\circ\varphi=\psi$.
\end{prop}
\begin{proof}
La única posibilidad es $\widetilde\psi(a/s)=\psi(a)\psi(s)^{-1}$. Si $a/s=a'/s'$, hay $t\in S$ con $t(s'a-sa')=0$; al aplicar $\psi$ y multiplicar por las unidades correspondientes se verifica que las dos imágenes coinciden. La misma fórmula muestra que $\widetilde\psi$ respeta suma, producto y unidad. La unicidad se sigue de que todo $a/s$ es el producto de $a/1$ por $(s/1)^{-1}$.
\end{proof}

\begin{obs}[Epimorfismo, no necesariamente monomorfismo]
\label{obs:s03-epimorfismo-no-necesariamente-monomorfismo}
La localización $A\to S^{-1}A$ siempre es un epimorfismo en la categoría de anillos: si dos morfismos desde $S^{-1}A$ coinciden después de componer con $\varphi$, coinciden en $a/1$ y en $s/1$, por tanto también en $1/s$ y en todo $a/s$. No tiene por qué ser inyectivo ni sobreyectivo. Por ejemplo, $\mathbb Z\to\mathbb Q$ es epimorfismo pero no sobreyectivo. Un cociente $A\to A/I$ también es epimorfismo. En la categoría de anillos con unidad, los monomorfismos son precisamente los morfismos inyectivos.
\end{obs}

\begin{nota}[Advertencia]
La localización $A\to S^{-1}A$ no siempre es inyectiva: puede matar elementos si se invierten divisores de cero.
\end{nota}

\begin{ejem}[Localizar donde hay divisores de cero]
\label{ejem:s03-localizar-divisores-cero}
En $A=k[x,y]/(xy)$, al invertir $x$, la relación $xy=0$ fuerza $y=0$ y se obtiene $A_x\cong k[x]_x$. Geométricamente, $D(x)$ elimina el eje $y$ y el origen del eje $x$.
\end{ejem}

\section{El teorema de correspondencia}

\begin{teo}[Correspondencia para localizaciones]
\label{teo:s03-correspondencia-localizaciones}
Sea $S\subseteq A$ multiplicativo. Las asignaciones
\[
\mathfrak q\longmapsto\varphi^{-1}(\mathfrak q),\qquad
\mathfrak p\longmapsto S^{-1}\mathfrak p
\]
establecen una biyección entre $\Spec(S^{-1}A)$ y $\set*{\mathfrak p\in\Spec A \given \mathfrak p\cap S=\varnothing}$. La correspondencia inversa también preserva inclusiones.
\end{teo}
\begin{proof}
Si $\mathfrak q$ es primo de $S^{-1}A$, su contracción es primo y no corta a $S$, pues cada elemento de $S$ se vuelve unidad. Si $\mathfrak p$ es primo que no corta a $S$, su extensión $S^{-1}\mathfrak p$ es propia y primo: si $(a/s)(b/t)\in S^{-1}\mathfrak p$, algún $u\in S$ satisface $uab\in\mathfrak p$, y como $u\notin\mathfrak p$, se tiene $a\in\mathfrak p$ o $b\in\mathfrak p$. La extensión seguida de contracción devuelve $\mathfrak p$ porque $\mathfrak p$ es saturado respecto de $S$. En sentido contrario, todo ideal $J\subseteq S^{-1}A$ es extendido desde su contracción: si $a/s\in J$, entonces $a/1=(s/1)(a/s)\in J$; y si $a/1\in J$, entonces $a/s=(a/1)(s/1)^{-1}\in J$. Así, la contracción de la extensión de $\mathfrak q$ es $\mathfrak q$.
\end{proof}

\begin{coro}
\label{coro:s03-inclusion-inducida-identifica-f-d}
La inclusión inducida identifica $\Spec A_f$ con $D(f)$, de hecho homeomórficamente.
\end{coro}
\begin{proof}
Para las potencias de $f$, la condición $\mathfrak p\cap S=\varnothing$ equivale a $f\notin\mathfrak p$. La identificación respeta la topología: todo ideal de $A_f$ es extendido, de modo que los cerrados se corresponden con los cerrados relativos de $D(f)$.
\end{proof}

\begin{coro}
\label{coro:s03-s-p-p-identifica-q}
Para $S=A\setminus\mathfrak p$, $\Spec A_{\mathfrak p}$ se identifica con $\set*{\mathfrak q\in\Spec A \given \mathfrak q\subseteq\mathfrak p}$, exactamente los primos que generizan a $\mathfrak p$.
\end{coro}

\begin{obs}[El epimorfismo se vuelve monomorfismo]
\label{obs:s03-epimorfismo-vuelve-monomorfismo}
Por ejemplo,
\[
\Spec\mathbb Z[1/6]=D(6)=D(2)\cap D(3).
\]
Este espacio contiene $(0)$ y los $(p)$ con $p\ne2,3$; se eliminan $(2)$ y $(3)$. En general, $\Spec A_f$ elimina los primos que contienen $f$. El mapa inducido por $A\to S^{-1}A$ es inyectivo.
\end{obs}

\section{Anillos locales}

\begin{defi}
\label{defi:s03-anillo-conmutativo-unidad-local-tiene}
Un anillo conmutativo con unidad es local si tiene un único ideal maximal.
\end{defi}

\begin{prop}
\label{prop:s03-p-anillo-p-local-unico}
Para $\mathfrak p\in\Spec A$, el anillo $A_{\mathfrak p}$ es local y su único ideal maximal es $\mathfrak pA_{\mathfrak p}$.
\end{prop}
\begin{proof}
Por la correspondencia para localizaciones, los primos de $A_{\mathfrak p}$ corresponden a los primos $\mathfrak q\subseteq\mathfrak p$. El ideal extendido $\mathfrak pA_{\mathfrak p}$ es maximal, pues su cociente es el cuerpo $\Frac(A/\mathfrak p)$. Todo maximal de $A_{\mathfrak p}$ debe corresponder a un primo de $A$ contenido en $\mathfrak p$ maximal entre tales primos; necesariamente es $\mathfrak p$. Así el maximal es único.
\end{proof}

\begin{obs}[El cuerpo residual]
\label{obs:s03-cuerpo-residual}
El cuerpo residual del punto se identifica con $\kappa(\mathfrak p)$:
\[
A_{\mathfrak p}/\mathfrak pA_{\mathfrak p}\cong\Frac(A/\mathfrak p)=\kappa(\mathfrak p).
\]
\end{obs}

\begin{ejem}[Un germen de curva suave]
\label{ejem:s03-germen-curva-suave}
Sea $A=k[x]$ y $\mathfrak p=(x)$. Los primos de $A_{\mathfrak p}$ son $(0)\subset(x)A_{\mathfrak p}$: el germen local de la recta en el origen tiene dimensión uno.
\end{ejem}

\begin{ejem}[El mismo germen en teoría de números]
\label{ejem:s03-mismo-germen-teoria-numeros}
Para $A=\mathbb Z$ y $\mathfrak p=(p)$, los primos de $\mathbb Z_{(p)}$ son $(0)\subset(p)\mathbb Z_{(p)}$. Es el análogo aritmético del \cref{ejem:s03-germen-curva-suave}.
\end{ejem}

\begin{nota}[Advertencia: no confundir con los $p$-ádicos]
$\mathbb Z_{(p)}$ es una localización, mientras que $\mathbb Z_p$ es la completación $p$-ádica. Se relacionan mediante una inclusión densa, pero no son el mismo anillo. Del mismo modo, $\mathbb Q=\mathbb Z_{(0)}$ no es $\mathbb Q_p$.
\end{nota}

\section{Hacia las funciones regulares}

\begin{defi}[Función regular en un punto]
\label{defi:s03-funcion-regular-punto}
Sea $\mathfrak p\in\Spec A$. Un elemento de $A_{\mathfrak p}$ es una función regular en $\mathfrak p$: una fracción $a/s$ con $s\notin\mathfrak p$.
\end{defi}

\begin{obs}
\label{obs:s03-fraccion-s-tiene-sentido-todo}
La fracción $a/s$ tiene sentido en todo el abierto $D(s)\ni\mathfrak p$, con el mismo denominador para cada punto del abierto. Este comportamiento local es el que permitirá pegar información en una gavilla.
\end{obs}

El espectro de $\mathbb Z[1/p]$ es $D(p)$, todo $\Spec\mathbb Z$ salvo $(p)$; el de $\mathbb Z_{(p)}$ retiene solo $(0)$ y $(p)$. En el otro extremo, $\mathbb Z_{(0)}=\mathbb Q$ tiene un solo punto, el genérico.

\section{Ejercicios}

\begin{ejer}
\label{ejer:s03-describe-isomorfismo-f-t-tf}
Describe un isomorfismo $A_f\cong A[t]/(tf-1)$.
\end{ejer}
\begin{ejer}
\label{ejer:s03-demuestra-morfismo-s-1-inyectivo}
Demuestra que el morfismo $A\to S^{-1}A$ es inyectivo si y solo si cada elemento de $S$ es no divisor de cero en $A$.
\end{ejer}
\begin{ejer}
\label{ejer:s03-completa-ejemplo-demuestra-k-x}
Completa el ejemplo y demuestra que $(k[x,y]/(xy))_x\cong k[x]_x$, con $y\mapsto0$.
\end{ejer}
\begin{ejer}
\label{ejer:s03-demuestra-todo-ideal-j-s}
Demuestra que para todo ideal $J\subseteq S^{-1}A$ se tiene $J=(\varphi^{-1}(J))^e$. Deduce que la biyección del \cref{coro:s03-inclusion-inducida-identifica-f-d} es un homeomorfismo, comparando $V(I\cdot A_f)$ con $V(I)\cap D(f)$ para todo ideal $I\subseteq A$.
\end{ejer}
\begin{ejer}
\label{ejer:s03-anillos-z-2-z-2}
Los anillos $\mathbb Z_{(2)}$ y $\mathbb Z_2$ son anillos de valuación discreta con cuerpo residual $\mathbb F_2$ y espectros de la misma forma. ¿Son el mismo anillo? Da un elemento de $\mathbb Z_2$ que no esté en $\mathbb Z_{(2)}$ y explica por qué $\mathbb Z_{(2)}\subsetneq\mathbb Z_2$.
\end{ejer}
\begin{ejer}
\label{ejer:s03-sean-p-q-primos-construye}
Sean $\mathfrak p\subseteq\mathfrak q$ primos de $A$. Construye un morfismo natural $A_{\mathfrak q}\to A_{\mathfrak p}$ usando la propiedad universal de la localización.
\end{ejer}


\chapter{Pregavillas y gavillas}
\section{\texorpdfstring{De vuelta a “pegar los $A_f$”}{De vuelta a pegar los Af}}

La sesión retoma la observación de que cada anillo local $A_{\mathfrak p}$ ve funciones definidas en vecindades de $\mathfrak p$. Antes de formalizar el pegado, se completa el material sobre localizaciones.

\begin{defi}[Localización de un módulo]
\label{defi:s04-localizacion-modulo}
Si $M$ es un $A$-módulo y $S\subseteq A$ es multiplicativo, $S^{-1}M$ está formado por fracciones $m/s$, con la relación $m_1/s_1=m_2/s_2$ si existe $t\in S$ tal que $t(s_2m_1-s_1m_2)=0$. Es un módulo sobre $S^{-1}A$.
\end{defi}

\begin{lema}[Localizar es exacto]
\label{lema:s04-localizar-exacto}
La localización es exacta: si $0\to M'\to M\to M''\to0$ es exacta, también lo es $0\to S^{-1}M'\to S^{-1}M\to S^{-1}M''\to0$.
\end{lema}
\begin{proof}
La localización conmuta con cocientes: $S^{-1}(M/N)\cong S^{-1}M/S^{-1}N$, pues las clases de $m/s$ módulo $S^{-1}N$ son exactamente las fracciones de las clases de $m$ en $M/N$. Si $M'\to M$ es inyectivo, la clase de $m'/s$ que se anula en $S^{-1}M$ satisface $tm'=0$ en $M$ para algún $t\in S$, y por inyectividad también en $M'$. Así $S^{-1}M'\to S^{-1}M$ es inyectivo. Junto con la compatibilidad con cocientes, esto prueba la exactitud.
\end{proof}

\begin{lema}
\label{lema:s04-s-p-entonces-s-1}
Si $S=A\setminus\mathfrak p$, entonces $S^{-1}\mathfrak p=\mathfrak pA_{\mathfrak p}$ dentro de $A_{\mathfrak p}$.
\end{lema}
\begin{proof}
Ambos ideales consisten en las fracciones $a/s$ con $a\in\mathfrak p$ y $s\notin\mathfrak p$: al extender $\mathfrak p$ se invierten precisamente los elementos fuera de $\mathfrak p$.
\end{proof}

\begin{prop}
\label{prop:s04-todo-primo-p-p-pa}
Para todo primo $\mathfrak p$, $A_{\mathfrak p}/\mathfrak pA_{\mathfrak p}\cong\Frac(A/\mathfrak p)=\kappa(\mathfrak p)$.
\end{prop}
\begin{proof}
Al cocientar por $\mathfrak pA_{\mathfrak p}$ se obtiene la localización de $A/\mathfrak p$ en todas sus clases no nulas, que es su cuerpo de fracciones.
\end{proof}

\section{El ejemplo universal: funciones suaves}

El ejemplo guía son las funciones suaves: a cada abierto $U\subseteq X$ se asigna $\mathcal O(U)$, el conjunto de funciones suaves $U\to\mathbb R$, con restricciones por inclusión. Las secciones compatibles se pegan de forma única.

\begin{obs}
\label{obs:s04-argumento-no-usa-funciones-sean}
El argumento vale también para funciones continuas, holomorfas o reales. Todas estas familias comparten una estructura que formalizaremos enseguida.
\end{obs}

\begin{ejem}[El germen de una función suave y su tallo]
\label{ejem:s04-germen-funcion-suave-tallo}
Fijemos $p\in X$. Dos funciones suaves definidas cerca de $p$ representan el mismo germen si coinciden en alguna vecindad de $p$. El tallo $\mathcal O_p$ consta de pares $(f,U)$, con $p\in U$, módulo esa relación. Es un anillo local: la evaluación $[f]\mapsto f(p)$ tiene como núcleo el único ideal maximal de los gérmenes que se anulan en $p$.
\end{ejem}

\section{Pregavillas}

\begin{defi}[Pregavilla]
\label{defi:s04-pregavilla}
Una pregavilla de conjuntos $\mathcal F$ en $X$ asigna un conjunto $\mathcal F(U)$ a cada abierto $U\subseteq X$ y, a cada inclusión $U\subseteq V$, un mapa de restricción $\res_{V,U}:\mathcal F(V)\to\mathcal F(U)$. Se requiere que $\res_{U,U}=\id$ y que las restricciones sean compatibles con composición.
\end{defi}

\begin{obs}[Descripción categórica]
\label{obs:s04-descripcion-categorica}
Una pregavilla equivale a un funtor contravariante de la categoría de abiertos de $X$ a la categoría de conjuntos.
\end{obs}

\begin{ejem}[Pregavilla constante]
\label{ejem:s04-pregavilla-constante}
Para un conjunto $S$, la asignación $\mathcal F(U)=S$ con todas las restricciones iguales a la identidad es una pregavilla.
\end{ejem}

\section{Gavillas: los dos axiomas}

\begin{defi}[Gavilla]
\label{defi:s04-gavilla}
Una pregavilla $\mathcal F$ es una gavilla si para todo abierto $U$ y todo cubrimiento abierto $U=\bigcup_iU_i$ satisface:
\begin{enumerate}
\item \emph{Identidad:} si $s,t\in\mathcal F(U)$ restringen a la misma sección en cada $U_i$, entonces $s=t$.
\item \emph{Pegado:} si $s_i\in\mathcal F(U_i)$ y $s_i|_{U_i\cap U_j}=s_j|_{U_i\cap U_j}$ para todos $i,j$, existe una sección $s\in\mathcal F(U)$ que restringe a $s_i$ en cada $U_i$.
\end{enumerate}
\end{defi}

\begin{obs}
\label{obs:s04-identidad-dice-mas-manera-pegar}
Identidad dice que hay a lo más una manera de pegar y pegado que hay al menos una; juntos dan existencia y unicidad. Una pregavilla que satisface identidad se llama separada.
\end{obs}

\begin{ejem}[Las funciones suaves forman una gavilla]
\label{ejem:s04-funciones-suaves-forman-gavilla}
Las funciones suaves (y también las continuas u holomorfas) forman una gavilla: identidad y pegado se verifican localmente y permiten recuperar las secciones globales.
\end{ejem}

\begin{ejem}[La gavilla constante]
\label{ejem:s04-gavilla-constante}
La gavilla constante asociada a un conjunto discreto $S$ asigna las funciones localmente constantes con valores en $S$. En general difiere de la pregavilla constante, que no satisface pegado.
\end{ejem}

\begin{ejem}[La gavilla rascacielos]
\label{ejem:s04-gavilla-rascacielos}
Fijados $p\in X$ y un conjunto $S$, se define $i_{p,*}S(U)=S$ si $p\in U$ y $i_{p,*}S(U)=\{*\}$ si $p\notin U$.
\end{ejem}

\section{Tallos y gérmenes, en general}

\begin{defi}[Tallo]
\label{defi:s04-tallo}
Si $\mathcal F$ es una pregavilla y $p\in X$, el tallo es
\[
\mathcal F_p=\set*{(s,U) \given p\in U,\ s\in\mathcal F(U)}/\sim,
\]
donde $(s,U)\sim(t,V)$ si existe un abierto $W$ con $p\in W\subseteq U\cap V$ tal que $s|_W=t|_W$. La clase de $(s,U)$ es el germen de $s$ en $p$.
\end{defi}

\begin{obs}[La definición como colímite]
\label{obs:s04-definicion-colimite}
Equivalentemente, $\mathcal F_p=\colim_{U\ni p}\mathcal F(U)$ sobre las vecindades ordenadas por restricción. La categoría es filtrada porque la intersección de dos vecindades vuelve a ser una vecindad.
\end{obs}

\begin{nota}[Advertencia]
Para una gavilla general no hay necesariamente una evaluación a un conjunto fijo: en su definición no aparecen valores puntuales.
\end{nota}

\begin{obs}[Evaluación para el tallo de la gavilla estructural]
\label{obs:s04-evaluacion-tallo-gavilla-estructural}
En $\Spec A$ habrá un mapa $A_{\mathfrak p}\to\kappa(\mathfrak p)$; el cuerpo de valores depende del punto. Si $a/s=a'/s'$ en $A_{\mathfrak p}$, existe $u\notin\mathfrak p$ con $u(s'a-sa')=0$; al pasar a $A/\mathfrak p$ se puede cancelar $u$, así que la evaluación está bien definida.
\end{obs}

\section{Secciones determinadas por sus gérmenes}

\begin{teo}
\label{teo:s04-f-gavilla-todo-abierto-u}
En toda gavilla, una sección sobre $U$ queda determinada por sus gérmenes. En particular, para cada abierto $U$, el mapa siguiente es inyectivo:
\[
\mathcal F(U)\longrightarrow\prod_{p\in U}\mathcal F_p,\qquad s\longmapsto(s_p)_{p\in U}.
\]
\end{teo}
\begin{proof}
Si $s,t\in\mathcal F(U)$ tienen el mismo germen en cada punto, para cada $p$ hay una vecindad $W_p\subseteq U$ donde coinciden. Las $W_p$ cubren $U$, así que el axioma de identidad implica $s=t$. Solo se usa identidad, por lo que vale también para pregavillas separadas.
\end{proof}

\begin{obs}[El recíproco y el nombre “gavilla”]
\label{obs:s04-reciproco-nombre-gavilla}
La existencia de una sección que pegue gérmenes localmente compatibles requiere el axioma de pegado. Las secciones son tallos atados de manera compatible; de ahí viene el nombre.
\end{obs}

\section{Morfismos de (pre)gavillas}

\begin{defi}[Morfismo de pregavillas]
\label{defi:s04-morfismo-pregavillas}
Un morfismo $\varphi:\mathcal F\to\mathcal G$ en $X$ es una familia de mapas $\varphi(U):\mathcal F(U)\to\mathcal G(U)$ compatible con las restricciones.
\end{defi}

\begin{ejem}
\label{ejem:s04-inclusion-funciones-suaves-funciones-continuas}
La inclusión de funciones suaves en funciones continuas es un morfismo de gavillas.
\end{ejem}

\begin{prop}[Un morfismo induce mapas de tallos]
\label{prop:s04-morfismo-induce-mapas-tallos}
Todo morfismo $\varphi:\mathcal F\to\mathcal G$ induce para cada $p\in X$ un mapa $\varphi_p:\mathcal F_p\to\mathcal G_p$, $[s,U]\mapsto[\varphi(U)(s),U]$.
\end{prop}
\begin{proof}
La compatibilidad con las restricciones asegura que la imagen no depende del representante del germen.
\end{proof}

\section{Gavillas de anillos y de módulos}

\begin{defi}[Espacio anillado]
\label{defi:s04-espacio-anillado}
Un espacio anillado es un par $(X,\mathcal O_X)$ donde $\mathcal O_X$ es una gavilla de anillos, llamada gavilla estructural. Una gavilla $\mathcal F$ es un $\mathcal O_X$-módulo si cada $\mathcal F(U)$ es un $\mathcal O_X(U)$-módulo y las restricciones respetan la acción.
\end{defi}

\begin{obs}
\label{obs:s04-estructura-necesaria-construir-o-partir}
Esta es la estructura necesaria para construir $\mathcal O_{\Spec A}$ a partir de las localizaciones $A_f$ en los abiertos básicos.
\end{obs}

\section{Hacia la gavilla estructural}

La meta siguiente es construir $\mathcal O_{\Spec A}$ a partir de $A_f$ en los abiertos básicos $D(f)$, y verificar que dichos datos se pegan en una gavilla.
\[
A_{\mathfrak p}/\mathfrak pA_{\mathfrak p}\cong\Frac(A/\mathfrak p)=\kappa(\mathfrak p).
\]
La igualdad identifica el cuerpo residual del anillo local con el cuerpo residual ya asociado al punto en la Sesión 1.

\section{Ejercicios}

\begin{ejer}
\label{ejer:s04-ejemplo-funciones-suaves-demuestra-ideal}
En el \cref{ejem:s04-funciones-suaves-forman-gavilla}, demuestra que el ideal de gérmenes que se anulan en $p$ es el único ideal maximal del tallo $\mathcal O_p$.
\end{ejer}
\begin{ejer}
\label{ejer:s04-x-variedad-diferenciable-p-x}
Si $X$ es una variedad diferenciable y $p\in X$, identifica $\mathfrak m_p/\mathfrak m_p^2$ con el espacio cotangente $T_p^*X$.
\end{ejer}
\begin{ejer}
\label{ejer:s04-muestra-funciones-acotadas-r-funciones}
Muestra que las funciones acotadas en $\mathbb R$, y las funciones holomorfas en $\mathbb C$ que admiten una raíz cuadrada holomorfa local, forman pregavillas pero no gavillas.
\end{ejer}
\begin{ejer}
\label{ejer:s04-mapas-continuos-u-espacio-topologico}
Los mapas continuos $U\to Y$, para $Y$ un espacio topológico fijo, forman una gavilla de conjuntos en $X$. Si $Y$ es un grupo topológico, forman también una gavilla de grupos.
\end{ejer}
\begin{ejer}
\label{ejer:s04-nucleo-morfismo-pregavillas-definido-abierto}
El núcleo de un morfismo de pregavillas, definido abierto por abierto, es una pregavilla. Demuestra que si las dos pregavillas son gavillas, el núcleo también es una gavilla.
\end{ejer}
\begin{ejer}
\label{ejer:s04-sean-x-c-o-x}
Sean $X=\mathbb C$, $\mathcal O_X$ la gavilla de funciones holomorfas y $\mathcal O_X^*$ la gavilla de funciones holomorfas que no se anulan. El mapa $\exp:\mathcal O_X\to\mathcal O_X^*$, $f\mapsto e^{2\pi if}$, es sobreyectivo en cada tallo, pero no en cada abierto. Encuentra un abierto $U$ y una sección de $\mathcal O_X^*(U)$ que no esté en la imagen de $\mathcal O_X(U)$.
\end{ejer}
\begin{ejer}[Para pensar; no se demuestra en clase]
\label{ejer:s04-pensar-no-demuestra-clase}
Un morfismo de gavillas es un isomorfismo si y solo si lo es en todos los tallos. La afirmación es falsa para pregavillas generales.
\end{ejer}


\chapter{La gavilla estructural}
\section{De vuelta al problema pendiente}

En la Sesión 4 quedaron dos piezas sueltas. De la Sesión 3 sabemos qué anillo debe aparecer localmente en $\mathfrak p$, a saber $A_{\mathfrak p}$, y se anticipó que las secciones sobre $D(f)$ serían $A_f$, sin probarlo. Pero una gavilla está definida en todos los abiertos de $\Spec A$, no solo en los básicos; aunque $D(f)\cap D(g)=D(fg)$ sigue siendo básico, una unión arbitraria de básicos no suele serlo.

El problema de hoy es general: dada información compatible solo en una base de la topología, ¿cuándo y cómo se extiende a una gavilla en todos los abiertos? Resolverlo en general en las secciones 2 y 3 permite después verificar dos condiciones algebraicas para la familia $\{A_f\}_{f\in A}$.

\section{Gavillas sobre una base}

\begin{defi}[Pregavilla y gavilla sobre una base]
\label{defi:s05-pregavilla-gavilla-sobre-base}
Sea $\mathcal B$ una base de la topología de $X$. Una pregavilla sobre $\mathcal B$ asigna secciones a los elementos de $\mathcal B$ y mapas de restricción compatibles. Es una gavilla sobre la base si, siempre que $B=\bigcup_\alpha B_\alpha$ con todos los conjuntos en $\mathcal B$, satisface identidad y pegado. En la condición de pegado, la compatibilidad entre secciones se verifica en cada elemento básico contenido en $B_\alpha\cap B_\beta$, pues la intersección no necesariamente pertenece a la base.
\end{defi}

\section{El teorema de extensión}

\begin{teo}[Toda gavilla sobre una base se extiende, de manera única]
\label{teo:s05-toda-gavilla-sobre-base-extiende}
Toda gavilla sobre una base $\mathcal B$ se extiende a una gavilla $\mathcal F$ en $X$, de modo compatible con sus secciones y restricciones en $\mathcal B$. La extensión es única salvo isomorfismo único. El mismo resultado vale para gavillas con valores en anillos, grupos o módulos.
\end{teo}
\begin{proof}[Construcción y unicidad]
Para cada $p\in X$ se define el tallo formal $\mathcal F_p=\colim_{B\in\mathcal B,\,p\in B}\mathcal F(B)$. Para un abierto $U$, se toman las familias $(s_p)_{p\in U}\in\prod_{p\in U}\mathcal F_p$ que localmente provienen de una sola sección básica: para cada $p$ existe $B\in\mathcal B$ con $p\in B\subseteq U$ y $s\in\mathcal F(B)$ cuyos gérmenes en todo $q\in B$ son los prescritos $s_q$.

Las restricciones son por restricción de coordenadas. Identidad se verifica punto a punto. Para pegar secciones compatibles sobre un cubrimiento, se pegan sus coordenadas en cada tallo; la compatibilidad en las intersecciones garantiza que la sección resultante está localmente representada. En un abierto básico $B$, la aplicación que envía una sección original a sus gérmenes es biyectiva: identidad da inyectividad y los testigos básicos locales se pegan mediante el axioma de pegado sobre $\mathcal B$. Esto prueba que la construcción extiende la gavilla dada.

Si $\mathcal G$ es otra extensión, las identificaciones en cada abierto básico fuerzan los mismos tallos y, por tanto, las mismas secciones sobre todo abierto: dos secciones de una gavilla que coinciden en cada germen son iguales. También fuerzan el único isomorfismo compatible con las identificaciones iniciales.
\end{proof}

\begin{obs}
\label{obs:s05-teorema-3-1-vale-palabra}
El \cref{teo:s05-toda-gavilla-sobre-base-extiende} vale palabra por palabra para gavillas con valores en anillos, grupos, módulos, etcétera: las construcciones de colímite, producto y subconjunto definido por compatibilidad se realizan en la categoría correspondiente.
\end{obs}

\begin{ejer}[Morfismos de gavillas sobre una base]
\label{ejer:s05-morfismos-gavillas-sobre-base}
Un morfismo de gavillas en $X$ queda determinado por su restricción a una base. Además, todo morfismo entre gavillas sobre la base se extiende de manera única a un morfismo entre sus extensiones.
\end{ejer}

\section{De vuelta a Spec A: los básicos son base}

Los abiertos distinguidos $D(f)$ forman una base de $\Spec A$, pues
\[
\Spec A\setminus V(I)=\bigcup_{f\in I}D(f).
\]
La asignación candidata se define así.
\begin{defi}[La asignación candidata en los básicos]
\label{defi:s05-asignacion-candidata-basicos}
Para cada $f\in A$, sea
\[
S_f=\set*{g\in A \given D(f)\subseteq D(g)},\qquad
\mathcal O_{\Spec A}(D(f)):=S_f^{-1}A.
\]
Si $D(f')\subseteq D(f)$, entonces $S_f\subseteq S_{f'}$, lo que induce el mapa de restricción de localización $S_f^{-1}A\to S_{f'}^{-1}A$; estos mapas son transitivos.
\end{defi}

\section{\texorpdfstring{El teorema prometido: $A_f\cong\mathcal O_{\Spec A}(D(f))$}{El teorema prometido: Af es OSpec A(D(f))}}

\begin{teo}
\label{teo:s05-mapa-natural-f-s-f}
El mapa natural $A_f\to S_f^{-1}A=\mathcal O_{\Spec A}(D(f))$ es un isomorfismo.
\end{teo}
\begin{proof}
Como $D(f)\subseteq D(f)$, se tiene $f\in S_f$, así que la propiedad universal da $A_f\to S_f^{-1}A$. Por otra parte, si $g\in S_f$, entonces $D(f)\subseteq D(g)$, lo que equivale a $f\in\sqrt{(g)}$. Existe $n$ con $f^n=ag$; como $f$ es invertible en $A_f$, también lo es $g$. Por la propiedad universal hay un mapa inverso $S_f^{-1}A\to A_f$. Ambos compuestos son la identidad porque los mapas están determinados por su restricción desde $A$ y por el envío de los elementos invertidos a sus inversos.
\end{proof}

\section{\texorpdfstring{El corazón técnico: identidad y pegado para $\{A_f\}$}{El corazón técnico: identidad y pegado para Af}}

\begin{lema}[Cuasicompacidad de los básicos]
\label{lema:s05-cuasicompacidad-basicos}
Todo abierto distinguido $D(f)$ es cuasicompacto.
\end{lema}
\begin{proof}
Supongamos $D(f)\subseteq\bigcup_iD(f_i)$. Entonces ningún primo de $A_f$ contiene todos los $f_i/1$, por lo que el ideal que generan en $A_f$ es la unidad. Una combinación de $1$ usa finitos términos y, al limpiar denominadores, produce un cubrimiento finito de $D(f)$ por algunos $D(f_i)$.
\end{proof}

\begin{teo}[$\{A_f\}$ es gavilla sobre la base $\{D(f)\}$]
\label{teo:s05-f-gavilla-sobre-base-d}
La asignación $D(f)\mapsto A_f$ es una gavilla sobre la base de abiertos distinguidos.
\end{teo}
\begin{proof}
Por cuasicompacidad basta verificar identidad y pegado para un cubrimiento finito $D(f)=\bigcup_{i=1}^nD(f_i)$. Al reemplazar $A$ por $A_f$ podemos suponer $f=1$. Entonces $(f_1,\ldots,f_n)=A$, así que hay una identidad $1=\sum_i c_if_i$.

Para identidad, sea $s\in A$ cuya imagen se anula en cada $A_{f_i}$. Para cada $i$ existe $m_i$ con $f_i^{m_i}s=0$. Aumentando a un exponente común $m$, se tiene $f_i^ms=0$ para todo $i$. Las potencias $f_i^m$ generan el ideal unidad: basta elevar una expresión de $1$ como combinación de los $f_i$ a una potencia suficientemente grande y agrupar cada monomio, que contiene alguna potencia $f_i^m$. Así $1=\sum_i d_if_i^m$ y $s=\sum_i d_i(f_i^ms)=0$.

Para pegado, sean $s_i=a_i/f_i^k\in A_{f_i}$ secciones compatibles, usando un exponente común $k$. La compatibilidad en $A_{f_if_j}$ equivale a $a_if_j^k=a_jf_i^k$ en $A$. Como las $f_j^k$ generan la unidad, elijamos $e_j$ con $1=\sum_je_jf_j^k$ y pongamos $s=\sum_je_ja_j\in A$. Para cada $i$,
\[
f_i^ks=\sum_je_ja_jf_i^k
=a_i\sum_je_jf_j^k=a_i,
\]
por compatibilidad. Por tanto $s$ restringe a $s_i$ en $A_{f_i}$.
\end{proof}

\section{La gavilla estructural}

\begin{coro}[La gavilla estructural]
\label{coro:s05-gavilla-estructural}
Existe una única, salvo isomorfismo único, gavilla de anillos $\mathcal O_{\Spec A}$ en $\Spec A$ tal que $\mathcal O_{\Spec A}(D(f))\cong A_f$, compatible con restricciones. En particular, $\mathcal O_{\Spec A}(\Spec A)\cong A$, y el tallo en $\mathfrak p$ es $\mathcal O_{\Spec A,\mathfrak p}\cong A_{\mathfrak p}$.
\end{coro}
\begin{proof}
El resultado se sigue del \cref{teo:s05-toda-gavilla-sobre-base-extiende} y de la identificación $\mathcal O_{\Spec A}(D(f))\cong A_f$. Para el tallo, se toma el colímite de los $A_f$ sobre los básicos que contienen a $\mathfrak p$; es precisamente la localización que invierte todos los elementos fuera de $\mathfrak p$.
\end{proof}

\begin{obs}
\label{obs:s05-esto-cierra-toda-formalidad-circulo}
Esto cierra, con toda formalidad, el círculo abierto desde la Sesión 3: los $A_f$ son los valores de una gavilla genuina y $A_{\mathfrak p}$ es el tallo de esa gavilla en $\mathfrak p$.
\end{obs}

\section{Adelanto: gavillas de módulos}

El mismo argumento, sustituyendo $A$ por un $A$-módulo $M$, da la asignación $D(f)\mapsto S_f^{-1}M=:M_f$. Las dos verificaciones de identidad y pegado se obtienen de las del \cref{teo:s05-f-gavilla-sobre-base-d} cambiando $A$ por $M$. Por el \cref{teo:s05-toda-gavilla-sobre-base-extiende}, esta gavilla sobre la base se extiende a una gavilla $\widetilde M$ en $\Spec A$ con $\widetilde M(D(f))=M_f$. Como $M_f$ es un módulo sobre $A_f=\mathcal O_{\Spec A}(D(f))$, de manera compatible con restricciones, $\widetilde M$ es una gavilla de $\mathcal O_{\Spec A}$-módulos: el prototipo de gavilla cuasicoherente.

\section{Ejercicios}

\begin{ejer}[Cosechando el sembrador de la Sesión 4]
\label{ejer:s05-cosechando-sembrador-sesion-4}
Retomamos $X=\mathbb C$, $\mathcal O_X$ las funciones holomorfas y $\mathcal O_X^*$ las que no se anulan, con $\exp(f)=e^{2\pi i f}$. Usando la base de abiertos contráctiles, verifica que $\exp$ es sobreyectiva en cada básico (todo abierto simplemente conexo admite logaritmo holomorfo) y concluye —sin repetir el contraejemplo concreto de la Sesión 4— que un morfismo de gavillas sobreyectivo en cada elemento de una base induce un epimorfismo de gavillas, pero el recíproco es falso.
\end{ejer}
\begin{ejer}
\label{ejer:s05-adapta-demostracion-b1-teorema-6}
Adapta la demostración de (B1) en el \cref{teo:s05-f-gavilla-sobre-base-d} al caso de un $D(f)$ arbitrario cubierto por básicos $D(f_i)$ con $f_i\in S_f$.
\end{ejer}
\begin{ejer}
\label{ejer:s05-adapta-demostracion-b2-teorema-6}
Adapta la demostración de (B2) en el \cref{teo:s05-f-gavilla-sobre-base-d} al caso de un $D(f)$ arbitrario cubierto por básicos $D(f_i)$ con $f_i\in S_f$.
\end{ejer}
\begin{ejer}
\label{ejer:s05-escribe-b1-b2-teorema-6}
Escribe (B1) y (B2) del \cref{teo:s05-f-gavilla-sobre-base-d} como la exactitud de $0\to A\to\prod_{i=1}^n A_{f_i}\to\prod_{i<j}A_{f_if_j}$ e identifica los mapas.
\end{ejer}
\begin{ejer}[Sembrador: se cosecha al estudiar gavillas cuasicoherentes]
\label{ejer:s05-sembrador-cosecha-estudiar-gavillas-cuasicoherentes}
Si $\varphi:M\to N$ es un morfismo de $A$-módulos, verifica que la asignación $M\mapsto\widetilde M$ es exacta: una sucesión corta exacta de módulos induce una sucesión corta exacta de gavillas.
\end{ejer}


\chapter{Esquemas afines y la equivalencia contravariante}
\section{De vuelta al puente de la Sesión 5}

Arrancamos la Sesión 6. La meta es cerrar la pregunta que quedó abierta: ¿todo morfismo de espacios localmente anillados entre dos espectros proviene de un único morfismo de anillos? Dado $\varphi:A\to B$, ya conocemos $\Spec\varphi:\Spec B\to\Spec A$. En cada básico $D(f_0)\subseteq\Spec A$, la propiedad universal de la localización da un mapa compatible $A_{f_0}\to B_{\varphi(f_0)}$. Como estos mapas son compatibles con las restricciones, se extienden a un morfismo de gavillas $\mathcal O_{\Spec A}\to (\Spec\varphi)_*\mathcal O_{\Spec B}$ por el \cref{teo:s05-toda-gavilla-sobre-base-extiende}.

\begin{ejer}
\label{ejer:s06-b-f-cada-f-0}
Sea $\varphi:A\to B$ y $f=\Spec\varphi$. Para cada $f_0\in A$, verifica que $f^{-1}(D(f_0))=D(\varphi(f_0))$ y que existe un único morfismo $A_{f_0}\to B_{\varphi(f_0)}$ que extiende $\varphi$ a través de las localizaciones.
\end{ejer}

Dos preguntas organizan la sesión. Primero, ¿el par $(f,f^\sharp)$ respeta la estructura local de los tallos, de modo que es un morfismo entre espacios localmente anillados? (§2–§3.) Segundo, ¿todo morfismo de espacios localmente anillados $\Spec B\to\Spec A$ proviene, de esta manera, de un único $\varphi:A\to B$? (§4.) La respuesta afirmativa significa que ningún morfismo se pierde ni aparece al pasar entre anillos y esquemas afines; esa es la equivalencia contravariante de la sección 5.

\section{Morfismos de espacios (localmente) anillados}

\begin{defi}[Morfismo de espacios anillados]
\label{defi:s06-morfismo-espacios-anillados}
Un morfismo $(X,\mathcal O_X)\to(Y,\mathcal O_Y)$ es un par $(f,f^\sharp)$ donde $f:X\to Y$ es continuo y $f^\sharp:\mathcal O_Y\to f_*\mathcal O_X$ es un morfismo de gavillas de anillos. En cada abierto $V\subseteq Y$, la componente es $\mathcal O_Y(V)\to\mathcal O_X(f^{-1}V)$.
\end{defi}

\begin{defi}[Morfismo de espacios localmente anillados]
\label{defi:s06-morfismo-espacios-localmente-anillados}
Si los espacios son localmente anillados, $(f,f^\sharp)$ es un morfismo localmente anillado si para todo $x\in X$ el morfismo de tallos $f_x^\sharp:\mathcal O_{Y,f(x)}\to\mathcal O_{X,x}$ es local: la preimagen del ideal maximal de $\mathcal O_{X,x}$ es el ideal maximal de $\mathcal O_{Y,f(x)}$.
\end{defi}

\begin{obs}
\label{obs:s06-contencion-f-x-1-m}
La contención $(f_x^\sharp)^{-1}(\mathfrak m_x)\supseteq\mathfrak m_{f(x)}$ debe verificarse. La contención opuesta es automática para un homomorfismo de anillos locales que envía $1$ a $1$: un elemento fuera del maximal es unidad, y su imagen también. La localidad exige además que ningún elemento maximal del anillo de origen salga del maximal del anillo de llegada.
\end{obs}

\section{\texorpdfstring{$\Spec$}{Spec} es funtor contravariante a \texorpdfstring{$\LocAn$}{espacios localmente anillados}}

\begin{teo}
\label{teo:s06-todo-morfismo-anillos-b-induce}
Todo morfismo de anillos $\varphi:A\to B$ induce un morfismo de espacios localmente anillados $(\Spec B,\mathcal O_{\Spec B})\to(\Spec A,\mathcal O_{\Spec A})$.
\end{teo}
\begin{proof}
Sea $\mathfrak q\in\Spec B$ y $\mathfrak p=\varphi^{-1}(\mathfrak q)$. En los tallos, el morfismo es $\varphi_{\mathfrak q}:A_{\mathfrak p}\to B_{\mathfrak q}$, $a/s\mapsto\varphi(a)/\varphi(s)$. Basta verificar que $\varphi_{\mathfrak q}^{-1}(\mathfrak qB_{\mathfrak q})=\mathfrak pA_{\mathfrak p}$.

Para todo primo $\mathfrak q\subseteq B$, se tiene $b/1\in\mathfrak qB_{\mathfrak q}$ si y solo si $b\in\mathfrak q$. Una dirección es inmediata. Para la otra, si $b/1=c/t$ con $c\in\mathfrak q$ y $t\notin\mathfrak q$, existe $u\notin\mathfrak q$ tal que $u(tb-c)=0$. Entonces $utb\in\mathfrak q$, y como $ut\notin\mathfrak q$, la primalidad da $b\in\mathfrak q$. Aplicando esto en $A$ y en $B$ se obtiene, para $a\in A$,
\[
a\in\mathfrak p\iff\varphi(a)\in\mathfrak q
\iff\varphi(a)/1\in\mathfrak qB_{\mathfrak q}
\iff a/1\in\mathfrak pA_{\mathfrak p}.
\]
Como $s\notin\mathfrak p$ se envía a una unidad, el mismo criterio vale para toda fracción $a/s$. Por tanto el morfismo de tallos es local.
\end{proof}

La composición de morfismos corresponde, por la contravariancia, a la composición inversa de morfismos de anillos: si $A\xrightarrow{\varphi}B\xrightarrow{\psi}C$, entonces $\Spec(\psi\circ\varphi)=\Spec\varphi\circ\Spec\psi$. En los mapas de gavillas, la igualdad se verifica en cada básico por la unicidad de la propiedad universal de localización.

\begin{ejer}[Funtorialidad]
\label{ejer:s06-funtorialidad}
Verifica que $\Spec(\id_A)=\id_{\Spec A}$ y que para $A\xrightarrow{\varphi}B\xrightarrow{\psi}C$ se tiene $\Spec(\psi\circ\varphi)=\Spec\varphi\circ\Spec\psi$, incluyendo la igualdad correspondiente de morfismos de gavillas.
\end{ejer}
\begin{proof}[Desarrollo en clase]
En puntos, ambos mapas coinciden porque
\[
\Spec(\psi\circ\varphi)(\mathfrak r)=(\psi\circ\varphi)^{-1}(\mathfrak r)=\varphi^{-1}(\psi^{-1}(\mathfrak r))=(\Spec\varphi\circ\Spec\psi)(\mathfrak r).
\]
En las gavillas, sobre cada abierto básico la unicidad de la propiedad universal de la localización fuerza la igualdad de los morfismos inducidos. Para la identidad, ambos mapas de anillos y los mapas de gavillas son identidades.
\end{proof}

\section{El teorema central: plena fidelidad de \texorpdfstring{$\Spec$}{Spec}}

\begin{teo}[Plena fidelidad de $\Spec$]
\label{teo:s06-plena-fidelidad}
Para cualesquiera anillos $A,B$, la asignación
\[
\Hom_{\mathbf{Ring}}(A,B)\xrightarrow{\sim}
\Hom_{\LocAn}(\Spec B,\Spec A),\qquad
\varphi\longmapsto(\Spec\varphi,\varphi^\sharp)
\]
\end{teo}
\begin{proof}
La construcción anterior da la aplicación. Para probar que es biyectiva, sea $(f,f^\sharp)$ un morfismo localmente anillado entre los dos espectros. En secciones globales define necesariamente
\[
\varphi=f^\sharp_{\Spec A}:A=\mathcal O_{\Spec A}(\Spec A)\longrightarrow
B=\mathcal O_{\Spec B}(\Spec B).
\]
Esto prueba la unicidad del candidato. Veamos que $f=\Spec\varphi$. Para $\mathfrak q\in\Spec B$, sea $\mathfrak p=f(\mathfrak q)$. El cuadrado de secciones globales y gérmenes conmuta. Para $a\in A$, la localidad de $f^\sharp_{\mathfrak q}$ da
\[
a\in\varphi^{-1}(\mathfrak q)
\iff \varphi(a)/1\in\mathfrak qB_{\mathfrak q}
\iff f^\sharp_{\mathfrak q}(a/1)\in\mathfrak qB_{\mathfrak q}
\iff a/1\in\mathfrak pA_{\mathfrak p}
\iff a\in\mathfrak p.
\]
Así $f(\mathfrak q)=\varphi^{-1}(\mathfrak q)$ para todo $\mathfrak q$.

Finalmente, sobre cada básico $D(f_0)$ el mapa $f^\sharp$ es una extensión $A_{f_0}\to B_{\varphi(f_0)}$ del mapa $\varphi$ y envía $f_0$ a una unidad. Por la propiedad universal de $A_{f_0}$, coincide con el mapa construido a partir de $\varphi$. Los básicos cubren $\Spec A$, así que los morfismos de gavillas coinciden globalmente. Esto prueba existencia y unicidad.
\end{proof}

\section{La equivalencia contravariante}

\begin{coro}[La equivalencia]
\label{coro:s06-equivalencia}
El funtor $\Spec:\mathbf{Ring}^{\mathrm{op}}\to\LocAn$ es plenamente fiel. Restringido a su imagen esencial, establece una anti-equivalencia entre anillos y esquemas afines.
\end{coro}

\begin{defi}[Esquema afín]
\label{defi:s06-esquema-afin}
Un esquema afín es un espacio localmente anillado isomorfo, como espacio localmente anillado, a $(\Spec A,\mathcal O_{\Spec A})$ para algún anillo conmutativo $A$.
\end{defi}

Por la plena fidelidad, dar un morfismo de esquemas afines equivale a dar un morfismo de anillos en dirección opuesta, y este último es único.

\begin{obs}
\label{obs:s06-imagen-esencial-corolario-5-1}
La imagen esencial del \cref{coro:s06-equivalencia} es, por definición, la categoría de esquemas afines. Dar un morfismo entre esquemas afines equivale a dar un único morfismo de anillos en la dirección opuesta; el \cref{teo:s06-plena-fidelidad} garantiza su existencia y unicidad.
\end{obs}

\section{Ejercicios}

\begin{ejer}[Automorfismos]
\label{ejer:s06-automorfismos}
Deduce que $\operatorname{Aut}(\Spec A)\cong\operatorname{Aut}(A)^{\mathrm{op}}$ como grupos. Explica por qué aparece el grupo opuesto.
\end{ejer}
\begin{ejer}[Un caso pequeño]
\label{ejer:s06-caso-pequeno}
Describe $\Spec k$ como espacio localmente anillado cuando $k$ es cuerpo. Usa la plena fidelidad para describir el conjunto de morfismos de espacios localmente anillados $\Spec k\to\Spec A$.
\end{ejer}


\chapter{Definición de esquema y pegado}
\section{De vuelta al cierre de la Sesión 6}

La Sesión 6 cerró la anti-equivalencia entre anillos y esquemas afines. Quedó la frase por formalizar: un esquema es un espacio localmente anillado que admite una cubierta por abiertos afines. Dos preguntas organizan esta sesión: si cada abierto básico $D(f)\subseteq\Spec A$ vuelve a ser afín y cómo construir un esquema a partir de piezas con reglas de traslape. El \cref{teo:s07-lema-pegado} responde la segunda pregunta: permite construir esquemas especificando sus piezas y las reglas de traslape.

\section{Restricción de un espacio localmente anillado a un abierto}

\begin{lema}[Restricción a un abierto]
\label{lema:s07-restriccion-abierto}
Sea $(X,\mathcal O_X)$ un espacio anillado y $U\subseteq X$ abierto. Para abiertos $V\subseteq U$, definamos $\mathcal O_X|_U(V)=\mathcal O_X(V)$, con las mismas restricciones. Entonces $(U,\mathcal O_X|_U)$ es un espacio anillado y, para cada $x\in U$, $ (\mathcal O_X|_U)_x\cong\mathcal O_{X,x}$. En particular, la restricción de un espacio localmente anillado sigue siendo localmente anillada.
\end{lema}
\begin{proof}
Los axiomas de gavilla se heredan de $\mathcal O_X$. Los abiertos de $U$ que contienen a $x$ son cofinales en los abiertos de $X$ que contienen a $x$: a una vecindad $W$ de $x$ en $X$ se le asocia $W\cap U$. Por ello los colímites que definen ambos tallos son canónicamente isomorfos.
\end{proof}

\begin{defi}[Subespacio localmente anillado abierto]
\label{defi:s07-subespacio-localmente-anillado-abierto}
En la situación descrita en el \cref{lema:s07-restriccion-abierto}, $(U,\mathcal O_X|_U)$ es el subespacio localmente anillado abierto determinado por $U$. La inclusión $j:U\hookrightarrow X$ induce $j^\#:\mathcal O_X\to j_*(\mathcal O_X|_U)$; sobre un abierto $V\subseteq X$, el mapa es la restricción $\mathcal O_X(V)\to\mathcal O_X(V\cap U)$. Es un morfismo de espacios localmente anillados cuando $X$ lo es.
\end{defi}

\section{Esquemas}

\begin{defi}[Esquema]
\label{defi:s07-esquema}
Un esquema es un espacio localmente anillado $(X,\mathcal O_X)$ en el que cada $x$ tiene un abierto $U\ni x$ isomorfo a $(\Spec A,\mathcal O_{\Spec A})$ como espacio localmente anillado, para algún anillo conmutativo $A$. Sus morfismos son los de espacios localmente anillados.
\end{defi}

\begin{obs}
\label{obs:s07-no-hace-falta-definicion-adicional}
No hace falta una definición adicional de morfismo de esquemas: es el morfismo de espacios localmente anillados entre los esquemas. Todo esquema afín es un esquema, tomando el abierto afín igual al espacio entero.
\end{obs}

\begin{ejem}[El espacio afín $\mathbb A_A^n$]
\label{ejem:s07-espacio-afin-n}
Para $A$ un anillo conmutativo y $n\geq0$, $\mathbb A_A^n=\Spec A[x_1,\ldots,x_n]$ ya es un esquema afín. Si $A=k$ es algebraicamente cerrado, sus puntos cerrados son los ideales $(x_1-a_1,\ldots,x_n-a_n)$ y se identifican con $k^n$. Para $A=\mathbb Z$, el espacio registra puntos sobre cada primo de $\Spec\mathbb Z$.
\end{ejem}

\section{Abiertos afines}

\begin{prop}[$D(f)$ es un abierto afín]
\label{prop:s07-d-f-abierto-afin}
Para $f\in A$, el abierto $D(f)\subseteq\Spec A$, con la gavilla estructural restringida, es isomorfo como espacio localmente anillado a $\Spec A_f$.
\end{prop}
\begin{proof}
El homeomorfismo se obtiene en el \cref{teo:s03-correspondencia-localizaciones} $\Spec A_f\cong D(f)$. En la base común formada por los abiertos $D(fg)$, la gavilla restringida de $\Spec A$ tiene secciones $A_{fg}$, mientras que la gavilla de $\Spec A_f$ tiene secciones $(A_f)_{\bar g}$. Ambos anillos son la $A$-álgebra universal en la que $f$ y $g$ son unidades, de modo que son canónicamente isomorfos. La unicidad del \cref{teo:s05-toda-gavilla-sobre-base-extiende} da el isomorfismo de gavillas y, por consiguiente, de espacios localmente anillados.
\end{proof}

\begin{prop}[Todo abierto de un esquema es un esquema]
\label{prop:s07-todo-abierto-esquema-esquema}
Todo abierto $V$ de un esquema $X$ es un esquema con estructura $\mathcal O_X|_V$.
\end{prop}
\begin{proof}
Para $x\in V$, elige un abierto afín $U\ni x$ en $X$. En $U\cap V$, que es abierto en $U\cong\Spec A$, existe un básico $D(g)$ que contiene a $x$. Por la \cref{prop:s07-d-f-abierto-afin}, $D(g)$ es afín. Estos abiertos dan una cubierta afín de $V$.
\end{proof}

\begin{defi}[Subesquema abierto, abierto afín, inmersión abierta]
\label{defi:s07-subesquema-abierto-abierto-afin-inmersion}
Si $V\subseteq X$ es abierto, $(V,\mathcal O_X|_V)$ es el subesquema abierto de $X$ determinado por $V$, y es un abierto afín si además es afín. Un morfismo de esquemas $Y\to X$ es una inmersión abierta si factoriza como un isomorfismo $Y\xrightarrow{\sim}V$ seguido de la inclusión de un subesquema abierto $V\subseteq X$.
\end{defi}

\begin{obs}
\label{obs:s07-abiertos-d-f-forman-base}
El homomorfismo $A\to A_f$ induce, por la \cref{prop:s07-d-f-abierto-afin}, el morfismo
\[
\Spec A_f\longrightarrow\Spec A
\]
es una inmersión abierta con imagen $D(f)$. Los abiertos $D(f)$ forman una base de abiertos afines en cada esquema afín. En cualquier esquema, cada punto tiene abiertos afines arbitrariamente pequeños: dentro de una carta afín se elige un básico contenido en el abierto dado. Por tanto los abiertos afines forman una base de la topología. El Lema de Comunicación Afín (Vakil, \S5.3), que permite verificar una propiedad en una sola cubierta afín, se enunciará más adelante.
\end{obs}

\begin{ejem}[La hipérbola]
\label{ejem:s07-hiperbola}
Sea $k$ un campo y $B=k[x,y]/(xy-1)$. El mapa $k[x]\to B$ induce una inmersión abierta $\Spec B\to\mathbb A_k^1$ con imagen $D(x)$, pues $B\cong k[x]_x$ y $y$ corresponde a $x^{-1}$.
\end{ejem}
\begin{obs}
\label{obs:s07-mismo-esquema-cerrado-k-2}
El mismo esquema es cerrado en $\mathbb A_k^2$, como conjunto de ceros de $xy-1$, y abierto en $\mathbb A_k^1$, como complemento del origen. Ser abierto o cerrado depende del esquema ambiente.
\end{obs}

\section{Datos de pegado}

Pegar objetos es construir un colímite. Para espacios topológicos $A,X,Y$ con mapas $A\to X$ y $A\to Y$, el pushout es el cociente $(X\amalg Y)/\sim$ con topología cociente, que identifica las dos imágenes de cada punto de $A$. Mapas continuos desde ese cociente equivalen a mapas desde $X$ y $Y$ compatibles sobre $A$.

\begin{ejem}[El círculo]
\label{ejem:s07-circulo}
Sean $I_1,I_2$ dos copias de $[0,1]$ y $A=\set{0,1}$, incluidas como extremos. El pushout $I_1\amalg_A I_2$ es homeomorfo a $S^1$: identifica los dos ceros y los dos unos, y los intervalos son los dos arcos.
\end{ejem}

En esquemas se pegan piezas por subesquemas abiertos y los cambios de carta son isomorfismos. La receta debe especificar qué partes se identifican y cómo. Por ejemplo, dos copias de $\mathbb A_k^1$ pueden pegarse a lo largo de $D(x)$ mediante la identidad o mediante $x\mapsto1/x$; son datos distintos y producen esquemas distintos.

\begin{defi}[Datos de pegado]
\label{defi:s07-datos-pegado}
Los datos de pegado para esquemas $\set{X_i}_{i\in I}$ consisten en subesquemas abiertos $X_{ij}\subseteq X_i$, con $X_{ii}=X_i$, e isomorfismos $\varphi_{ij}:X_{ij}\xrightarrow{\sim}X_{ji}$ que satisfacen:
\begin{enumerate}[label=(P\arabic*), leftmargin=2.1em]
\item $\varphi_{ii}=\id_{X_i}$ para todo $i$;
\item $\varphi_{ji}=\varphi_{ij}^{-1}$ para todo $i,j$;
\item para toda terna $i,j,k$, $\varphi_{ij}(X_{ij}\cap X_{ik})=X_{ji}\cap X_{jk}$ y, en $X_{ij}\cap X_{ik}$, $\varphi_{ik}=\varphi_{jk}\circ\varphi_{ij}$.
\end{enumerate}
\end{defi}
\begin{obs}
\label{obs:s07-primera-parte-p3-garantiza-composicion}
La primera parte de (P3) garantiza que la composición $\varphi_{jk}\circ\varphi_{ij}$ está definida en el traslape; la segunda dice que los dos cambios de carta coinciden.
\end{obs}
\begin{obs}[Con dos piezas no hay cociclo que verificar]
\label{obs:s07-dos-piezas-no-cociclo-verificar}
Si $I=\set{1,2}$, basta elegir $X_{12}\subseteq X_1$, $X_{21}\subseteq X_2$ y un isomorfismo $\varphi_{12}:X_{12}\to X_{21}$. Las condiciones diagonales e inversas determinan $\varphi_{11}$, $\varphi_{22}$ y $\varphi_{21}$; la condición (P3) se sigue de ellas.
\end{obs}

\section{El teorema central: el Lema de Pegado}

\begin{teo}[Lema de Pegado]
\label{teo:s07-lema-pegado}
Dados datos de pegado como en la \cref{defi:s07-datos-pegado}, existe un esquema $X$ y unas inmersiones abiertas $\psi_i:X_i\to X$ tales que
\begin{enumerate}[label=(\alph*)]
\item $X=\bigcup_i\psi_i(X_i)$;
\item $\psi_i(X_i)\cap\psi_j(X_j)=\psi_i(X_{ij})$ para todo $i,j$;
\item $\psi_j\circ\varphi_{ij}=\psi_i$ en $X_{ij}$.
\end{enumerate}
Además, si $(X',\psi_i')$ es otro esquema con inmersiones abiertas que satisfacen las mismas propiedades, existe un único isomorfismo $h:X\to X'$ con $h\circ\psi_i=\psi_i'$ para todo $i$.
\end{teo}
La demostración del \cref{teo:s07-lema-pegado} se desarrollará a continuación, en los cinco pasos de los \cref{sec:s07-paso-1,sec:s07-paso-2,sec:s07-paso-3,sec:s07-paso-4,sec:s07-paso-5}:
\begin{enumerate}
\item Construir el conjunto $X$ y las aplicaciones $\psi_i:X_i\to X$.
\item Dar a $X$ la topología cociente y probar que cada $\psi_i$ es un homeomorfismo sobre un abierto.
\item Construir $\mathcal O_X$ con familias compatibles de secciones y verificar los axiomas de gavilla.
\item Probar que $(X,\mathcal O_X)$ es un esquema y que las $\psi_i$ son inmersiones abiertas.
\item Probar la unicidad del pegado y del isomorfismo compatible con las $\psi_i$.
\end{enumerate}
\emph{Referencias clásicas:} Hartshorne, II.2, ejercicio 2.12; EGA I, \S2.3.

Como ejemplo guía, podemos pensar en dos copias de $\mathbb A_k^1$ pegadas por la identidad en $D(x)$. El conjunto resultante tendrá dos orígenes y una sola copia de cada punto de $D(x)$. El \cref{ejem:s07-recta-origen-doble} retomará esta construcción y probará que no es afín.

\subsection{Paso 1: el conjunto subyacente}
\label{sec:s07-paso-1}

Tomemos $X_0=\coprod_iX_i$ y definamos, para $x\in X_i$, $y\in X_j$,
\[
x\sim y\qquad\Longleftrightarrow\qquad x\in X_{ij}\quad\text{y}\quad y=\varphi_{ij}(x).
\]
La condición diagonal da reflexividad y la inversa $\varphi_{ji}=\varphi_{ij}^{-1}$ da simetría. Para transitividad, supongamos $x\sim y$ y $y\sim z$, con $x\in X_i$, $y\in X_j$, $z\in X_k$. Entonces $y=\varphi_{ij}(x)$ pertenece a $X_{ji}\cap X_{jk}$. Por la condición de traslape, existe $x'\in X_{ij}\cap X_{ik}$ con $y=\varphi_{ij}(x')$. La inyectividad de $\varphi_{ij}$ implica $x=x'$, así que $x\in X_{ik}$. La condición de cociclo da $\varphi_{ik}(x)=\varphi_{jk}(\varphi_{ij}(x))=z$, y por tanto $x\sim z$.

Definimos $X=X_0/\sim$ y $\psi_i$ como la composición de la inclusión de $X_i$ con la proyección al cociente. Cada $\psi_i$ es inyectiva: si $x,x'\in X_i$ son equivalentes, entonces $x'=\varphi_{ii}(x)=x$. Se tiene $\psi_j\circ\varphi_{ij}=\psi_i$. Toda clase tiene representante en alguna pieza, de modo que se cumple (a). Para (b), la contención $\psi_i(X_{ij})\subseteq\psi_i(X_i)\cap\psi_j(X_j)$ se sigue de (c). Recíprocamente, si $\psi_i(x)=\psi_j(y)$, entonces $x\sim y$, así que $x\in X_{ij}$. Esto prueba (a)--(c) para los conjuntos subyacentes.

\subsection{Paso 2: la topología}
\label{sec:s07-paso-2}

La topología de $X$ es la cociente de $q:X_0\to X$, la más fina que hace continua a $q$. Así, $W\subseteq X$ es abierto si y solo si $\psi_i^{-1}(W)$ es abierto en cada $X_i$. Para $V\subseteq X_i$ abierto,
\[
\psi_j^{-1}(\psi_i(V))=\varphi_{ij}(V\cap X_{ij}).
\]
Un punto $y\in X_j$ está en la preimagen de la izquierda si y solo si $y\sim x$ para algún $x\in V$. Esto equivale a que $y\in X_{ji}$ y $\varphi_{ji}(y)\in V$, o bien $y\in\varphi_{ij}(V\cap X_{ij})$.

Como $\varphi_{ij}$ es un isomorfismo entre abiertos, el lado derecho es abierto en $X_j$. La topología cociente implica que $\psi_i(V)$ es abierto en $X$. Por tanto $\psi_i$ es inyectiva, continua y abierta: un homeomorfismo sobre $\psi_i(X_i)$.

\subsection{Paso 3: la gavilla estructural}
\label{sec:s07-paso-3}

Para un isomorfismo $\varphi:Y\to Z$ y un abierto $U'\subseteq Z$, el transporte de secciones es
\[
\varphi^\sharp:\mathcal O_Z(U')\xrightarrow{\sim}\mathcal O_Y(\varphi^{-1}(U'))
\]
compatible con restricciones. El transporte lleva secciones de $Z$ a secciones de $Y$; esta dirección determina la condición de compatibilidad en los traslapes.

\begin{defi}[Imagen directa]
\label{defi:s07-imagen-directa}
Si $\psi:Y\to X$ es continua y $\mathcal F$ es una gavilla en $Y$, su imagen directa es la gavilla $\psi_*\mathcal F$ en $X$ dada por $(\psi_*\mathcal F)(W)=\mathcal F(\psi^{-1}(W))$, con las restricciones inducidas por las de $\mathcal F$.
\end{defi}

\begin{obs}
\label{obs:s07-imagen-directa-pushout}
La gavilla $(\psi_i)_*\mathcal O_{X_i}$ se llama imagen directa o \emph{pushforward}. El pushout topológico de la construcción anterior es un colímite que identifica puntos. Las secciones de la \cref{defi:s07-gavilla-estructural-pegado} se describen como familias compatibles en el producto
\[
\prod_i\bigl((\psi_i)_*\mathcal O_{X_i}\bigr)(W)
\]
\end{obs}

\begin{defi}[La gavilla estructural del pegado]
\label{defi:s07-gavilla-estructural-pegado}
Para un abierto $W\subseteq X$, definimos
\[
\mathcal O_X(W):=\set*{(s_i)_{i\in I}\in\prod_{i\in I}\mathcal O_{X_i}(\psi_i^{-1}(W))\given
s_i|_{\psi_i^{-1}(W)\cap X_{ij}}=\varphi_{ij}^{\sharp}\bigl(s_j|_{\psi_j^{-1}(W)\cap X_{ji}}\bigr)\text{ para todo }i,j}.
\]
\end{defi}

\begin{obs}
\label{obs:s07-condicion-gavilla-pegado}
Por la propiedad (c) del \cref{teo:s07-lema-pegado}, $\varphi_{ij}$ lleva $\psi_i^{-1}(W)\cap X_{ij}$ sobre $\psi_j^{-1}(W)\cap X_{ji}$. Así, ambos lados de la condición de la \cref{defi:s07-gavilla-estructural-pegado} son secciones sobre el mismo abierto. Además, $\mathcal O_X(W)$ es un subanillo del producto, pues las condiciones son igualdades entre homomorfismos de anillos. Falta verificar que sea una pregavilla y satisfaga los axiomas de gavilla.
\end{obs}

\begin{prop}
\label{prop:s07-gavilla-estructural}
La asignación $W\mapsto\mathcal O_X(W)$ de la \cref{defi:s07-gavilla-estructural-pegado}, con las restricciones coordenada a coordenada, es una gavilla de anillos en $X$.
\end{prop}
\begin{proof}
\emph{Pregavilla.} Si $W_0\subseteq W$, restringir cada coordenada lleva una familia compatible sobre $W$ a otra sobre $W_0$, pues los mapas $\varphi_{ij}^{\sharp}$ conmutan con las restricciones. La identidad y la composición se verifican coordenada a coordenada, porque son las de las gavillas $\mathcal O_{X_i}$. Por tanto, $\mathcal O_X$ es una pregavilla.

\emph{Identidad.} Sean dos secciones de $\mathcal O_X(W)$ que coinciden en cada abierto $W_\alpha$ de una cubierta de $W$. Fijemos $i$. Sus coordenadas $i$ coinciden en cada $\psi_i^{-1}(W_\alpha)$, que cubren $\psi_i^{-1}(W)$. El axioma de identidad para $\mathcal O_{X_i}$ implica que coinciden en todo $\psi_i^{-1}(W)$. Esto vale para cada $i$, así que las secciones son iguales.

\emph{Pegado.} Sean familias compatibles $(s_i^\alpha)_{i\in I}\in\mathcal O_X(W_\alpha)$ sobre una cubierta $W=\bigcup_\alpha W_\alpha$. Para cada $i$, el axioma de pegado de $\mathcal O_{X_i}$ produce una única sección $s_i\in\mathcal O_{X_i}(\psi_i^{-1}(W))$ cuyas restricciones son $s_i^\alpha$. Falta verificar la condición entre $s_i$ y $s_j$. En $\psi_i^{-1}(W)\cap X_{ij}$, los dos lados coinciden al restringirse a cada $\psi_i^{-1}(W_\alpha)\cap X_{ij}$, pues cada familia $(s_i^\alpha)_i$ ya cumplía esa condición. Estos abiertos cubren el traslape; el axioma de identidad para $\mathcal O_{X_i}$ da la igualdad en todo él. Así $(s_i)_i\in\mathcal O_X(W)$ y es la única sección que restringe a las familias dadas. Por tanto, $\mathcal O_X$ es una gavilla.
\end{proof}

\subsection{Paso 4: el pegado es un esquema}
\label{sec:s07-paso-4}

Fijemos $i\in I$ y un abierto $V\subseteq X_i$, y pongamos $W=\psi_i(V)$. Por el \cref{sec:s07-paso-2},
\[
V_j:=\psi_j^{-1}(W)=\varphi_{ij}(V\cap X_{ij})
\]
es abierto en $X_j$, y $V_i=V$.

\begin{lema}
\label{lema:s07-proyeccion-coordenada-iso}
La proyección
\[
\pi_i:\mathcal O_X(\psi_i(V))\longrightarrow\mathcal O_{X_i}(V),
\qquad (s_j)_{j\in I}\longmapsto s_i,
\]
es un isomorfismo de anillos.
\end{lema}
\begin{proof}
\emph{Inyectividad.} La condición de compatibilidad para el par $(i,j)$ determina $s_j$ a partir de $s_i$. Como $\varphi_{ij}$ es un isomorfismo, dos familias con igual coordenada $i$ son iguales.

\emph{Sobreyectividad.} Sea $s\in\mathcal O_{X_i}(V)$. La única elección posible, despejando la condición para el par $(i,j)$, es
\[
s_i=s,\qquad s_j=(\varphi_{ij}^{\sharp})^{-1}\bigl(s|_{V\cap X_{ij}}\bigr).
\]
La definición de $\mathcal O_X$ exige comprobar la compatibilidad para todos los pares.

\emph{Compatibilidad del par $(j,k)$.} Sean $T=V_j\cap X_{jk}$ y $T'=\varphi_{ji}(T)$, y pongamos $\sigma=s|_{T'}$. Por (P3), $T'\subseteq V\cap X_{ij}\cap X_{ik}$ y $\varphi_{jk}(T)=\varphi_{ik}(T')$. La definición de $s_j$ y $s_k$ da
\[
s_j|_T=\varphi_{ji}^{\sharp}(\sigma),\qquad
s_k|_{\varphi_{jk}(T)}=\varphi_{ki}^{\sharp}(\sigma).
\]
Por tanto,
\[
\varphi_{jk}^{\sharp}\bigl(s_k|_{\varphi_{jk}(T)}\bigr)
=(\varphi_{ki}\circ\varphi_{jk})^{\sharp}(\sigma)
=\varphi_{ji}^{\sharp}(\sigma)=s_j|_T.
\]
La segunda igualdad usa $\varphi_{jk}=\varphi_{ik}\circ\varphi_{ji}$ y $\varphi_{ki}=\varphi_{ik}^{-1}$. Así, la familia construida pertenece a $\mathcal O_X(W)$ y $\pi_i$ es sobreyectiva.
\end{proof}

\begin{prop}
\label{prop:s07-pegado-esquema}
El espacio $(X,\mathcal O_X)$ es un esquema, y cada $\psi_i:X_i\to X$ es una inmersión abierta que cumple las propiedades (a), (b) y (c) del \cref{teo:s07-lema-pegado}.
\end{prop}
\begin{proof}
\emph{De $\pi_i$ a $\psi_i^\sharp$.} Las proyecciones a la coordenada $i$ conmutan con las restricciones y definen, para cualquier abierto $W\subseteq X$, el mapa estructural
\[
\psi_i^\sharp(W):\mathcal O_X(W)\longrightarrow
\mathcal O_{X_i}(\psi_i^{-1}(W)),\qquad (s_j)_j\longmapsto s_i.
\]
Sobre $W=\psi_i(V)$, este mapa es precisamente $\pi_i$. El \cref{lema:s07-proyeccion-coordenada-iso} muestra que ese mapa es un isomorfismo; junto con el homeomorfismo sobre un abierto del \cref{sec:s07-paso-2}, prueba que $\psi_i$ es una inmersión abierta. Las propiedades (a), (b) y (c) ya valían para los conjuntos subyacentes, y la compatibilidad de las coordenadas prueba (c) como igualdad de morfismos: en cada traslape,
\[
(\psi_j\circ\varphi_{ij})^\sharp
=\varphi_{ij}^\sharp\circ\psi_j^\sharp=\psi_i^\sharp.
\]

\emph{$X$ es un esquema.} Para $x\in X_i$, el tallo $\mathcal O_{X,\psi_i(x)}$ es isomorfo a $\mathcal O_{X_i,x}$, que es local, pues $\psi_i(X_i)$ es abierto y $\psi_i$ es un isomorfismo sobre su imagen. Así, $X$ es localmente anillado. Los abiertos $\psi_i(X_i)$ cubren $X$. Como cada $X_i$ es un esquema y $\psi_i$ es un isomorfismo sobre su imagen, cada punto de $X$ tiene una base de abiertos afines. Por la \cref{defi:s07-esquema}, queda probada la existencia.
\end{proof}

\subsection{Paso 5: unicidad}
\label{sec:s07-paso-5}

\begin{prop}
\label{prop:s07-unicidad-pegado}
Si $(X',\psi_i')$ es otro esquema con inmersiones abiertas que cumplen las propiedades (a), (b) y (c) del \cref{teo:s07-lema-pegado}, existe un único isomorfismo $h:X\to X'$ tal que $h\circ\psi_i=\psi_i'$ para todo $i$.
\end{prop}
\begin{proof}
Definimos $h(\psi_i(x))=\psi_i'(x)$. Si $\psi_i(x)=\psi_j(y)$, la relación del primer paso dice que $x\in X_{ij}$ e $y=\varphi_{ij}(x)$. Por (c) en $X'$, $\psi_i'(x)=\psi_j'(y)$; así, la definición no depende del representante.

Para construir la inversa hay que probar el recíproco. Si $\psi_i'(x)=\psi_j'(y)$, la propiedad (b) en $X'$ da un $z\in X_{ij}$ con $\psi_i'(x)=\psi_i'(z)$. La inyectividad de $\psi_i'$ implica $x=z\in X_{ij}$. Por (c),
\[
\psi_j'(\varphi_{ij}(x))=\psi_i'(x)=\psi_j'(y),
\]
y la inyectividad de $\psi_j'$ da $y=\varphi_{ij}(x)$. Por tanto $x\sim y$ y queda bien definida $\bar h(\psi_i'(x))=\psi_i(x)$. Las identidades $h\circ\psi_i=\psi_i'$ y $\bar h\circ\psi_i'=\psi_i$ prueban que $h$ y $\bar h$ son inversas.

En cada abierto $\psi_i(X_i)$, el mapa es $\psi_i'\circ\psi_i^{-1}$, un isomorfismo de esquemas. Estos mapas coinciden en los traslapes por (c), así que se pegan a un isomorfismo $h:X\to X'$. La unicidad se sigue de que los $\psi_i(X_i)$ cubren $X$.
\end{proof}

\begin{proof}[Conclusión de la demostración]
Los \cref{sec:s07-paso-1,sec:s07-paso-2} construyen el espacio topológico $X$ y las aplicaciones $\psi_i$. La \cref{prop:s07-gavilla-estructural} verifica los axiomas de gavilla para la \cref{defi:s07-gavilla-estructural-pegado}; el \cref{lema:s07-proyeccion-coordenada-iso} y la \cref{prop:s07-pegado-esquema} muestran que $X$ es un esquema y que las $\psi_i$ son las inmersiones abiertas requeridas. La \cref{prop:s07-unicidad-pegado} da la unicidad. Esto demuestra el \cref{teo:s07-lema-pegado}.
\end{proof}

\section{El pegado como colímite}

\begin{prop}[El pegado es un coigualador]
\label{prop:s07-pegado-coigualador}
Sean $E=\coprod_{i,j}X_{ij}$ y $X_\amalg=\coprod_iX_i$, con mapas $s,t:E\rightrightarrows X_\amalg$ dados respectivamente por la inclusión de $X_{ij}$ en $X_i$ y por $\varphi_{ij}$ seguida de la inclusión en $X_j$. Los $\psi_i$ definen $\psi:X_\amalg\to X$, con $\psi\circ s=\psi\circ t$. El diagrama $E\rightrightarrows X_\amalg\xrightarrow{\psi}X$ es un coigualador en la categoría de esquemas. Es decir, para todo esquema $Y$ y todo $f:X_\amalg\to Y$ con $f\circ s=f\circ t$, existe un único $\bar f:X\to Y$ con $\bar f\circ\psi=f$.
\end{prop}
\begin{proof}
Un morfismo $X_\amalg\to Y$ equivale a una familia $f_i:X_i\to Y$. La igualdad $f\circ s=f\circ t$ dice exactamente que $f_i|_{X_{ij}}=f_j\circ\varphi_{ij}$. En los abiertos $\psi_i(X_i)$, los mapas $f_i\circ\psi_i^{-1}$ coinciden en los traslapes por esta igualdad, así que se pegan a un único morfismo $X\to Y$. La unicidad se sigue de que los $\psi_i(X_i)$ cubren $X$.
\end{proof}

\begin{obs}
\label{obs:s07-dar-morfismo-x-equivale-dar}
Dar un morfismo $X\to Y$ equivale a dar una familia $f_i:X_i\to Y$ con $f_i|_{X_{ij}}=f_j\circ\varphi_{ij}$. El argumento usa solo las propiedades (a), (b), (c), por lo que vale para cualquier pegado que las satisfaga y expresa de nuevo la unicidad.
\end{obs}

La distinción con los límites se mantiene: en la categoría de esquemas, los límites finitos existen por el objeto final $\Spec\mathbb Z$ junto con los productos fibrados. El objeto inicial es el esquema vacío.

\section{Consistencia: pegado y cubierta}

\begin{coro}
\label{coro:s07-x-esquema-cubierta-abiertos-u}
Sea $(X,\mathcal O_X)$ un esquema con una cubierta por abiertos $\set{U_i}_{i\in I}$. Tomando $U_{ij}=U_i\cap U_j$ y los mapas de transición identidad, los datos satisfacen la \cref{defi:s07-datos-pegado} y el esquema producido por el \cref{teo:s07-lema-pegado} es canónicamente isomorfo a $X$.
\end{coro}
\begin{proof}
El mapa de cada pieza $U_i$ a $X$ es una inmersión abierta. Como los $U_i$ cubren $X$ y sus intersecciones son exactamente $U_{ij}$, estos mapas satisfacen las tres propiedades del \cref{teo:s07-lema-pegado}. Su unicidad da el isomorfismo canónico.
\end{proof}
\begin{obs}
\label{obs:s07-definicion-esquema-cubiertas-afines-teorema}
En particular, el corolario se aplica a una cubierta por abiertos afines. La \cref{defi:s07-esquema} y el \cref{teo:s07-lema-pegado} describen la misma clase de objetos desde dos direcciones. Todo esquema resulta de pegar su propia cubierta afín; el pegado también permite construir esquemas cuando las piezas no están previamente dentro de uno común.
\end{obs}

\begin{samepage}
\section{Un ejemplo: la recta con origen doble}

\begin{ejem}[La recta con origen doble]
\label{ejem:s07-recta-origen-doble}
Sea $k$ un campo, $X_1=X_2=\mathbb A_k^1=\Spec k[x]$, $X_{12}=X_{21}=D(x)$ y $\varphi_{12}=\id$. Por la \cref{obs:s07-dos-piezas-no-cociclo-verificar} son datos de pegado. Sea $X$ el esquema que produce el \cref{teo:s07-lema-pegado}, con $\psi_1,\psi_2$. Los puntos $O_1:=\psi_1(0)$ y $O_2:=\psi_2(0)$ son distintos porque $0\notin X_{12}$. Fuera de los orígenes hay una sola copia de cada punto. Es la recta con origen doble, un esquema que no es afín.

\begin{center}
\begin{tikzpicture}[>=to, every node/.style={font=\small}]
  \draw[->] (-3,0)--(-0.25,0);
  \draw[->] (0.25,0)--(3,0);
  \draw[->] (-3,-0.45)--(-0.25,-0.45);
  \draw[->] (0.25,-0.45)--(3,-0.45);
  \fill[rojoCIMAT] (0,0) circle (2pt);
  \fill[grisCIMAT] (0,-0.45) circle (2pt);
  \node[above] at (0,0) {$O_1$};
  \node[below] at (0,-0.45) {$O_2$};
  \node[right] at (3,0) {$X_1$};
  \node[right] at (3,-0.45) {$X_2$};
  \node at (0.9,-1.05) {Las dos copias se identifican en $D(x)$};
\end{tikzpicture}
\end{center}

Por la \cref{defi:s07-gavilla-estructural-pegado}, una sección global es un par de polinomios en $k[x]$ que coinciden en $k[x]_x$; la inclusión $k[x]\hookrightarrow k[x]_x$ es inyectiva, así que los dos polinomios son iguales. Por tanto $\Gamma(X,\mathcal O_X)=k[x]$. Si $X$ fuera afín, sería isomorfo a $\Spec\Gamma(X,\mathcal O_X)=\Spec k[x]$. El mapa canónico al espectro de secciones globales viene dado por
\[
p\longmapsto\set*{s\in\Gamma(X,\mathcal O_X)\given s_p\in\mathfrak m_p},
\]
donde $\mathfrak m_p$ es el ideal maximal de $\mathcal O_{X,p}$. Ambos orígenes van al ideal $(x)$, de modo que este mapa no es inyectivo y $X$ no es afín.
\end{ejem}
\end{samepage}

\begin{obs}[El colímite entre esquemas afines]
\label{obs:s07-colimite-entre-esquemas-afines}
La categoría de esquemas afines es equivalente a la opuesta de la categoría de anillos. El pushout de $\Spec T\to\Spec R$ y $\Spec T\to\Spec S$ calculado entre afines es $\Spec(R\times_T S)$. Para $R=S=k[x]$, $T=k[x]_x$ y las inclusiones, $k[x]\times_{k[x]_x}k[x]\cong k[x]$. Por tanto, el pushout afín del diagrama
\[
\Spec k[x]\longleftarrow\Spec k[x]_x\longrightarrow\Spec k[x]
\]
es $\Spec k[x]$; en esquemas, el pushout es la recta con origen doble. Por la \cref{prop:s07-pegado-coigualador}, para dos piezas se obtiene la propiedad universal del pushout. La recta con origen doble se aplica sobre $\mathbb A_k^1$ sin ser isomorfa a ella. La inclusión de esquemas afines en esquemas no preserva este colímite. Por eso la demostración del \cref{teo:s07-lema-pegado} construye primero el espacio y la gavilla, y reconoce después los abiertos afines.
\end{obs}

\section{Ejercicios}

\begin{ejer}[Pegado con una sola pieza]
\label{ejer:s07-pegado-sola-pieza}
Verifica que si $I=\set{1}$ (un solo esquema $X_1$, con $X_{11}=X_1$ y $\varphi_{11}=\id$), el esquema producido por el \cref{teo:s07-lema-pegado} es canónicamente isomorfo a $X_1$, con $\psi_1$ la identidad.
\end{ejer}

\begin{samepage}
\begin{ejer}[Dos copias de $\Spec A$ pegadas por la identidad]
\label{ejer:s07-dos-copias-pegadas-identidad}
Sean $X_1=X_2=\Spec A$ dos copias, $f\in A$, $X_{12}=X_{21}=D(f)$ y el isomorfismo de pegado la identidad. Describe el esquema pegado en cada caso.
\begin{enumerate}
\item $f$ es invertible en $A$.
\item $A=k[x]$ con $k$ un campo y $f=x$. Pista: los puntos de $\Spec k[x]$ que no están en $D(x)$ no se identifican entre las dos copias.
\end{enumerate}
\end{ejer}
\end{samepage}

\begin{ejer}[El plano con recta doble]
\label{ejer:s07-plano-recta-doble}
Sea $k$ un campo, $X_1=X_2=\mathbb A_k^2=\Spec k[x,y]$, $X_{12}=X_{21}=D(x)$ y $\varphi_{12}=\id_{D(x)}$. Sea $X$ el esquema pegado.
\begin{enumerate}
\item Calcula $\Gamma(X,\mathcal O_X)$.
\item Muestra que $X$ no es afín, adaptando el argumento del \cref{ejem:s07-recta-origen-doble}. Pista: considera el ideal primo $(x)\subseteq k[x,y]$ visto en cada una de las dos copias.
\item Explica por qué esto no contradice el \cref{coro:s07-x-esquema-cubierta-abiertos-u}.
\end{enumerate}
\end{ejer}

\begin{ejer}[La imagen directa de una gavilla]
\label{ejer:s07-imagen-directa-gavilla}
Sean $\psi:Y\to X$ continua y $\mathcal F$ una gavilla en $Y$. Verifica que $\psi_*\mathcal F$ de la \cref{defi:s07-imagen-directa} es una gavilla y que la \cref{defi:s07-gavilla-estructural-pegado} da el subanillo de $\prod_i((\psi_i)_*\mathcal O_{X_i})(W)$ formado por las familias compatibles.
\end{ejer}

\begin{ejer}[Restricción de un morfismo a un abierto]
\label{ejer:s07-restriccion-morfismo-abierto}
Sea $f:X\to Y$ un morfismo de esquemas y $V\subseteq Y$ abierto. Verifica que $f^{-1}(V)$ es abierto y que la restricción induce un morfismo de esquemas $f|_{f^{-1}(V)}:f^{-1}(V)\to V$.
Usa la \cref{prop:s07-todo-abierto-esquema-esquema} para que $f^{-1}(V)$ y $V$ sean esquemas, y el \cref{lema:s07-restriccion-abierto} para sus estructuras locales. \emph{[Cf. Vakil, ejercicio de restricción de morfismos.]}
\end{ejer}


\chapter{Ejemplos globales}
\section{Lo que traemos de la Sesión 7}

\emph{Advertencia. Dos correcciones de la sesión pasada.} El morfismo canónico entre el límite de un diagrama y el producto de sus objetos va del límite al producto: el límite es un cono sobre el diagrama sin flechas y se factoriza por el producto. La gavilla $\mathcal O_X$ del pegado se obtiene, sobre cada abierto, como un límite inverso o proyectivo de anillos. Un límite directo es un colímite sobre un conjunto dirigido; el pegado no es de ese tipo.

Para los tallos, en cambio, $\mathcal O_{X,x}=\varinjlim_{U\ni x}\mathcal O_X(U)$ es un límite directo o colímite de las restricciones a vecindades cada vez menores.

La sesión anterior probó el \myref{teo:s07-lema-pegado}. Hoy lo usamos con dos piezas, caso en el que el cociclo se cumple automáticamente.

\begin{defi}[Datos de pegado]
\label{defi:s08-datos-pegado}
Unos datos de pegado para una familia de esquemas $\set{X_i}_{i\in I}$ consisten, para cada par $(i,j)\in I\times I$, en un subesquema abierto $X_{ij}\subseteq X_i$, con $X_{ii}=X_i$, y un isomorfismo de esquemas $\varphi_{ij}:X_{ij}\to X_{ji}$ que cumplen:
\begin{enumerate}[label=(P\arabic*), leftmargin=2.1em]
\item $\varphi_{ii}=\id_{X_i}$ para todo $i$;
\item $\varphi_{ji}=\varphi_{ij}^{-1}$ para todo $i,j$;
\item para toda terna $i,j,k$, $\varphi_{ij}(X_{ij}\cap X_{ik})=X_{ji}\cap X_{jk}$ y $\varphi_{ik}=\varphi_{jk}\circ\varphi_{ij}$ en $X_{ij}\cap X_{ik}$.
\end{enumerate}
\end{defi}

Con dos piezas, $I=\set{1,2}$, basta dar abiertos $X_{12}\subseteq X_1$, $X_{21}\subseteq X_2$ y un isomorfismo $\varphi_{12}:X_{12}\to X_{21}$. Se pone $\varphi_{11}=\varphi_{22}=\id$ y $\varphi_{21}=\varphi_{12}^{-1}$: (P1) y (P2) valen por construcción y (P3) se sigue de ellas. \emph{[Hartshorne, II.2, ejercicio 2.12; \cref{obs:s07-dos-piezas-no-cociclo-verificar}.]}

\begin{teo}[Lema de Pegado]
\label{teo:s08-lema-pegado}
Dados datos de pegado para $\set{X_i}_{i\in I}$, existe un esquema $X$ con inmersiones abiertas $\psi_i:X_i\to X$ tales que:
\begin{enumerate}[label=(\alph*)]
\item $X=\bigcup_{i\in I}\psi_i(X_i)$;
\item $\psi_i(X_i)\cap\psi_j(X_j)=\psi_i(X_{ij})$ para todo $i,j$;
\item $\psi_j\circ\varphi_{ij}=\psi_i$ en $X_{ij}$ para todo $i,j$.
\end{enumerate}
Si $(X',\psi_i')$ es otro esquema con inmersiones abiertas que cumplen estas propiedades, existe un único isomorfismo $h:X\to X'$ con $h\circ\psi_i=\psi_i'$ para todo $i$.
\end{teo}
La demostración está en el \cref{teo:s07-lema-pegado}. \emph{[Hartshorne, II.2, ejercicio 2.12; EGA I, \S2.3.]}

\begin{obs}[Morfismos desde un pegado]
\label{obs:s08-morfismos-desde-pegado}
Dar un morfismo $f:X\to Y$ equivale a dar morfismos $f_i:X_i\to Y$ tales que $f_i|_{X_{ij}}=f_j\circ\varphi_{ij}$ para todo $i,j$, por la \cref{prop:s07-pegado-coigualador}. Hoy usaremos esta propiedad con $Y=\Spec A$.
\end{obs}

% Presentación S08, diapositiva 6: coigualador y morfismos desde el pegado.
\begin{center}
\begin{tikzcd}[column sep=huge]
\displaystyle\coprod_{i,j}X_{ij}
  \arrow[r,shift left=.65ex,"\text{inclusión}"]
  \arrow[r,shift right=.65ex,swap,"\varphi_{ij}\text{ e inclusión}"]
& \displaystyle\coprod_iX_i \arrow[r,"\psi"] & X \arrow[r,"f"] & Y
\end{tikzcd}
\end{center}

\section{El plano afín menos el origen}

\begin{ejem}[El plano afín menos el origen]
\label{ejem:s08-plano-menos-origen}
Sea $k$ un campo, $A=k[x,y]$ y $\mathbb A_k^2=\Spec A$. El punto cerrado $(x,y)$ es el único primo que contiene a $(x,y)$. Su complemento
\[
U=\mathbb A_k^2\setminus\set{(x,y)}=D(x)\cup D(y)
\]
es abierto, pues $V(x)\cap V(y)=V(x,y)=\set{(x,y)}$. Con su estructura de subesquema abierto, $U$ es un esquema. Vamos a calcular $\Gamma(U,\mathcal O_U)$. \emph{[Vakil, \S4.4.1 a \S4.4.3.]}
\end{ejem}

% Presentación S08, diapositiva 8: los cuatro dibujos de abiertos del plano.
\begin{center}
\begin{tikzpicture}[every node/.style={font=\small}]
\foreach \a/\col in {0/blue,3.6/green,7.2/green,10.8/orange}{
  \fill[\col!12] (\a,0) rectangle ++(2,2);
  \draw[grisCIMAT] (\a,0) rectangle ++(2,2);
}
\draw[rojoCIMAT,dashed,thick] (1,0)--(1,2);
\draw[rojoCIMAT,dashed,thick] (3.6,1)--(5.6,1);
\draw[rojoCIMAT,dashed,thick] (8.2,0)--(8.2,2);
\draw[rojoCIMAT,dashed,thick] (7.2,1)--(9.2,1);
\draw[grisCIMAT] (11.8,0)--(11.8,2) (10.8,1)--(12.8,1);
\fill[white] (11.8,1) circle (3pt);
\draw[rojoCIMAT,thick] (11.8,1) circle (3pt);
\node at (1,-.4) {$D(x)$};
\node at (1,-.8) {$x\ne0$};
\node at (4.6,-.4) {$D(y)$};
\node at (4.6,-.8) {$y\ne0$};
\node at (8.2,-.4) {$D(x)\cap D(y)=D(xy)$};
\node at (8.2,-.8) {$xy\ne0$};
\node at (11.8,-.4) {$U=D(x)\cup D(y)$};
\node at (11.8,-.8) {todo salvo el origen};
\end{tikzpicture}
\end{center}

\begin{prop}[El plano menos el origen]
\label{prop:s08-secciones-plano-menos-origen}
La restricción
\[
\Gamma(\mathbb A_k^2,\mathcal O)\longrightarrow\Gamma(U,\mathcal O_U)
\]
es un isomorfismo. En particular, $\Gamma(U,\mathcal O_U)=k[x,y]$.
\end{prop}
\begin{proof}
\emph{Las secciones.} La cubierta $U=D(x)\cup D(y)$ tiene traslape $D(xy)$. Por los axiomas de identidad y pegado, una sección global es un par $(f,g)\in A_x\times A_y$ cuyas imágenes coinciden en $A_{xy}$. Como $A$ es un dominio, las localizaciones son inyectivas y podemos considerar
\[
A\subseteq A_x,A_y\subseteq A_{xy},
\qquad \Gamma(U,\mathcal O_U)=A_x\cap A_y.
\]
Si $h=p/x^a=q/y^b$ pertenece a la intersección, con $p,q\in A$ y $a,b\geq0$, entonces $py^b=qx^a$. El anillo $A$ es un dominio de factorización única, y $x,y$ son primos distintos. De $x^a\mid py^b$ y $x\nmid y$ se obtiene $x^a\mid p$, luego $h\in A$. La inclusión opuesta es inmediata: la restricción es un isomorfismo.

\emph{Advertencia.} La inclusión $A\hookrightarrow A_x$ usa que $A$ es un dominio; para un anillo cualquiera la localización puede perder elementos.
\end{proof}

\begin{obs}[Lectura geométrica]
\label{obs:s08-hartogs}
Los polos de $p/x^a$ están en la recta $x=0$, y los de $q/y^b$ en la recta $y=0$. Si representan la misma función, sus polos tendrían que estar en el origen. En dimensión dos, los polos de una función forman una curva y un punto no lo es. Toda función regular en $\mathbb A_k^2$ menos el origen se extiende al plano: es el análogo algebraico del lema de Hartogs. En $\mathbb A_k^1\setminus\set{0}$ sí hay funciones nuevas, como $1/x$. Para $n\geq2$ se usa la cubierta $D(x_1),\ldots,D(x_n)$ de $\mathbb A_k^n$ menos el origen; véase el \cref{ejer:s08-afin-menos-origen}.
\end{obs}

\emph{Comparación topológica.} El plano real $\mathbb R^2$ es contráctil, mientras que $\mathbb R^2\setminus\set{0}$ tiene el tipo de homotopía de $S^1$ y grupo fundamental $\mathbb Z$. En el caso complejo, $\mathbb C^2\setminus\set{0}$ tiene el tipo de homotopía de $S^3$, que es simplemente conexa; lo distingue de $\mathbb C^2$ su tercer grupo de homología $H_3=\mathbb Z$ (o su $\pi_3$).

\clearpage
\subsection{El morfismo canónico}
\label{sec:s08-morfismo-canonico}

\begin{defi}[El morfismo canónico $\eta_X$]
\label{defi:s08-morfismo-canonico}
Sea $X$ un esquema y $R=\Gamma(X,\mathcal O_X)$. El morfismo canónico $\eta_X:X\to\Spec R$ se define de la siguiente manera. Para $x\in X$, denotamos por $\mathfrak m_x$ el ideal maximal de $\mathcal O_{X,x}$ y por $\kappa(x)=\mathcal O_{X,x}/\mathfrak m_x$ su cuerpo residual.
\begin{enumerate}[label=(\roman*)]
\item En puntos,
\[
\eta_X(x)=\set*{a\in R\given a_x\in\mathfrak m_x}
=\ker\bigl(R\longrightarrow\mathcal O_{X,x}\longrightarrow\kappa(x)\bigr).
\]
Es un primo porque $\kappa(x)$ es un campo.
\item En abiertos, para $f\in R$,
\[
\eta_X^{-1}(D(f))=X_f
:=\set*{x\in X\given f_x\notin\mathfrak m_x},
\]
el abierto donde $f$ no se anula.
\item En gavillas, sobre cada básico,
\[
\eta_f^\sharp:\mathcal O_{\Spec R}(D(f))=R_f
\longrightarrow\mathcal O_X(X_f),
\qquad
\frac{a}{f^n}\longmapsto\frac{a|_{X_f}}{(f|_{X_f})^n}.
\]
\end{enumerate}
El conjunto $X_f$ es abierto porque un germen invertible tiene una inversa en algún entorno. Las inversas locales se pegan, de modo que $f|_{X_f}$ es invertible y la fórmula tiene sentido. En tallos, $a/s\mapsto a_x/s_x$ es un homomorfismo local. \emph{[Hartshorne, II, ejercicio 2.4; EGA I, 1.6.3.]}
\end{defi}

\begin{ejer}[Verificación del morfismo canónico]
\label{ejer:s08-verificacion-morfismo-canonico}
Sea $X$ un esquema y $R=\Gamma(X,\mathcal O_X)$. Muestra que $\eta_X:X\to\Spec R$ es un morfismo de esquemas verificando:
\begin{enumerate}[label=(\roman*)]
\item $\eta_X(x)$ es un ideal primo de $R$.
\item $\eta_X$ es continuo, pues $\eta_X^{-1}(D(f))=X_f$ es abierto.
\item $\eta_X^\sharp$ está bien definido: $f|_{X_f}$ es invertible en $\mathcal O_X(X_f)$ y los homomorfismos $R_f\to\mathcal O_X(X_f)$ son compatibles con las restricciones.
\item En tallos, $(\eta_X)_x^\sharp:R_{\eta_X(x)}\to\mathcal O_{X,x}$ es un homomorfismo local.
\end{enumerate}
\end{ejer}
% S08b, diapositivas 3--5: esta verificación se desarrolla en clase.
\begin{proof}[Desarrollo en clase]
\emph{Puntos y topología.} El ideal $\eta_X(x)$ es primo como núcleo de la evaluación en el cuerpo $\kappa(x)$; es propio porque $1$ se envía a $1\ne0$. Si $x\in X_f$, el germen $f_x$ es unidad. Elige un representante local $g$ de su inversa. La igualdad $g_xf_x=1$ vale en un entorno $V$ de $x$, pues dos secciones con el mismo germen coinciden cerca de $x$. Así $V\subseteq X_f$, luego $X_f$ es abierto. La igualdad $\eta_X^{-1}(D(f))=X_f$ y la base de abiertos $D(f)$ prueban la continuidad.

\emph{Gavillas.} Las inversas locales de $f$ en $X_f$ son únicas y por ello se pegan en una inversa de $f|_{X_f}$. La propiedad universal de la localización factoriza la restricción de manera única:
% S08b, diapositiva 4: localización y compatibilidad con restricciones.
\begin{center}
\begin{minipage}{.45\textwidth}
\centering
\begin{tikzcd}[column sep=large]
R\arrow[r,"\mathrm{res}"]\arrow[d,swap,"\mathrm{loc}"]
 &\mathcal O_X(X_f)\\
R_f\arrow[ur,dashed,swap,"\eta_f^\sharp"]&
\end{tikzcd}
\end{minipage}\hfill
\begin{minipage}{.51\textwidth}
\centering
\begin{tikzcd}[column sep=large]
R_f\arrow[r,"\eta_f^\sharp"]\arrow[d]
 &\mathcal O_X(X_f)\arrow[d,"\mathrm{res}"]\\
R_g\arrow[r,swap,"\eta_g^\sharp"]&\mathcal O_X(X_g)
\end{tikzcd}
\end{minipage}
\end{center}
Si $D(g)\subseteq D(f)$, también $X_g\subseteq X_f$. Ambos caminos son mapas $R_f\to\mathcal O_X(X_g)$ que coinciden en $R$. La unicidad de la localización da la igualdad. Por el \cref{teo:s05-toda-gavilla-sobre-base-extiende}, los mapas sobre los básicos definen un morfismo de gavillas. La compatibilidad con las restricciones es indispensable.

\emph{Tallos.} Sea $\mathfrak p=\eta_X(x)$. Si $s\notin\mathfrak p$, el germen $s_x$ es unidad; por tanto el mapa de gérmenes se factoriza por $R_{\mathfrak p}$:
\nopagebreak
% S08b, diapositiva 5: gérmenes, localización y cuerpo residual.
\begin{center}
\begin{tikzcd}[column sep=large]
R\arrow[r,"\text{germen}"]\arrow[d,swap,"\mathrm{loc}"]
 &\mathcal O_{X,x}\arrow[r]&\kappa(x)\\
R_{\mathfrak p}\arrow[ur,dashed,swap,"(\eta_X)_x^\sharp"]&&
\end{tikzcd}
\end{center}
El mapa inducido es
\[
(\eta_X)_x^\sharp:R_{\mathfrak p}\longrightarrow\mathcal O_{X,x},
\qquad \frac{a}{s}\longmapsto\frac{a_x}{s_x}.
\]
Como $s_x$ es unidad, la imagen de $a/s$ está en $\mathfrak m_x$ si y solo si $a_x\in\mathfrak m_x$, es decir, si y solo si $a\in\mathfrak p$. La preimagen de $\mathfrak m_x$ es $\mathfrak pR_{\mathfrak p}$: el homomorfismo es local. Así $\eta_X$ es un morfismo de espacios localmente anillados y, entre esquemas, un morfismo de esquemas.
\end{proof}

En una carta afín $V\cong\Spec B$, el morfismo es $\Spec(\rho_V)$ para la restricción
\[
\rho_V:R\longrightarrow\Gamma(V,\mathcal O_X)\cong B.
\]
En puntos manda $\mathfrak q$ a $\rho_V^{-1}(\mathfrak q)$ y en básicos induce las localizaciones de $\rho_V$, como en el \cref{teo:s06-todo-morfismo-anillos-b-induce}. Esta descripción muestra también que la construcción es independiente de la cubierta.

\begin{prop}[Qué dice el morfismo canónico]
\label{prop:s08-propiedades-morfismo-canonico}
Sea $X$ un esquema y $R=\Gamma(X,\mathcal O_X)$.
\begin{enumerate}[label=(\alph*)]
\item $\eta_X$ es un morfismo de espacios localmente anillados que en secciones globales es la identidad de $R$.
\item Para todo anillo $B$ y todo $\varphi:B\to R$ hay un único $g:X\to\Spec B$ que induce $\varphi$ en secciones globales, a saber $g=\Spec(\varphi)\circ\eta_X$. Es decir,
\[
\Hom(X,\Spec B)\cong\Hom_{\mathbf{Ring}}(B,\Gamma(X,\mathcal O_X)).
\]
\item Si $h:X\to Y$ es un morfismo, entonces
\[
\eta_Y\circ h=\Spec(\Gamma(h))\circ\eta_X.
\]
\item Si $\psi:X\to Y$ es un isomorfismo, $\eta_X$ es un isomorfismo si y solo si $\eta_Y$ lo es.
\item $X$ es afín si y solo si $\eta_X$ es un isomorfismo.
\end{enumerate}
\emph{[Hartshorne, II, ejercicio 2.4; EGA I, 1.6.3.]}
\end{prop}
\begin{proof}
\emph{La identidad en secciones globales.} La verificación está en el \cref{ejer:s08-verificacion-morfismo-canonico}. Con $f=1$, se tiene $X_1=X$ y el mapa $R\to\mathcal O_X(X)$ es la identidad; esto prueba (a).

\emph{Naturalidad y propiedad universal.} Para (c), sea $h:X\to Y$, con $h(x)=y$. Como el homomorfismo $h_x^\sharp:\mathcal O_{Y,y}\to\mathcal O_{X,x}$ es local, para $a\in\Gamma(Y,\mathcal O_Y)$,
\[
\Gamma(h)(a)\in\eta_X(x)
\iff h_x^\sharp(a_y)\in\mathfrak m_x
\iff a_y\in\mathfrak m_y
\iff a\in\eta_Y(y).
\]
Así $\Gamma(h)^{-1}(\eta_X(x))=\eta_Y(h(x))$, que da la igualdad en puntos. En gavillas, sobre $D(a)$ ambos caminos envían una fracción $b/a^n$ a la misma fracción de secciones restringidas, por la compatibilidad de $h^\sharp$ con restricciones y gérmenes. Luego los mapas de gavillas también coinciden.
% S08b, diapositiva 8: naturalidad general y con destino afín.
\begin{center}
\begin{minipage}{.48\textwidth}
\centering
\begin{tikzcd}[column sep=large]
X\arrow[r,"\eta_X"]\arrow[d,swap,"h"]
 &\Spec\Gamma(X,\mathcal O_X)\arrow[d,"\Spec\Gamma(h)"]\\
Y\arrow[r,swap,"\eta_Y"]&\Spec\Gamma(Y,\mathcal O_Y)
\end{tikzcd}
\end{minipage}\hfill
\begin{minipage}{.48\textwidth}
\centering
\begin{tikzcd}[column sep=large]
X\arrow[r,"\eta_X"]\arrow[d,swap,"g"]
 &\Spec\Gamma(X,\mathcal O_X)\arrow[d,"\Spec\Gamma(g)"]\\
\Spec B\arrow[r,swap,"\id"]&\Spec B
\end{tikzcd}
\end{minipage}
\end{center}
Aquí hemos identificado $\Gamma(\Spec B,\mathcal O)$ con $B$. En efecto, para $X=\Spec B$, la fórmula en puntos es
\[
\eta_{\Spec B}(\mathfrak p)
=\set*{b\in B\given b/1\in\mathfrak pB_{\mathfrak p}}
=\mathfrak p,
\]
y en gavillas $B_f\to B_f$ es la identidad. Por tanto $\eta_{\Spec B}=\id$.

Dado $\varphi:B\to R$, el compuesto $g=\Spec(\varphi)\circ\eta_X$ induce $\varphi$, por (a). Recíprocamente, la naturalidad aplicada a cualquier $g:X\to\Spec B$ da $g=\Spec(\Gamma(g))\circ\eta_X$. Por tanto el homomorfismo en secciones globales determina $g$, lo que prueba (b).
% S08b, diapositiva 7: factorización universal y correspondencia de Hom.
\begin{center}
\begin{minipage}{.47\textwidth}
\centering
\begin{tikzcd}[column sep=large,row sep=large]
X\arrow[r,"\eta_X"]\arrow[dr,swap,"g"]
 &\Spec R\arrow[d,dashed,"\exists!\,\Spec(\varphi)"]\\
&\Spec B
\end{tikzcd}
\end{minipage}\hfill
\begin{minipage}{.49\textwidth}
\centering
\begin{tikzcd}[row sep=large]
\Hom(X,\Spec B)\arrow[d,shift left=.8ex,"\Gamma"]\\
\Hom_{\mathbf{Ring}}(B,R)
 \arrow[u,shift left=.8ex,"\varphi\mapsto\Spec(\varphi)\circ\eta_X"]
\end{tikzcd}
\end{minipage}
\end{center}
Todo morfismo de $X$ a un esquema afín se factoriza de manera única por $\eta_X$. Es la unidad de la adjunción entre $\Gamma$ y $\Spec$: la correspondencia anterior envía $g$ a $\Gamma(g)$ y $\varphi$ a $\Spec(\varphi)\circ\eta_X$.

\emph{Isomorfismos y afinidad.} Si $\psi:X\to Y$ es un isomorfismo, también lo son $\Gamma(\psi)$ y $\Spec(\Gamma(\psi))$. El cuadrado de naturalidad tiene flechas verticales invertibles:
% S08b, diapositiva 9: transporte del morfismo canónico por un isomorfismo.
\begin{center}
\begin{tikzcd}[column sep=large]
X\arrow[r,"\eta_X"]\arrow[d,swap,"\psi","\sim"']
 &\Spec\Gamma(X,\mathcal O_X)
 \arrow[d,"\Spec\Gamma(\psi)","\sim"']\\
Y\arrow[r,swap,"\eta_Y"]&\Spec\Gamma(Y,\mathcal O_Y)
\end{tikzcd}
\end{center}
Por ello
\[
\eta_X=\Spec(\Gamma(\psi))^{-1}\circ\eta_Y\circ\psi,
\]
y $\eta_X$ es un isomorfismo si y solo si lo es $\eta_Y$, que es (d). En puntos, si $Y=\Spec B$ y $\psi(x)=\mathfrak p$, para $a=\Gamma(\psi)(b)$ se tiene $a_x=\psi_x^\sharp(b/1)$. Este isomorfismo local lleva $b/1$ al maximal si y solo si $b\in\mathfrak p$. Así $\eta_X(x)$ es el ideal transportado $\Gamma(\psi)(\mathfrak p)$.

Si $X$ es afín, elige $\psi:X\xrightarrow{\sim}\Spec B$. Como $\eta_{\Spec B}=\id$, (d) muestra que $\eta_X$ es un isomorfismo. A la inversa, si $\eta_X$ es un isomorfismo, identifica $X$ con el esquema afín $\Spec\Gamma(X,\mathcal O_X)$.

La presentación da también el argumento mediante Yoneda. Para todo esquema $Z$, conmuta el triángulo
% S08b, diapositiva 10: argumento de adjunción y Yoneda.
\begin{center}
\begin{tikzcd}[column sep=large,row sep=large]
\Hom(Z,X)\arrow[r,"\eta_X\circ-"]\arrow[dr,swap,"\Gamma"]
 &\Hom(Z,\Spec\Gamma(X,\mathcal O_X))
 \arrow[d,"\sim"]\\
&\Hom_{\mathbf{Ring}}(\Gamma(X,\mathcal O_X),\Gamma(Z,\mathcal O_Z)).
\end{tikzcd}
\end{center}
La flecha vertical es biyectiva por (b). Si $X$ es afín, la diagonal también lo es, aplicando (b) tras identificar $X$ con un espectro. Por tanto $\eta_X\circ-$ induce biyecciones en todos los conjuntos $\Hom(Z,-)$, y Yoneda implica que $\eta_X$ es un isomorfismo.
\end{proof}

\subsection{Aplicación al plano menos el origen}

\begin{obs}[El caso afín: identidad e inclusión]
\label{obs:s08-canonico-identidad-inclusion}
Para $X=\Spec A$, el morfismo canónico es la identidad bajo $\Gamma(X,\mathcal O_X)=A$. Si $U\subseteq\Spec A$ es abierto, con inclusión $j$ y restricción $\rho:A\to\Gamma(U,\mathcal O_U)$, la naturalidad da
\[
j=\Spec(\rho)\circ\eta_U.
\]
\end{obs}
% S08b, diapositiva 11: distinguir el canónico y la inclusión.
\begin{center}
\begin{tikzcd}[column sep=large,row sep=large]
U\arrow[r,"\eta_U"]\arrow[d,swap,"j"]
 &\Spec\Gamma(U,\mathcal O_U)\arrow[d,"\Spec(\rho)"]\\
\Spec A\arrow[r,swap,"\eta_{\Spec A}=\id"]&\Spec A
\end{tikzcd}
\end{center}
En puntos, para $\mathfrak p\in U$,
\[
\Spec(\rho)(\eta_U(\mathfrak p))
=\rho^{-1}\set*{g\in\Gamma(U,\mathcal O_U)\given g_{\mathfrak p}\in\mathfrak m_{\mathfrak p}}
=\mathfrak p.
\]
El germen de $a|_U$ en $\mathfrak p$ es el de $a$. Si $\rho$ es un isomorfismo, $\eta_U$ se identifica con la inclusión $j$ al identificar sus destinos mediante $\Spec(\rho)$. Antes de hacerlo, los destinos son distintos.

\begin{coro}
\label{coro:s08-plano-menos-origen-no-afin}
El abierto $U=\mathbb A_k^2\setminus\set{(x,y)}$ no es afín.
\end{coro}
\begin{proof}
Por la \cref{prop:s08-secciones-plano-menos-origen}, la restricción $\rho:k[x,y]\to\Gamma(U,\mathcal O_U)$ es un isomorfismo. La \cref{obs:s08-canonico-identidad-inclusion} da el cuadrado
% S08b, diapositiva 12: el canónico del plano perforado y la inclusión.
\begin{center}
\begin{tikzcd}[column sep=large,row sep=large]
U\arrow[r,"\eta_U"]\arrow[d,swap,"j"]
 &\Spec\Gamma(U,\mathcal O_U)\arrow[d,"\Spec(\rho)","\sim"']\\
\mathbb A_k^2\arrow[r,swap,"\eta_{\mathbb A_k^2}=\id"]&\mathbb A_k^2
\end{tikzcd}
\end{center}
Si $U$ fuera afín, $\eta_U$ sería un isomorfismo por la \cref{prop:s08-propiedades-morfismo-canonico}(e). Como $\Spec(\rho)$ lo es, también lo sería $j=\Spec(\rho)\circ\eta_U$. Pero $j$ omite el origen $(x,y)$ y no es sobreyectiva.

En la identificación $\Gamma(U,\mathcal O_U)=A=k[x,y]$, la fórmula es
\[
\eta_U(\mathfrak p)
=\set*{a\in A\given a_{\mathfrak p}\in\mathfrak pA_{\mathfrak p}}
=\mathfrak p,
\]
pues $a_{\mathfrak p}\in\mathfrak pA_{\mathfrak p}$ equivale a $a\in\mathfrak p$. Las funciones globales de $U$ son las del plano entero, cuyo espectro recupera el origen ausente.
\end{proof}

\section{La recta proyectiva por cartas}

Con la técnica anterior, calculemos secciones globales de un esquema construido por pegado. Primero construiremos sus cartas y sus transiciones.

\begin{defi}[La recta proyectiva $\mathbb P_A^1$]
\label{defi:s08-recta-proyectiva}
Sea $A$ un anillo. Sean $U_0=\Spec A[t]$ y $U_1=\Spec A[s]$, dos copias de la recta afín $\mathbb A_A^1$, llamadas cartas. Los abiertos $U_{01}=D(t)\subseteq U_0$ y $U_{10}=D(s)\subseteq U_1$ tienen anillos de funciones $A[t,t^{-1}]$ y $A[s,s^{-1}]$. El isomorfismo de anillos
\[
A[s,s^{-1}]\longrightarrow A[t,t^{-1}],\qquad s\longmapsto t^{-1},
\]
define un isomorfismo de esquemas $\varphi_{01}:U_{01}\to U_{10}$. Por la \cref{defi:s08-datos-pegado}, son datos de pegado. La recta proyectiva $\mathbb P_A^1$ es el esquema del \cref{teo:s08-lema-pegado}, con inmersiones abiertas $\psi_0:U_0\to\mathbb P_A^1$ y $\psi_1:U_1\to\mathbb P_A^1$. \emph{[Vakil, \S4.4.6 y \S4.4.10; Hartshorne, II.2, ejercicio 2.12.]}
\end{defi}

\begin{obs}[Coordenadas]
\label{obs:s08-coordenadas}
Un punto $[x_0:x_1]$ con $x_0\ne0$ tiene coordenada $t=x_1/x_0$; cuando $x_1\ne0$, tiene coordenada $s=x_0/x_1$. En el traslape se cumple $s=t^{-1}$, la transición de la \cref{defi:s08-recta-proyectiva}. Si $A=k$ y $a\in k^\times$, los puntos $t=a$ y $s=a^{-1}$ de ambas cartas se identifican.
\end{obs}

% Presentación S08, diapositiva 13: círculo y las dos proyecciones.
\begin{center}
\begin{minipage}{.47\textwidth}
\centering
\begin{tikzpicture}[scale=1,every node/.style={font=\small}]
\draw[grisCIMAT,thick] (0,0) circle (1.2);
\draw[blue!60!black,thick] (-2.4,-1.2)--(2.4,-1.2) node[right] {$t$};
\draw[green!45!black,thick] (-2.4,1.2)--(2.4,1.2) node[right] {$s$};
\begin{scope}
\clip (-2.4,-1.2) rectangle (2.4,1.2);
\foreach \x/\y in {-.98/.69,.98/.69,-.98/-.69,.98/-.69}{
 \pgfmathsetmacro{\projt}{2.4*\x/(1.2-\y)}
 \pgfmathsetmacro{\projs}{2.4*\x/(1.2+\y)}
 \draw[blue!60!black,dashed] (0,1.2)--(\projt,-1.2);
 \draw[green!45!black,dashed] (0,-1.2)--(\projs,1.2);
}
\end{scope}
\foreach \x/\y in {-.98/.69,.98/.69,-.98/-.69,.98/-.69}{
 \fill[rojoCIMAT] (\x,\y) circle (1.8pt);
}
\fill[rojoCIMAT] (0,1.2) circle (2.5pt);
\fill[rojoCIMAT] (0,-1.2) circle (2.5pt);
\node[above] at (0,1.2) {$s=0$};
\node[below] at (0,-1.2) {$t=0$};
\end{tikzpicture}
\end{minipage}\hfill
\begin{minipage}{.49\textwidth}
\small
La carta $t$ proyecta desde $s=0$ sobre la recta tangente en $t=0$. La carta $s$ proyecta desde $t=0$ sobre la tangente en $s=0$. El traslape es el círculo menos los dos polos, donde $ts=1$. Sobre $\mathbb R$ se ve un círculo; sobre $\mathbb C$, la esfera de Riemann.
\end{minipage}
\end{center}

\begin{obs}[El morfismo estructural]
\label{obs:s08-morfismo-estructural}
Los morfismos $U_0\to\Spec A$ y $U_1\to\Spec A$ inducidos por $A\to A[t]$ y $A\to A[s]$ coinciden en el traslape, pues $s\mapsto t^{-1}$ es un homomorfismo de $A$-álgebras. La \cref{obs:s08-morfismos-desde-pegado} da $\pi:\mathbb P_A^1\to\Spec A$; sus fibras se estudian en el \cref{sec:s08-puntos}.
\end{obs}

% S08b, diapositiva 15: propiedad universal del pegado y triángulo de anillos.
Sean $\iota_0:U_{01}\hookrightarrow U_0$ y $\iota_1:U_{10}\hookrightarrow U_1$. La condición para pegar $f_0:U_0\to Y$ y $f_1:U_1\to Y$ es $f_0\circ\iota_0=f_1\circ\iota_1\circ\varphi_{01}$. La \cref{obs:s08-morfismos-desde-pegado} da entonces un único $f:\mathbb P_A^1\to Y$:
\begin{center}
\begin{minipage}{.59\textwidth}
\centering
\begin{tikzcd}[column sep=large,row sep=large]
U_{01}\arrow[r,"\iota_0"]\arrow[d,swap,"\iota_1\circ\varphi_{01}"]
 &U_0\arrow[d,"\psi_0"]\arrow[ddr,bend left=25,"f_0"]&\\
U_1\arrow[r,swap,"\psi_1"]\arrow[drr,bend right=20,swap,"f_1"]
 &\mathbb P_A^1\arrow[dr,dashed,"\exists!\,f"]&\\
&&Y
\end{tikzcd}
\end{minipage}\hfill
\begin{minipage}{.37\textwidth}
\centering
\begin{tikzcd}[column sep=large,row sep=large]
A\arrow[r]\arrow[dr,swap]
 &A[s,s^{-1}]\arrow[d,"s\mapsto t^{-1}"]\\
&A[t,t^{-1}]
\end{tikzcd}
\end{minipage}
\end{center}
Para $Y=\Spec A$, el triángulo de anillos verifica la compatibilidad de los morfismos estructurales: la transición es un homomorfismo de $A$-álgebras. De este pegado se obtiene $\pi$.

\begin{prop}[$\mathbb P_A^1$ como conjunto]
\label{prop:s08-proyectiva-como-conjunto}
Se tiene $\mathbb P_A^1=\psi_0(U_0)\amalg\psi_1(V(s))$, con
\[
V(s)=\Spec(A[s]/(s))\cong\Spec A.
\]
Así, $\mathbb P_A^1$ es la recta afín $\mathbb A_A^1$ más una copia de $\Spec A$, la sección del infinito. Esta sección es un punto cuando $A$ es un campo.
\end{prop}
\begin{proof}
Por (a) y (b) del \cref{teo:s08-lema-pegado}, el complemento de $\psi_0(U_0)$ es
\[
\psi_1(U_1\setminus D(s))=\psi_1(V(s)).
\]
El cociente $A[s]/(s)\cong A$ identifica este conjunto con $\Spec A$.
\end{proof}

% Presentación S08, diapositiva 14: mismas piezas, dos pegados.
\begin{center}
\begin{tikzpicture}[every node/.style={font=\small}]
\begin{scope}
\node at (0,1.7) {recta con origen doble};
\draw[blue!60!black,thick] (-2.2,1)--(2.2,1) node[right] {$X_1$};
\draw[green!45!black,thick] (-2.2,0)--(2.2,0) node[right] {$X_2$};
\foreach \x in {-1.8,-1.2,-.6,.6,1.2,1.8}{\draw[grisCIMAT,dotted,semithick] (\x,0)--(\x,1);}
\fill[rojoCIMAT] (0,1) circle (2.5pt) node[above] {$O_1$};
\fill[rojoCIMAT] (0,0) circle (2.5pt) node[below] {$O_2$};
\node at (0,-.9) {$\varphi^\sharp(y)=x$};
\end{scope}
\begin{scope}[xshift=7cm]
\node at (0,1.7) {$\mathbb P_k^1$};
\draw[blue!60!black,thick] (-2.2,1)--(2.2,1) node[right] {$t$};
\draw[green!45!black,thick] (-2.2,0)--(2.2,0) node[right] {$s$};
\draw[grisCIMAT,dotted,semithick] (-2,1)--(-.4,0) (-.4,1)--(-2,0) (.4,1)--(2,0) (2,1)--(.4,0);
\fill[rojoCIMAT] (0,1) circle (2.5pt) node[above] {$t=0$};
\fill[rojoCIMAT] (0,0) circle (2.5pt) node[below] {$s=0$};
\node at (0,-.9) {$\varphi^\sharp(y)=x^{-1}$};
\end{scope}
\end{tikzpicture}
\end{center}
Las piezas $\mathbb A_k^1$ y el traslape $D(x)$ son los mismos. Con la identidad, cerca de $O_1$ y de $O_2$ aparecen los mismos puntos. Con la inversión, cerca de $t=0$ en una carta es lejos en la otra. El morfismo de pegado distingue las dos construcciones.
\clearpage
\section{Secciones globales y no afinidad}

\begin{prop}
\label{prop:s08-secciones-proyectiva}
Para todo anillo $A$, se tiene $\Gamma(\mathbb P_A^1,\mathcal O)=A$: el homomorfismo $A\to\Gamma(\mathbb P_A^1,\mathcal O)$ inducido por $\pi$ es un isomorfismo. \emph{[Vakil, \S4.4.8.]}
\end{prop}
\begin{proof}
Una sección global es un par $(f,g)\in A[t]\times A[s]$ con $f=g(t^{-1})$ en $A[t,t^{-1}]$, por la \cref{defi:s07-gavilla-estructural-pegado}. Es el igualador de las dos restricciones:
% S08b, diapositiva 19: las secciones globales como igualador.
\begin{center}
\begin{tikzcd}[column sep=huge]
\Gamma(\mathbb P_A^1,\mathcal O)\arrow[r,hook]
 &A[t]\times A[s]
 \arrow[r,shift left=.8ex,"{(f,g)\mapsto f}"]
 \arrow[r,shift right=.8ex,swap,"{(f,g)\mapsto g(t^{-1})}"]
 &A[t,t^{-1}]
\end{tikzcd}
\end{center}
Escribamos
\[
f=\sum_{n\geq0}a_nt^n,\qquad g(t^{-1})=\sum_{n\geq0}b_nt^{-n}.
\]
El anillo de Laurent es libre como $A$-módulo con base $\set{t^n}_{n\in\mathbb Z}$. Comparar coeficientes da $a_n=b_n=0$ para $n>0$ y $a_0=b_0$. Luego $f=g=a_0\in A$. Cada $a\in A$ da, a su vez, la sección compatible $(a,a)$.

\emph{Advertencia.} $A[t]\hookrightarrow A[t,t^{-1}]$ es inyectivo para cualquier $A$, sin suponer que sea un dominio.
\end{proof}

\begin{coro}
\label{coro:s08-proyectiva-no-afin}
Si $A\ne0$, la recta proyectiva $\mathbb P_A^1$ no es afín.
\end{coro}
\begin{proof}
Sea $R=\Gamma(\mathbb P_A^1,\mathcal O)$ y $\varphi:A\to R$ el isomorfismo de la \cref{prop:s08-secciones-proyectiva}. Por la \cref{prop:s08-propiedades-morfismo-canonico}(b), aplicada a $g=\pi$, se tiene
\[
\pi=\Spec(\varphi)\circ\eta_{\mathbb P_A^1}.
\]
% S08b, diapositiva 21: identificar pi con el morfismo canónico.
\begin{center}
\begin{tikzcd}[column sep=large,row sep=large]
\mathbb P_A^1\arrow[r,"\eta_{\mathbb P_A^1}"]\arrow[dr,swap,"\pi"]
 &\Spec R\arrow[d,"\Spec(\varphi)","\sim"']\\
&\Spec A
\end{tikzcd}
\end{center}
Si $\mathbb P_A^1$ fuera afín, su morfismo canónico sería un isomorfismo por la \cref{prop:s08-propiedades-morfismo-canonico}(e). Como $\Spec(\varphi)$ también lo es, $\pi$ sería un isomorfismo y, en particular, biyectivo.

Como $A\ne0$, existe un primo $\mathfrak r\subseteq A$. En la carta $U_0$ consideremos
\[
P_0=\mathfrak r A[t]+(t),\qquad
P_1=\mathfrak r A[t]+(t-1).
\]
Ambos son primos de la misma carta $U_0$, pues $A[t]/P_i\cong A/\mathfrak r$. Son distintos: la imagen de $t-1$ en $A[t]/P_0$ es $-1\ne0$. La inmersión abierta $\psi_0$ es inyectiva, de modo que también son distintos en $\mathbb P_A^1$; no hace falta cambiar de carta. Además, $P_i\cap A=\mathfrak r$. En $U_0$, el morfismo $\pi$ es $\Spec$ de la inclusión $A\to A[t]$, por lo que envía cada primo a su contracción. Así,
\[
\pi(P_0)=\mathfrak r=\pi(P_1),
\]
y $\pi$ no es inyectivo. En términos de gérmenes, $a\in A$ pertenece al maximal del tallo en $P_i$ si y solo si $a\in\mathfrak r$. Esto confirma la fórmula del morfismo canónico.
\end{proof}

% Presentación S08, diapositiva 18: no sobreyectividad y no inyectividad.
\begin{center}
\begin{tikzpicture}[every node/.style={font=\small}]
\fill[orange!12] (0,0) rectangle (1.6,1.6);
\draw[grisCIMAT] (0,0) rectangle (1.6,1.6);
\draw[rojoCIMAT,thick,fill=white] (.8,.8) circle (2.5pt);
\draw[->,blue!60!black,thick] (1.95,.8)--(3.15,.8) node[midway,above] {$\eta_U$};
\fill[orange!12] (3.5,0) rectangle (5.1,1.6);
\draw[grisCIMAT] (3.5,0) rectangle (5.1,1.6);
\fill[rojoCIMAT] (4.3,.8) circle (2.5pt) node[above right,xshift=-2.5pt,yshift=-2.5pt] {$\scriptstyle(x,y)$};
\node at (.8,-.4) {$U$};
\node at (4.3,-.4) {$\mathbb A_k^2=\Spec\Gamma(U,\mathcal O_U)$};
\node[align=center] at (2.55,-1.25) {\footnotesize Falta el origen: el mapa no es sobreyectivo.};
\begin{scope}[xshift=11cm]
\draw[green!45!black,thick] (0,.8)--(0,2.3);
\node[left] at (0,1.6) {$\mathbb P_A^1$};
\foreach \y/\lab in {2.1/P_0,1.6/P_1,1.1/\infty}{\fill[rojoCIMAT] (0,\y) circle (2.5pt) node[right] {$\lab$};}
\draw[->,blue!60!black,thick] (0,.7)--(0,.05) node[midway,left] {$\pi$};
\draw[grisCIMAT] (-1.4,-.2)--(1.4,-.2) node[right] {$\Spec A$};
\fill[rojoCIMAT] (0,-.2) circle (2.5pt) node[below] {$\mathfrak r$};
\node[align=center] at (0,-1.25) {\footnotesize $P_0\ne P_1$ tienen misma imagen: el mapa no es inyectivo.};
\end{scope}
\end{tikzpicture}
\end{center}
Un esquema afín se reconstruye con sus funciones. En $U$ las funciones no detectan el hueco del origen; en $\mathbb P_A^1$ no distinguen los dos puntos elegidos.

\begin{obs}
\label{obs:s08-constantes-compacidad}
Es el análogo de que las únicas funciones holomorfas en $\mathbb{CP}^1$ son constantes; en el caso complejo, esto se explica por la compacidad. Una función entera no constante tiene en $\infty$ un polo si es un polinomio y una singularidad esencial si no lo es, como $e^z$. Si fuera holomorfa también en $\infty$, sería constante por Liouville. La versión algebraica de la compacidad es la propiedad, que se estudiará en la Sesión 15. La recta con origen doble tiene secciones globales $k[x]$, que la distinguen de $\mathbb P_k^1$; véase el \cref{sec:s08-dos-rectas}. \emph{[Vakil, ejercicio 4.4.B.]}
\end{obs}

\section{Los puntos de la recta proyectiva}
\label{sec:s08-puntos}

\begin{obs}[Qué buscamos]
\label{obs:s08-objetivos-puntos}
Describiremos los puntos de la recta proyectiva sobre un cuerpo, sobre $\mathbb Z$ y sobre un anillo de valoración discreta. Las cartas reducen el problema a clasificar primos de anillos de polinomios. Después estudiaremos la reducción de puntos racionales a una fibra especial, construiremos $\mathbb P_A^n$ y clasificaremos los pegados de dos rectas.
\end{obs}

\subsection{Lo que se usa hoy}

\subsubsection{Definiciones y resultados previos}
Recordemos las \cref{defi:s01-espectro-primo,defi:s02-s-definimos-v-s-p,defi:s02-f-abierto-distinguido-asociado-f}: los puntos de $\Spec A$ son sus ideales primos, los cerrados son $V(I)=\set{\mathfrak p\in\Spec A\given I\subseteq\mathfrak p}$ y los abiertos básicos son $D(f)=\set{\mathfrak p\given f\notin\mathfrak p}$. Por el \cref{coro:s05-gavilla-estructural}, $\mathcal O(D(f))=A_f$; además, $D(f)\cong\Spec A_f$ por el \cref{coro:s03-inclusion-inducida-identifica-f-d}.

Un punto $x$ es cerrado cuando $\set{x}$ es cerrado, y es genérico de un cerrado $Z$ cuando $\overline{\set{x}}=Z$. Por la \cref{prop:s02-todo-p-cerradura-punto-p}, $\overline{\set{\mathfrak p}}=V(\mathfrak p)$: los puntos cerrados del espectro son los ideales máximos y, si $A$ es un dominio, $(0)$ es su punto genérico. La \cref{defi:s02-decimos-q-especializacion-p-q} describe las especializaciones.

\begin{obs}[La especialización es local]
\label{obs:s08-especializacion-local}
Si $y\in\overline{\set{x}}$ y un abierto $U$ contiene a $y$, contiene también a $x$: toda vecindad de $y$ corta a $\set{x}$. Así, $x\notin U$ implica $\overline{\set{x}}\cap U=\varnothing$.
\end{obs}

Por la \cref{obs:s03-cuerpo-residual}, el cuerpo residual de $\mathfrak p$ es
\[
\kappa(\mathfrak p)=A_{\mathfrak p}/\mathfrak pA_{\mathfrak p}=\Frac(A/\mathfrak p).
\]
En un esquema, $\kappa(x)=\mathcal O_{X,x}/\mathfrak m_x$, y una carta afín conserva este cuerpo. Un homomorfismo $\varphi:A\to B$ induce $\Spec\varphi$, dado por $\mathfrak q\mapsto\varphi^{-1}(\mathfrak q)$; véanse el \cref{lema:s01-preimagen-primo} y el \cref{ejer:s02-b-morfismo-anillos-demuestra-1}. Además,
\[
(\Spec\varphi)^{-1}(D(f))=D(\varphi(f)).
\]

Una inmersión abierta identifica su fuente con un subesquema abierto de la meta; en particular, es inyectiva en puntos. Las cartas de la \cref{defi:s08-recta-proyectiva} son inmersiones abiertas. Su morfismo estructural se construyó en la \cref{obs:s08-morfismo-estructural}; para sus secciones y su no afinidad, véanse \cref{prop:s08-secciones-proyectiva,coro:s08-proyectiva-no-afin}. El criterio de afinidad y el plano menos el origen se trataron en \cref{prop:s08-propiedades-morfismo-canonico,obs:s08-canonico-identidad-inclusion,prop:s08-secciones-plano-menos-origen,coro:s08-plano-menos-origen-no-afin}.

\begin{defi}[Esquema sobre $A$, morfismo estructural y fibra]
\label{defi:s08-esquema-sobre-anillo-fibra}
Un esquema sobre $A$ es un esquema $X$ con un morfismo $\pi:X\to\Spec A$, llamado morfismo estructural. Para $\mathbb A_A^1=\Spec A[t]$, la inclusión $A\to A[t]$ induce $\pi(\mathfrak q)=\mathfrak q\cap A$. La fibra sobre $y\in\Spec A$ es el conjunto $\pi^{-1}(y)$. Aquí solo se usa como conjunto de puntos; la fibra como esquema, $X\times_{\Spec A}\Spec\kappa(y)$, se construirá en las Sesiones 12 y 14.
\end{defi}

\begin{defi}[Grado de un punto]
\label{defi:s08-grado-punto}
Sea $k$ un campo y $x$ un punto de un esquema sobre $k$ cuyo cuerpo residual es una extensión finita de $k$. Su grado es $[\kappa(x):k]$, donde la estructura de $k$-álgebra viene del morfismo estructural. Un punto de grado uno se llama racional.
\end{defi}

\begin{defi}[Polinomios mónicos, irreducibles y primitivos]
\label{defi:s08-polinomios-primitivos}
Sea $B$ un dominio. Un polinomio de $B[t]$ es mónico si su coeficiente principal es uno. Es irreducible si no es una unidad y no es producto de dos no unidades. El contenido de un polinomio no nulo de $\mathbb Z[t]$ es el máximo común divisor de sus coeficientes; el polinomio es primitivo si su contenido es uno.
\end{defi}

\subsubsection{Hechos de álgebra}

\begin{prop}[Hecho A, polinomios sobre un campo]
\label{prop:s08-polinomios-campo}
Si $k$ es un campo, $k[t]$ es un dominio de ideales principales. Sus primos son $(0)$ y los $(f)$, con $f$ irreducible; estos últimos son máximos y tienen un único generador mónico. Si $\deg f=d$, el campo $k[t]/(f)$ tiene $k$-base $1,t,\ldots,t^{d-1}$. En cualquier dominio de ideales principales, todo primo no nulo es máximo. \emph{[Lang, Algebra, cap.~IV, \S1; Dummit--Foote, \S8.2.]}
\end{prop}

\begin{prop}[Hecho B, primos de una localización]
\label{prop:s08-primos-localizacion}
Sea $B$ un anillo y $S\subseteq B$ multiplicativo. La aplicación $\mathfrak q\mapsto S^{-1}\mathfrak q$ es una biyección entre los primos de $B$ disjuntos de $S$ y $\Spec S^{-1}B$. Su inversa es la contracción por $B\to S^{-1}B$. Es el \cref{teo:s03-correspondencia-localizaciones}. \emph{[Atiyah--Macdonald, prop.~3.11(iv).]}
\end{prop}

\begin{prop}[Hecho C, primos de un cociente]
\label{prop:s08-primos-cociente}
Sea $B$ un anillo e $I$ un ideal. La aplicación $\mathfrak q\mapsto\mathfrak q/I$ es una biyección entre los primos de $B$ que contienen a $I$ y $\Spec(B/I)$. Su inversa es la contracción por $B\to B/I$; véase el \cref{ejer:s01-i-cociente-demuestra-i-inyectivo}. \emph{[Atiyah--Macdonald, cap.~1.]}
\end{prop}

\begin{prop}[Hecho D, lema de Gauss]
\label{prop:s08-gauss}
El producto de dos polinomios primitivos de $\mathbb Z[t]$ es primitivo. \emph{[Lang, Algebra, cap.~IV, \S2.]} En consecuencia:
\begin{enumerate}[label=(G\arabic*)]
\item Todo irreducible de $\mathbb Q[t]$ de grado positivo es, salvo un factor de $\mathbb Q^\times$, un polinomio primitivo $g\in\mathbb Z[t]$, único salvo signo e irreducible en $\mathbb Z[t]$. Recíprocamente, un primitivo de grado positivo irreducible en $\mathbb Z[t]$ es irreducible en $\mathbb Q[t]$. Como $\mathbb Z[t]$ es un dominio de factorización única, $(g)$ es primo.
\item Si $g\in\mathbb Z[t]$ es primitivo y divide a $h\in\mathbb Z[t]$ en $\mathbb Q[t]$, también lo divide en $\mathbb Z[t]$; equivalentemente,
\[
g\mathbb Q[t]\cap\mathbb Z[t]=g\mathbb Z[t].
\]
\end{enumerate}
\end{prop}

\begin{prop}[Hecho E, cuerpos finitos]
\label{prop:s08-cuerpos-finitos}
Sea $K$ un campo de característica $p>0$ y $q$ una potencia de $p$.
\begin{enumerate}[label=(E\arabic*)]
\item El Frobenius $x\mapsto x^q$ es un homomorfismo de anillos $K\to K$.
\item Si $K$ tiene $q$ elementos, todo $x\in K$ cumple $x^q=x$.
\item Si $K$ es finito, el grupo multiplicativo $K^\times$ es cíclico.
\end{enumerate}
\emph{[Lang, Algebra, cap.~V, \S5; Dummit--Foote, \S14.3.]}
\end{prop}

\begin{lema}[Fibra de la recta afín]
\label{lema:s08-fibra-recta-afin}
Sean $A$ un anillo, $\mathfrak p\in\Spec A$ y $\kappa=\kappa(\mathfrak p)$. La fibra sobre $\mathfrak p$ de la contracción $\Spec A[t]\to\Spec A$ está en biyección con $\Spec\kappa[t]$. La biyección lleva $\mathfrak r\in\Spec\kappa[t]$ a su imagen inversa por $A[t]\to\kappa[t]$. Lo mismo vale con $s$ en lugar de $t$.
\end{lema}
\begin{proof}
Localizamos primero en $S=A\setminus\mathfrak p$. Los primos $\mathfrak q$ de $A[t]$ con $\mathfrak q\cap A=\mathfrak p$ corresponden a los primos de $A_{\mathfrak p}[t]$ que contienen a $\mathfrak pA_{\mathfrak p}[t]$, por la \cref{prop:s08-primos-localizacion}. Al reducir por ese ideal, la \cref{prop:s08-primos-cociente} los identifica con los primos de
\[
A_{\mathfrak p}[t]/\mathfrak pA_{\mathfrak p}[t]=\kappa[t].
\]
Recíprocamente, un primo que contiene a $\mathfrak pA_{\mathfrak p}[t]$ contrae a $\mathfrak pA_{\mathfrak p}$ en $A_{\mathfrak p}$, porque este ideal es máximo. Su contracción a $A$ es $\mathfrak p$. Las dos correspondencias inversas son contracciones y su compuesto es el mapa anunciado.
\end{proof}

\begin{lema}[Cerrados y cartas]
\label{lema:s08-cerrados-cartas}
Si $X$ está cubierto por abiertos $U_i$, un subconjunto $Z$ es cerrado si y solo si $Z\cap U_i$ es cerrado en $U_i$ para todo $i$. En particular, un punto es cerrado si y solo si es cerrado en cada carta que lo contiene.
\end{lema}
\begin{proof}
Una intersección de un cerrado con $U_i$ es cerrada en $U_i$. Recíprocamente, si todas las intersecciones son cerradas, cada $U_i\setminus Z$ es abierto en $U_i$ y en $X$; su unión es $X\setminus Z$. Para un punto, las cartas que no lo contienen tienen intersección vacía con él.
\end{proof}

\subsection{Sobre un cuerpo}

\begin{prop}[Los puntos de $\mathbb A_k^1$]
\label{prop:s08-puntos-recta-afin-campo}
Sea $k$ un campo. Los puntos de $\mathbb A_k^1=\Spec k[t]$ son $(0)$ y los $(f)$ con $f$ mónico irreducible, uno por cada tal polinomio. \emph{[Vakil, \S4.4.6; \cref{prop:s01-recta-afin}.]}
\begin{enumerate}[label=(\alph*)]
\item $(0)$ es el punto genérico, no es cerrado y $\kappa((0))=k(t)$.
\item Cada $(f)$ es cerrado, con $\kappa((f))=k[t]/(f)$ y grado $\deg f$.
\item Los puntos de grado uno son los $(t-a)$, con $a\in k$.
\end{enumerate}
\end{prop}
\begin{proof}
La \cref{prop:s08-polinomios-campo} lista los primos y sus cocientes. Como $k[t]$ es un dominio, la cerradura de $(0)$ es todo el espectro; no es máximo, pues $(0)\subsetneq(t)$. Los demás primos son máximos y sus cocientes ya son campos, de dimensión $\deg f$ sobre $k$. Un mónico de grado uno es precisamente $t-a$.
\end{proof}

\begin{prop}[Los puntos de $\mathbb P_k^1$]
\label{prop:s08-puntos-proyectiva-cuerpo}
\label{ejem:s08-puntos-proyectiva-cuerpo}
Sea $k$ un campo e $\infty=\psi_1((s))$. Por la \cref{prop:s08-proyectiva-como-conjunto},
\[
\mathbb P_k^1=\mathbb A_k^1\amalg\set{\infty}.
\]
Sus puntos son el genérico $(0)$, los puntos cerrados $(f)$ con $f\in k[t]$ mónico irreducible y el punto cerrado $\infty$. Se tiene $\kappa((0))=k(t)$, $\kappa((f))=k[t]/(f)$, de grado $\deg f$, y $\kappa(\infty)=k$, de grado uno. \emph{[Vakil, \S4.4.6 y \S4.4.7.]}
\end{prop}
\begin{proof}
La lista procede de las \cref{prop:s08-proyectiva-como-conjunto,prop:s08-puntos-recta-afin-campo}. El genérico $(0)$ está en $D(t)$ y corresponde al genérico $(0)$ de la segunda carta; su cerradura contiene ambas cartas. No es cerrado, pues ya no lo es en $U_0$.

Para $(f)\ne(t)$, la localización por $t$ da un ideal máximo de $k[t,t^{-1}]$. En la segunda carta corresponde a un irreducible no divisible por $s$: al sustituir $t=s^{-1}$ y limpiar denominadores se obtiene un polinomio de término constante no nulo. Así el punto es cerrado en ambas cartas. El punto $(t)$ no está en la segunda carta y $\infty=(s)$ no está en la primera. El \cref{lema:s08-cerrados-cartas} prueba su cerradura global. Los cuerpos residuales se calculan en las cartas afines.
\end{proof}

\begin{coro}[Los puntos cerrados de $\mathbb P_{\mathbb C}^1$ y de $\mathbb P_{\mathbb R}^1$]
\label{coro:s08-puntos-complejos-reales}
\begin{enumerate}[label=(\alph*)]
\item Sobre $\mathbb C$, los puntos cerrados son los $(t-z)$, con $z\in\mathbb C$, e $\infty$, todos de grado uno. Como conjunto son $\mathbb C\cup\set{\infty}$, la esfera de Riemann.
\item Sobre $\mathbb R$, son los $(t-a)$ con $a\in\mathbb R$ y $\infty$, de grado uno, y los $(t^2+bt+c)$ con $b^2<4c$, de grado dos y cuerpo residual $\mathbb C$. Asociar a cada polinomio sus raíces complejas, y a $\infty$ el conjunto $\set{\infty}$, da una biyección con las órbitas de $\operatorname{Gal}(\mathbb C/\mathbb R)=\set{\id,z\mapsto\bar z}$ en $\mathbb C\cup\set{\infty}$.
\end{enumerate}
\emph{[Vakil, \S4.4.7.]}
\end{coro}
\begin{proof}
Por el teorema fundamental del álgebra, los irreducibles complejos son lineales y los reales son lineales o cuadráticos de discriminante negativo. En este último caso, evaluar en una raíz no real identifica el cociente con $\mathbb C$. Las raíces forman un par conjugado; las raíces reales y el infinito son órbitas de un solo elemento. Aplicamos la \cref{prop:s08-puntos-proyectiva-cuerpo}.
\end{proof}

% Presentación S08, diapositiva 21: órbita de la conjugación.
\begin{center}
\begin{tikzpicture}[every node/.style={font=\small}]
\draw[grisCIMAT] (-2.5,0)--(2.5,0) node[right] {$\mathbb R$};
\draw[grisCIMAT] (0,-1.5)--(0,1.5);
\foreach \x in {-1.8,.7,1.8}{\fill[rojoCIMAT] (\x,0) circle (2pt);}
\fill[rojoCIMAT] (-.8,1.1) circle (2.5pt) node[left] {$z$};
\fill[rojoCIMAT] (-.8,-1.1) circle (2.5pt) node[left] {$\bar z$};
\draw[->,blue!60!black,dashed,thick] (-.8,1.05) to[bend left=25] (-.8,-1.05);
\node[right,blue!60!black] at (.25,.7) {$\operatorname{Gal}(\mathbb C/\mathbb R)$};
\end{tikzpicture}
\end{center}
\emph{Advertencia.} Un punto cerrado de grado dos es el par $\set{z,\bar z}$, no dos puntos. El polinomio $t^2+1$ da un solo punto de $\mathbb P_{\mathbb R}^1$, aunque $i$ y $-i$ son dos números.

\begin{lema}[Factorización de $t^{q^n}-t$]
\label{lema:s08-factorizacion-frobenius}
Sea $q$ una potencia de un primo y $n\geq1$. En $\mathbb F_q[t]$, el polinomio $t^{q^n}-t$ es el producto de todos los mónicos irreducibles cuyo grado divide a $n$, cada uno una sola vez.
\end{lema}
\begin{proof}
La derivada de $P=t^{q^n}-t$ es $-1$, por lo que no hay factores repetidos. Sea $f$ mónico irreducible de grado $d$, $K=\mathbb F_q[t]/(f)$ y $\alpha$ la clase de $t$. Por la \cref{prop:s08-polinomios-campo}, $K$ tiene $q^d$ elementos. El Frobenius $F(x)=x^q$ fija $\mathbb F_q$, y $F^d=\id_K$ por la \cref{prop:s08-cuerpos-finitos}.

Si $d\mid n$, resulta $\alpha^{q^n}=\alpha$ y $f\mid P$. Recíprocamente, si $f\mid P$, entonces $F^n(\alpha)=\alpha$. Como $\alpha$ genera $K$ sobre $\mathbb F_q$, $F^n=\id_K$. La ciclicidad de $K^\times$ implica que $q^d-1$ divide a $q^n-1$. Escribamos $n=ad+r$, con $0\leq r<d$. Módulo $q^d-1$, tenemos $q^n-1\equiv q^r-1$; si $r>0$, este es un entero positivo menor que $q^d-1$, contradicción. Por tanto $r=0$, es decir, $d\mid n$.
\end{proof}

\begin{prop}[Puntos cerrados de $\mathbb P_{\mathbb F_q}^1$]
\label{prop:s08-puntos-cerrados-campo-finito}
Sea $q$ una potencia de un primo, y sea $M_d$ el número de puntos cerrados de $\mathbb P_{\mathbb F_q}^1$ de grado $d$, donde el grado de $x$ es $[\kappa(x):\mathbb F_q]$. Para todo $n\geq1$,
\[
\sum_{d\mid n}dM_d=q^n+1.
\]
En particular, $M_1=q+1$ y $M_2=(q^2-q)/2$.
\end{prop}
\begin{proof}
Sea $N_d$ el número de polinomios mónicos irreducibles de grado $d$ en $\mathbb F_q[t]$. Por el \cref{lema:s08-factorizacion-frobenius}, comparar grados da
\[
\sum_{d\mid n}dN_d=q^n.
\]
Los puntos cerrados de grado $d$ de $\mathbb A_{\mathbb F_q}^1$ son esos $N_d$, por la \cref{prop:s08-puntos-recta-afin-campo}. El infinito añade un punto de grado uno, de modo que $M_1=N_1+1$ y $M_d=N_d$ para $d\geq2$. Así se obtiene la fórmula. Para $n=1$ da $M_1=q+1$, y para $n=2$ da $M_1+2M_2=q^2+1$, de donde resulta el valor de $M_2$. \emph{[Lang, Algebra, cap.~V, \S5.]}
\end{proof}

\begin{obs}[Lectura geométrica]
\label{obs:s08-orbitas-frobenius}
Un punto cerrado de grado $d$ es una órbita de Frobenius de $d$ puntos de $\mathbb P^1(\overline{\mathbb F}_q)$, como $\set{z,\bar z}$ es una órbita de conjugación sobre $\mathbb R$. La órbita cae en $\mathbb P^1(\mathbb F_{q^n})$ cuando $d\mid n$ y aporta $d$ puntos. Como este conjunto tiene $q^n+1$ puntos $[a:b]$, la suma cuenta las órbitas por su tamaño.
\end{obs}

\subsection{Las fibras de \texorpdfstring{$\mathbb P_A^1$}{P1A} sobre un anillo cualquiera}

\begin{prop}[Fibras de $\mathbb P_A^1$]
\label{prop:s08-fibras-proyectiva}
Sean $A$ un anillo, $\mathfrak p\in\Spec A$, $\kappa=\kappa(\mathfrak p)$ y $\pi:\mathbb P_A^1\to\Spec A$ el morfismo estructural. Hay una biyección entre el conjunto $\pi^{-1}(\mathfrak p)$ y el conjunto de puntos de $\mathbb P_\kappa^1$, que lleva:
\begin{itemize}
\item el primo $\mathfrak pA[t]$ de $U_0$ al genérico $(0)$ de $\mathbb P_\kappa^1$;
\item cada primo de $U_0$ sobre $\mathfrak p$ correspondiente a $(\bar f)\subseteq\kappa[t]$, con $\bar f$ mónico irreducible, a ese punto $(\bar f)$;
\item el único punto $\infty_{\mathfrak p}$ de $\psi_1(V(s))$ sobre $\mathfrak p$ al punto $\infty$.
\end{itemize}
\end{prop}
\begin{proof}
En $U_0$, la fibra se identifica con $\Spec\kappa[t]$ por el \cref{lema:s08-fibra-recta-afin}. El núcleo de $A[t]\to\kappa[t]$ es $\mathfrak pA[t]$, de modo que este primo corresponde a $(0)$. Por la \cref{prop:s08-proyectiva-como-conjunto}, el complemento de $U_0$ es una copia de $\Spec A$ y aporta exactamente $\infty_{\mathfrak p}$. Este corresponde en $U_1$ a $(\mathfrak p,s)$. Las identificaciones de las dos cartas respetan $s=t^{-1}$, por lo que recuperan la lista de la \cref{prop:s08-puntos-proyectiva-cuerpo} sobre $\kappa$. Solo afirmamos aquí una biyección de conjuntos; la construcción de la fibra como esquema queda para más adelante.
\end{proof}

\subsection{Sobre \texorpdfstring{$\mathbb Z$}{Z}}

\begin{prop}[{Los primos de $\mathbb Z[t]$}]
\label{prop:s08-primos-z-t}
Los ideales primos de $\mathbb Z[t]$ son exactamente los siguientes:
\begin{enumerate}[label=(\alph*)]
\item $(0)$;
\item $(g)$, con $g\in\mathbb Z[t]$ primitivo, irreducible y de grado positivo, único salvo signo;
\item $(p)$, con $p$ un número primo;
\item $(p,g)$, con $p$ primo y $g\in\mathbb Z[t]$ mónico cuya reducción $\bar g\in\mathbb F_p[t]$ es irreducible.
\end{enumerate}
Los dos primeros tipos contraen a $(0)$ en $\mathbb Z$ y los dos últimos a $(p)$. En el último, el ideal depende solo de $\bar g$, no del levantamiento mónico $g$. \emph{[Vakil, ejercicio 3.2.Q; Sesión 1.]}
\end{prop}
\begin{proof}
Sea $\mathfrak q\subseteq\mathbb Z[t]$ primo. Su contracción a $\mathbb Z$ es $(0)$ o $(p)$.

Si $\mathfrak q\cap\mathbb Z=(0)$, localizamos en $\mathbb Z\setminus\set{0}$. Por las \cref{prop:s08-primos-localizacion,prop:s08-polinomios-campo}, $\mathfrak q$ es la contracción de $(0)$ o de $(f)\subseteq\mathbb Q[t]$, con $f$ irreducible. En el primer caso es $(0)$. En el segundo, la \cref{prop:s08-gauss}(G1) permite sustituir $f$ por un primitivo irreducible $g\in\mathbb Z[t]$, único salvo signo. Por (G2),
\[
\mathfrak q=g\mathbb Q[t]\cap\mathbb Z[t]=g\mathbb Z[t].
\]
Estos ideales son primos y no contienen enteros no nulos.

Si $\mathfrak q\cap\mathbb Z=(p)$, el primo contiene a $(p)$. La \cref{prop:s08-primos-cociente} lo identifica con un primo de $\mathbb F_p[t]$: es $(0)$ o $(\bar g)$, con $\bar g$ mónico irreducible. Al contraer obtenemos $(p)$ o $(p,g)$ para cualquier levantamiento mónico $g$. Dos levantamientos difieren por un elemento de $p\mathbb Z[t]$, de modo que dan el mismo ideal. La unicidad del generador mónico en $\mathbb F_p[t]$ y las correspondencias usadas prueban que no hay más identificaciones.
\end{proof}

\begin{samepage}
\begin{prop}[Los puntos de $\mathbb P_{\mathbb Z}^1$]
\label{prop:s08-puntos-sobre-z}
\label{ejem:s08-puntos-sobre-z}
Los puntos de $\mathbb P_{\mathbb Z}^1$ son los primos de $U_0=\Spec\mathbb Z[t]$ de la \cref{prop:s08-primos-z-t} y los puntos de la sección del infinito $\psi_1(V(s))\cong\Spec\mathbb Z$: $(s)$ y los $(s,p)$, con $p$ primo. Sus cuerpos residuales y las imágenes por $\pi$ son:
\begin{center}
\normalfont
\small
\renewcommand{\arraystretch}{1.3}
\begin{tabular}{lccc}
\hline
Punto & Ideal & Cuerpo residual & Sobre\\
\hline
genérico & $(0)$ & $\mathbb Q(t)$ & $(0)$\\
cerrado en la fibra racional & $(g)$ & $\mathbb Q[t]/(g)$ & $(0)$\\
genérico de la fibra en $p$ & $(p)$ & $\mathbb F_p(t)$ & $(p)$\\
cerrado en la fibra en $p$ & $(p,g)$ & $\mathbb F_p[t]/(\bar g)$ & $(p)$\\
infinito, fibra racional & $(s)$ & $\mathbb Q$ & $(0)$\\
infinito, fibra en $p$ & $(s,p)$ & $\mathbb F_p$ & $(p)$\\
\hline
\end{tabular}
\end{center}
Las cuatro primeras filas son primos de $\mathbb Z[t]$ en $U_0$; las dos últimas son primos de $\mathbb Z[s]$ en $U_1$.
\end{prop}
\end{samepage}
\begin{proof}
La lista se obtiene de las \cref{prop:s08-primos-z-t,prop:s08-proyectiva-como-conjunto}. El cuerpo residual de un primo $\mathfrak q$ es $\Frac(B/\mathfrak q)$ en su carta $\Spec B$. Para $(g)$ se puede localizar primero en los enteros no nulos y obtener el campo $\mathbb Q[t]/(g)$. Para $(p)$ se obtiene $\mathbb F_p(t)$ y para $(p,g)$ el cociente $\mathbb F_p[t]/(\bar g)$, que ya es un campo. Los cocientes por $(s)$ y $(s,p)$ son $\mathbb Z$ y $\mathbb F_p$. Finalmente, el morfismo estructural contrae cada primo a $\mathbb Z$, lo que da la última columna.
\end{proof}

% Presentación S08, diapositiva 26: superficie aritmética y secciones.
\begin{center}
\begin{tikzpicture}[every node/.style={font=\small},x=1.05cm,y=.85cm]
\draw[grisCIMAT,->] (-.4,0)--(9.4,0) node[right] {$\Spec\mathbb Z$};
\foreach \x/\p in {2/2,4/3,6/5,8/7}{
 \draw[green!45!black,thick] (\x,.5)--(\x,3.7);
 \fill[grisCIMAT] (\x,0) circle (2pt) node[below] {$(\p)$};
}
\draw[blue!60!black,thick] plot[domain=.5:3.7,samples=75] ({.05*sin(700*\x)},\x);
\node[below] at (0,0) {$(0)$};
\node[above right,blue!60!black] at (0,3.7) {$\mathbb P_{\mathbb Q}^1$};
\node[above,green!45!black] at (2,3.7) {$\mathbb P_{\mathbb F_2}^1$};
\draw[orange!70!black,thick] (-.35,3.35)--(9,3.35) node[right] {$s=0$};
\draw[rojoCIMAT,thick] (-.35,1.35)--(9,1.35) node[right] {$t=0$};
\draw[orange!70!black,thick] (-.35,2.35) .. controls (.8,2.1) and (1.3,1.7) .. (2,1.35)
 .. controls (4,.6) and (6,.55) .. (9,.55) node[right] {$t=2$};
\fill[rojoCIMAT] (2,1.35) circle (3pt) node[above right] {$(2,t)$};
\node at (8.8,.65) {$\cdots$};
\draw[->,blue!60!black,thick] (-1,2.25)--(-1,.3) node[midway,left] {$\pi$};
\end{tikzpicture}
\end{center}
La base $\Spec\mathbb Z$ es una curva y cada fibra es una recta proyectiva: el dibujo representa $\mathbb P_{\mathbb Z}^1$ como superficie aritmética. Las secciones $t=0$ y $t=2$ se cortan solo sobre $(2)$, en el punto $(2,t)$.

\begin{prop}[Puntos cerrados de $\mathbb P_{\mathbb Z}^1$]
\label{prop:s08-puntos-cerrados-sobre-z}
\label{obs:s08-puntos-cerrados-sobre-z}
Un punto de $\mathbb P_{\mathbb Z}^1$ es cerrado si y solo si es de tipo $(p,g)$ o $(s,p)$, es decir, si cae sobre $(p)$ y es cerrado en la fibra $\mathbb P_{\mathbb F_p}^1$ de la \cref{prop:s08-fibras-proyectiva}. Su cuerpo residual es $\mathbb F_{p^d}$, donde $d$ es su grado en esa fibra. Para cada primo $p$ hay $p+1$ puntos cerrados con cuerpo residual $\mathbb F_p$ y $(p^2-p)/2$ con cuerpo residual $\mathbb F_{p^2}$.
\end{prop}
\begin{proof}
Los ideales $(p,g)$ son máximos porque sus cocientes son los campos finitos $\mathbb F_p[t]/(\bar g)$. Si el punto está en ambas cartas, el cambio $s=t^{-1}$ identifica también su cuerpo residual con un cociente de $\mathbb F_p[s]$; este es un campo, pues la inversa de un elemento no nulo de un campo finito es un polinomio en él. Así es cerrado en ambas cartas. El punto $(s,p)$ solo está en $U_1$ y es máximo allí. Aplicamos el \cref{lema:s08-cerrados-cartas}.

Los primos $(0)$ y $(p)$ no son máximos en $\mathbb Z[t]$. Tampoco lo es $(g)$: elijamos un primo $p$ que no divida al coeficiente principal de $g$. Entonces $\bar g$ es no constante y tiene un factor irreducible $\bar h\in\mathbb F_p[t]$. El ideal propio $(p,h)$ contiene estrictamente a $(g)$, pues $p\notin(g)$. Por último, $(s)$ no es máximo en $\mathbb Z[s]$. Esto excluye todos los demás puntos de la \cref{prop:s08-puntos-sobre-z}. Los grados y el conteo se deducen de las \cref{prop:s08-fibras-proyectiva,prop:s08-puntos-cerrados-campo-finito}.
\end{proof}

\begin{obs}[Dimensión]
\label{obs:s08-dimension-proyectiva-z}
La cadena de primos $(0)\subsetneq(p)\subsetneq(p,t)$ en $\mathbb Z[t]$ tiene longitud dos. Por tanto, la dimensión de Krull de $\mathbb Z[t]$ y la de $\mathbb P_{\mathbb Z}^1$, que contiene $\Spec\mathbb Z[t]$ como abierto, son al menos dos. Esto motiva la imagen de superficie; la dimensión se estudiará en la Sesión 16.
\end{obs}

\emph{Advertencia.} Un punto $(f)$ de la fibra racional es cerrado en $\mathbb P_{\mathbb Q}^1$, pero no en $\mathbb P_{\mathbb Z}^1$. Por ejemplo,
\[
\mathbb Z[t]/(2t-1)\cong\mathbb Z[1/2]
\]
no es un cuerpo. La cerradura de $(2t-1)$ es una curva horizontal que corta cada fibra una vez. En $U_0$ es $V(2t-1)$ y sus puntos corresponden, por la \cref{prop:s08-primos-cociente}, a los primos de $\mathbb Z[1/2]$: son el propio $(2t-1)$ y los $(p,2t-1)$ con $p$ impar. En la fibra sobre $2$ no hay punto en esta carta, porque $(2,2t-1)=(1)$. En $U_1$, la ecuación es $s=2$, de modo que la cerradura llega al infinito $(s,2)$.

\begin{samepage}
\subsection{Sobre un anillo de valoración discreta}

\begin{defi}[Anillo de valoración discreta]
\label{defi:s08-dvr}
Un anillo de valoración discreta (DVR) es un dominio de ideales principales local que no es un cuerpo. Sean $K$ su cuerpo de fracciones, $\mathfrak m=(\varpi)$ su único ideal máximo y $\kappa=A/\mathfrak m$ su cuerpo residual. El generador $\varpi$ se llama uniformizante.

Ejemplos son $A=\mathbb Z_{(p)}=\set{a/b\in\mathbb Q\given p\nmid b}$, con $K=\mathbb Q$, $\mathfrak m=(p)$ y $\kappa=\mathbb F_p$, y $A=k[u]_{(u)}$, con $K=k(u)$, $\varpi=u$ y $\kappa=k$. En $\mathbb Z_{(p)}$, las no unidades son los $a/b$ con $p\mid a$, así que su único ideal máximo es $p\mathbb Z_{(p)}$. Es un dominio de ideales principales por ser una localización de $\mathbb Z$.
\end{defi}
\end{samepage}

\begin{prop}[Estructura de un DVR]
\label{prop:s08-estructura-dvr}
Sea $A$ un DVR con uniformizante $\varpi$.
\begin{enumerate}[label=(\alph*)]
\item $\Spec A=\set{(0),\mathfrak m}$, $D(\varpi)=\set{(0)}$ y $\set{\mathfrak m}$ es cerrado.
\item Cada $x\in K^\times$ se escribe de manera única como $x=u\varpi^n$, con $u\in A^\times$ y $n\in\mathbb Z$. Definimos $v(x)=n$ y $v(0)=\infty$.
\end{enumerate}
\emph{[Atiyah--Macdonald, cap.~9; Sesión 3.]}
\end{prop}
\begin{proof}
Todo primo no nulo de un dominio de ideales principales es máximo, por la \cref{prop:s08-polinomios-campo}; como $A$ es local, debe ser $\mathfrak m$. Esto prueba (a). Además, un dominio de ideales principales tiene factorización única y todo irreducible de $A$ genera el único primo no nulo, por lo que es asociado a $\varpi$. Así, todo elemento no nulo de $A$ es una unidad por una potencia no negativa de $\varpi$. Escribiendo un elemento de $K^\times$ como cociente obtenemos (b). La unicidad de la factorización da la del exponente y, fijado este, la de la unidad.
\end{proof}

\begin{prop}[Las fibras de $\mathbb P_A^1$ sobre un DVR]
\label{prop:s08-fibras-dvr}
\label{ejem:s08-proyectiva-dvr}
Sea $A$ un DVR y $\pi:\mathbb P_A^1\to\Spec A$ el morfismo estructural.
\begin{enumerate}[label=(\alph*)]
\item La fibra genérica $\pi^{-1}((0))$ es abierta y, como conjunto, es $\mathbb P_K^1$.
\item La fibra especial $\pi^{-1}(\mathfrak m)$ es cerrada y, como conjunto, es $\mathbb P_\kappa^1$.
\item El punto $(0)$ de $A[t]$ es el genérico de $\mathbb P_A^1$, y la fibra genérica es densa.
\end{enumerate}
\end{prop}
\begin{proof}
La \cref{prop:s08-fibras-proyectiva} da las dos descripciones de puntos. La \cref{prop:s08-estructura-dvr}(a) y la continuidad de $\pi$ muestran que la primera fibra es abierta y la segunda cerrada. Como $A[t]$ es un dominio, $(0)$ es genérico en $U_0$; pertenece a $D(t)$ y se identifica con el genérico $(0)$ de $U_1$. Su cerradura contiene ambas cartas y es todo $\mathbb P_A^1$; como el punto pertenece a la fibra genérica, esta es densa.
\end{proof}

\begin{obs}[El disco]
\label{obs:s08-disco-dvr}
Se piensa $\Spec A$ como un disco pequeño. El punto cerrado $\mathfrak m$ es el centro, y el genérico $(0)$ representa al disco sin su centro. Así, $\mathbb P_A^1$ es una familia de rectas proyectivas sobre el disco: la fibra especial vive sobre el centro y la genérica sobre el resto.
\end{obs}

\begin{defi}[Puntos racionales y reducción]
\label{defi:s08-reduccion}
Para un campo $F$, $\mathbb P^1(F)$ es el conjunto de clases $[a:b]$ de pares $(a,b)\in F^2\setminus\set{(0,0)}$ módulo $F^\times$. El punto asociado de $\mathbb P_F^1$ es
\[
x_{[a:b]}=
\begin{cases}
(t-b/a),&a\ne0,\\
\infty,&a=0.
\end{cases}
\]
Si $b\ne0$, también tiene coordenada $s=a/b$ en $U_1$. Por la \cref{prop:s08-puntos-proyectiva-cuerpo}, esta aplicación es una biyección sobre los puntos de grado uno.

Sea ahora $A$ un DVR. Un representante $(a,b)$ de un punto de $\mathbb P^1(K)$ es normalizado si $a,b\in A$ y al menos uno es una unidad, es decir, si $\min(v(a),v(b))=0$. Su reducción es $[\bar a:\bar b]\in\mathbb P^1(\kappa)$, donde las barras indican reducción módulo $\mathfrak m$.
\end{defi}

\begin{lema}[La reducción está bien definida]
\label{lema:s08-reduccion-bien-definida}
\label{obs:s08-reduccion-bien-definida}
Todo punto de $\mathbb P^1(K)$ tiene un representante normalizado, y dos representantes normalizados difieren por una unidad de $A$. En consecuencia, $[\bar a:\bar b]$ es un punto de $\mathbb P^1(\kappa)$ independiente del representante.
\end{lema}
\begin{proof}
Sea $r=\min(v(a),v(b))$, que es finito porque las coordenadas no son ambas cero. Multiplicar por $\varpi^{-r}$ hace que ambas valoraciones sean no negativas y una sea cero. Por la \cref{prop:s08-estructura-dvr}(b), el nuevo par está normalizado.

Si $(a',b')=\lambda(a,b)$ y ambos pares están normalizados, entonces
\[
0=\min(v(a'),v(b'))=v(\lambda)+\min(v(a),v(b))=v(\lambda).
\]
Así $\lambda\in A^\times$ y $\bar\lambda\in\kappa^\times$. Los pares reducidos son no nulos y difieren por $\bar\lambda$, por lo que representan el mismo punto.
\end{proof}

Para $A=\mathbb Z_{(p)}$, esta es la reducción de un punto racional módulo $p$. Que todo punto $K$-racional se extienda a un punto con valores en $A$ es la parte de existencia del criterio valuativo de propiedad de la Sesión 15.

% S08b, diapositiva 33 y esqueleto, sección 5.3: extensión sobre el DVR.
\begin{center}
\begin{tikzcd}[column sep=large,row sep=large]
\Spec K\arrow[r,"{[a:b]}"]\arrow[d]
 &\mathbb P_A^1\arrow[d,"\pi"]\\
\Spec A\arrow[r,swap,"\id"]\arrow[ur,dashed]&\Spec A
\end{tikzcd}
\end{center}
La flecha discontinua representa la extensión a $\Spec A$ del punto definido sobre $K$; su valor en el punto cerrado es la reducción.

\begin{prop}[La cerradura corta la fibra especial en la reducción]
\label{prop:s08-cerradura-reduccion}
Sea $[a:b]\in\mathbb P^1(K)$ y $x=x_{[a:b]}$ el punto correspondiente de la fibra genérica de $\mathbb P_A^1$. La cerradura de $\set{x}$ corta la fibra especial en exactamente un punto: el que la \cref{prop:s08-fibras-proyectiva} identifica con $x_{[\bar a:\bar b]}\in\mathbb P_\kappa^1$, para un representante normalizado $(a,b)$.
\end{prop}
\begin{proof}
Normalizamos las coordenadas mediante el \cref{lema:s08-reduccion-bien-definida}. Supongamos primero que $a$ es una unidad y pongamos $c=b/a\in A$. En $U_0$, el punto $x$ es el primo $(t-c)$ de $A[t]$: el cociente es $A$ y su contracción a $A$ es $(0)$. Su cerradura en $U_0$ es $V(t-c)\cong\Spec A$ y corta la fibra especial en el único primo
\[
\mathfrak mA[t]+(t-c),
\]
que corresponde al punto $t=\bar c=\bar b/\bar a$.

No aparece otro punto en la segunda carta. Si $b=0$, entonces $x\notin U_1$ y la \cref{obs:s08-especializacion-local} muestra que su cerradura tampoco corta a $U_1$. Si $b\ne0$, entonces $x\in U_1$ y la ecuación $bs-a=0$ se cumple en $x$, luego también en su cerradura dentro de esa carta. Cuando $b\in\mathfrak m$, su reducción es $-\bar a=0$, imposible, así que no hay punto de la fibra especial en $U_1$. Si $b$ es unidad, la cerradura es $V(s-a/b)$; su punto especial $s=\bar a/\bar b$ es el ya encontrado en el traslape.

Si $a$ no es unidad, $b$ sí lo es. El mismo argumento, intercambiando las cartas, da un único punto especial con $s=\bar a/\bar b$. En todos los casos es el punto $[\bar a:\bar b]$.
\end{proof}

% Presentación S08, diapositiva 30: especialización de un punto racional.
\begin{center}
\begin{minipage}{.45\textwidth}
\centering
\begin{tikzpicture}[every node/.style={font=\small}]
\draw[grisCIMAT] (-.7,0)--(3.3,0);
\draw[blue!60!black,thick] plot[domain=.35:2.9,samples=65] ({.05*sin(700*\x)},\x);
\draw[green!45!black,thick] (2.3,.35)--(2.3,2.9);
\node[above,blue!60!black] at (0,2.9) {$\mathbb P_K^1$};
\node[above,green!45!black] at (2.3,2.9) {$\mathbb P_\kappa^1$};
\fill[blue!60!black] (0,0) circle (2pt) node[below] {$(0)$};
\fill[green!45!black] (2.3,0) circle (2pt) node[below] {$\mathfrak m$};
\draw[rojoCIMAT,thick] (0,1.85) .. controls (.8,2.2) and (1.7,2.15) .. (2.3,1.6);
\fill[rojoCIMAT] (0,1.85) circle (2.5pt) node[left] {$[a:b]$};
\fill[rojoCIMAT] (2.3,1.6) circle (2.5pt) node[right] {$[\bar a:\bar b]$};
\end{tikzpicture}
\end{minipage}\hfill
\begin{minipage}{.51\textwidth}
\small
La cerradura de un punto $K$-racional corta la fibra especial en su reducción. Si $A=k[u]_{(u)}$, se interpreta como el límite cuando $u\to0$. Para $A=\mathbb Z_{(p)}$, con $p$ impar, $1/2=[2:1]$ se reduce a $t=2^{-1}$ en $\mathbb F_p$.
\end{minipage}
\end{center}
\emph{Advertencia.} Para $A=\mathbb Z_{(2)}$, el punto $1/2=[2:1]$ se reduce a $[0:1]=\infty$. Hay que usar un representante con una coordenada unidad: $[1:1/2]$ no sirve.

\section{\texorpdfstring{$\mathbb P_A^n$}{PnA} por cartas}

\begin{defi}[Las cartas de $\mathbb P_A^n$]
\label{defi:s08-cartas-proyectiva-n}
Sea $n\geq1$. Para $i=0,\ldots,n$, sea $A_i=A[x_{k/i}:k\ne i]$, el anillo de polinomios en las $n$ variables $x_{k/i}$, y $U_i=\Spec A_i$. Se pone $x_{i/i}=1$, $U_{ii}=U_i$ y $\varphi_{ii}=\id$. Para $j\ne i$, sea $U_{ij}=D(x_{j/i})\subseteq U_i$, con anillo de funciones $A_i[x_{j/i}^{-1}]$, y
\[
\varphi_{ij}^\sharp:A_j[x_{i/j}^{-1}]\longrightarrow A_i[x_{j/i}^{-1}],
\qquad x_{k/j}\longmapsto\frac{x_{k/i}}{x_{j/i}}.
\]
Es un isomorfismo: envía $x_{i/j}$ a $x_{j/i}^{-1}$, con inversa $x_{k/i}\mapsto x_{k/j}/x_{i/j}$. \emph{[Vakil, \S4.4.9 y \S4.4.10.]}
\end{defi}

\begin{prop}
\label{prop:s08-cociclo-proyectiva-n}
Los $U_i$, $U_{ij}$ y $\varphi_{ij}$ son datos de pegado, incluida la igualdad de dominios de (P3):
\[
\varphi_{ij}(U_{ij}\cap U_{ik})=U_{ji}\cap U_{jk}.
\]
El esquema del \cref{teo:s08-lema-pegado} es el espacio proyectivo $\mathbb P_A^n$.
\end{prop}

\emph{Advertencia.} (P3) exige dos cosas: que $\varphi_{ij}$ lleve exactamente $U_{ij}\cap U_{ik}$ sobre $U_{ji}\cap U_{jk}$, y que el triángulo de transiciones conmute. La primera no se sigue de la segunda.

% Presentación S08, diapositiva 32: dominios y triángulo del cociclo.
\begin{center}
\begin{tikzcd}[column sep=small,row sep=large]
& U_{ij}\cap U_{ik}\subseteq U_i
  \arrow[dl,swap,"\varphi_{ij}"] \arrow[dr,"\varphi_{ik}"] &\\
U_{ji}\cap U_{jk}\subseteq U_j \arrow[rr,swap,"\varphi_{jk}"]
&& U_{ki}\cap U_{kj}\subseteq U_k
\end{tikzcd}
\end{center}
\begin{proof}
Sea
\[
L=A[x_0^{\pm1},\ldots,x_n^{\pm1}],\qquad
L_0=\bigoplus_{\substack{m\in\mathbb Z^{n+1}\\ \sum_\ell m_\ell=0}}Ax^m,
\quad x^m=x_0^{m_0}\cdots x_n^{m_n}.
\]
El submódulo $L_0$ es un subanillo. Para cada $i$, definimos
\[
\iota_i:A_i\hookrightarrow L_0,\qquad x_{k/i}\longmapsto x_kx_i^{-1}.
\]
Es inyectivo porque monomios distintos se envían a monomios distintos. La imagen de $A_i[x_{j/i}^{-1}]$ es
\[
R_{ij}=\bigoplus_{\substack{\sum_\ell m_\ell=0\\m_\ell\geq0\text{ si }\ell\notin\set{i,j}}}Ax^m.
\]
En efecto, los exponentes fuera de $i,j$ son no negativos; el exponente en $j$ puede ser cualquiera, y el grado total cero fija el exponente en $i$. La descripción es simétrica en $i,j$. Además,
\[
\iota_j(x_{k/j})=\iota_i(x_{k/i}/x_{j/i}),
\qquad \varphi_{ij}^\sharp=\iota_i^{-1}\iota_j
\]
en las localizaciones identificadas con $R_{ij}$. Así se verifican (P1) y (P2).

El anillo del traslape triple en la carta $i$ se identifica con
\[
R_{ijk}=\bigoplus_{\substack{\sum_\ell m_\ell=0\\m_\ell\geq0\text{ si }\ell\notin\set{i,j,k}}}Ax^m,
\]
al invertir $x_{j/i}$ y $x_{k/i}$. En la carta $j$, los abiertos $U_{ji}\cap U_{jk}$ tienen el mismo anillo: bajo $\varphi_{ij}^\sharp$, invertir $x_{k/j}$ equivale a invertir $x_{k/i}/x_{j/i}$, y $x_{j/i}$ ya es unidad. Por tanto la transición lleva exactamente el primer abierto triple al segundo. Las tres transiciones inducen la identidad de $R_{ijk}$; su composición coincide con la transición directa. Esto prueba ambas partes de (P3).
\end{proof}

\begin{obs}
\label{obs:s08-hiperplano-infinito}
Geométricamente, $U_i$ es el abierto donde $x_i\ne0$. La carta $U_0$ es isomorfa a $\mathbb A_A^n$. Su complemento es el cerrado donde se anula $x_0$, descrito en $U_i$ por $V(x_{0/i})$. Es una copia de $\mathbb P_A^{n-1}$, el hiperplano al infinito de la Sesión 10. Para $n=1$ se recupera la \cref{defi:s08-recta-proyectiva}, con $t=x_{1/0}$ y $s=x_{0/1}$. La construcción formal mediante $\operatorname{Proj}$ llegará en la Sesión 21.
\end{obs}

\section{Los pegados de dos rectas a lo largo de \texorpdfstring{$D(x)$}{D(x)}}
\label{sec:s08-dos-rectas}

\begin{lema}[Funtorialidad del pegado]
\label{lema:s08-funtorialidad-pegado}
Sean $(X_i,X_{ij},\varphi_{ij})$ y $(X_i',X_{ij}',\varphi_{ij}')$ datos de pegado con los mismos índices, y pegados $(X,\psi_i)$ y $(X',\psi_i')$. Sean $\nu_i:X_i\to X_i'$ isomorfismos tales que $\nu_i(X_{ij})=X_{ij}'$ y $\nu_j\circ\varphi_{ij}=\varphi_{ij}'\circ\nu_i$ en $X_{ij}$. Existe un único isomorfismo $h:X\to X'$ con $h\circ\psi_i=\psi_i'\circ\nu_i$.
\end{lema}

% Presentación S08, diapositiva 34: los dos cuadrados de funtorialidad.
\begin{center}
\begin{minipage}{.43\textwidth}
\centering
\begin{tikzcd}
X_i\arrow[r,"\psi_i"]\arrow[d,swap,"\nu_i"]
 &X\arrow[d,"h"]\\
X_i'\arrow[r,swap,"\psi_i'"]&X'
\end{tikzcd}
\par\smallskip{\small La conclusión, para todo $i$.}
\end{minipage}\hfill
\begin{minipage}{.53\textwidth}
\centering
\begin{tikzcd}
X_{ij}\arrow[r,"\varphi_{ij}"]\arrow[d,swap,"\nu_i"]
 &X_{ji}\arrow[d,"\nu_j"]\\
X_{ij}'\arrow[r,swap,"\varphi_{ij}'"]&X_{ji}'
\end{tikzcd}
\par\smallskip{\small La hipótesis.}
\end{minipage}
\end{center}
\begin{proof}
Los morfismos $\psi_i'\circ\nu_i:X_i\to X'$ son compatibles, pues
\[
\psi_j'\circ\nu_j\circ\varphi_{ij}
=\psi_j'\circ\varphi_{ij}'\circ\nu_i
=\psi_i'\circ\nu_i.
\]
La \cref{obs:s08-morfismos-desde-pegado} da un único $h:X\to X'$ con la propiedad requerida. Los $\nu_i^{-1}$ satisfacen las mismas condiciones en sentido inverso y producen $h':X'\to X$. Sus composiciones restringen a la identidad en cada carta; como las cartas cubren, $h'h=\id_X$ y $hh'=\id_{X'}$.
\end{proof}

\begin{prop}[Pegados de dos rectas]
\label{prop:s08-pegados-dos-rectas}
Sea $k$ un campo, $X_1=\Spec k[x]$, $X_2=\Spec k[y]$, $X_{12}=D(x)$ y $X_{21}=D(y)$. Para un isomorfismo de $k$-esquemas $\varphi:X_{12}\to X_{21}$, el pegado es isomorfo a la recta con origen doble si $\varphi^\sharp(y)=cx$, o a $\mathbb P_k^1$ si $\varphi^\sharp(y)=cx^{-1}$, con $c\in k^\times$. No hay otros casos, y los dos esquemas no son isomorfos.
\end{prop}
\begin{proof}
El isomorfismo de $k$-álgebras
\[
\varphi^\sharp:k[y,y^{-1}]\xrightarrow{\sim}k[x,x^{-1}]
\]
envía $y$ a una unidad. Las unidades de $k[x,x^{-1}]$ son los monomios $cx^n$, con $c\in k^\times$ y $n\in\mathbb Z$. En un producto de polinomios de Laurent sobre un campo, los exponentes mínimo y máximo se suman. Si el producto es uno, cada factor tiene un solo exponente.

La inversa también envía $x$ a una unidad, digamos $dy^m$. Al componer,
\[
x=dc^m x^{nm},
\]
de modo que $nm=1$ y $n=\pm1$. Tomemos $\nu_1=\id$ y el cambio de coordenada en la segunda carta definido por $\nu_2^\sharp(y')=c^{-1}y$. Entonces
\[
\varphi^\sharp\nu_2^\sharp(y')=c^{-1}cx^{\pm1}=x^{\pm1}.
\]
Además, $\nu_2^\sharp$ lleva $(y')^{\pm1}$ a $c^{\mp1}y^{\pm1}$, de modo que $\nu_2(X_{21})=X_{21}'$. La compatibilidad para el par inverso se obtiene invirtiendo la del par $(1,2)$. El \cref{lema:s08-funtorialidad-pegado} identifica el pegado con el de $y'=x$ o el de $y'=x^{-1}$. Son la recta con origen doble del \cref{ejem:s07-recta-origen-doble} y la recta proyectiva de la \cref{defi:s08-recta-proyectiva}. La constante se absorbe en el cambio de carta; el exponente distingue los casos. Que no son isomorfos se demuestra en la \cref{prop:s08-secciones-origen-doble}.
\end{proof}

\begin{prop}[Secciones globales de la recta con origen doble]
\label{prop:s08-secciones-origen-doble}
Sea $X$ la recta con origen doble sobre un campo $k$, con cartas $X_1=\Spec k[x]$ y $X_2=\Spec k[y']$ pegadas por $y'=x$ en los abiertos donde la coordenada es invertible. Entonces $\Gamma(X,\mathcal O_X)=k[x]$. En consecuencia, $X$ y $\mathbb P_k^1$ no son isomorfos, y $X$ no es afín. \emph{[Vakil, ejercicio 4.4.B.]}
\end{prop}
\begin{proof}
Una sección es un par $(f(x),g(y'))$ que coincide en $k[x,x^{-1}]$ tras sustituir $y'=x$. Como $k[x]\hookrightarrow k[x,x^{-1}]$ es inyectivo, la condición equivale a $f=g$. Así, las secciones son los pares $(f,f)$ y su anillo es $k[x]$. Un isomorfismo con $\mathbb P_k^1$ induciría un isomorfismo entre $k[x]$ y $k$, por la \cref{prop:s08-secciones-proyectiva}; esto es imposible, pues $k$ es un campo y $k[x]$ no.

Las restricciones $\Gamma(X,\mathcal O_X)\to\Gamma(X_i,\mathcal O_{X_i})$ son isomorfismos. Por la construcción de la \cref{defi:s08-morfismo-canonico}, $\eta_X:X\to\Spec k[x]$ restringe a la identidad en cada carta, identificando $y'$ con $x$. Envía los dos orígenes distintos al mismo punto $(x)$, de modo que no es inyectivo. La \cref{prop:s08-propiedades-morfismo-canonico}(e) implica que $X$ no es afín.
\end{proof}

\begin{obs}[Los dos pegados sobre un anillo]
\label{obs:s08-dos-pegados-anillo}
El mismo cálculo de secciones de la recta con origen doble vale sobre cualquier anillo $A$: $A[x]\hookrightarrow A[x,x^{-1}]$ sigue siendo inyectivo y el anillo de secciones es $A[x]$. Si $A\ne0$, para cualquier $\mathfrak p\in\Spec A$ los dos puntos $\mathfrak pA[x]+(x)$, uno en cada carta, son distintos y tienen la misma imagen por el morfismo canónico. Por tanto, la recta con origen doble tampoco es afín sobre $A$. Junto con el \cref{coro:s08-proyectiva-no-afin}, esto da dos pegados no afines de las mismas piezas. Sobre un campo los distingue la \cref{prop:s08-secciones-origen-doble}; para un anillo no nulo cualquiera los distinguirá la separación: $\mathbb P_A^1$ es separado y la recta con origen doble no lo es. Esta propiedad se estudiará en la Sesión 15.
\end{obs}

\section{Ejercicios}

\begin{ejer}[El espacio afín menos el origen]
\label{ejer:s08-afin-menos-origen}
Sea $k$ un campo, $n\geq2$ y $U=\mathbb A_k^n\setminus\set{(x_1,\ldots,x_n)}$. Muestra que $\Gamma(U,\mathcal O_U)=k[x_1,\ldots,x_n]$ y que $U$ no es afín. Pista: cubre $U$ con $D(x_1),\ldots,D(x_n)$ y adapta la \cref{prop:s08-secciones-plano-menos-origen}. \emph{[Cf. Vakil, \S4.4.1 a \S4.4.3 (caso $n=2$) y \S4.4.2.]}
\end{ejer}

\begin{ejer}[La recta menos el origen]
\label{ejer:s08-recta-menos-origen}
Sea $k$ un campo y $V=\mathbb A_k^1\setminus\set{0}=D(x)$. Calcula $\Gamma(V,\mathcal O_V)$ y describe $\eta_V$. Justifica por qué es un isomorfismo, a diferencia del plano menos el origen.
\end{ejer}

\begin{ejer}[El automorfismo que intercambia las cartas]
\label{ejer:s08-intercambio-cartas}
Muestra que $t\leftrightarrow s$ define un automorfismo de $\mathbb P_A^1$ que intercambia $\psi_0(U_0)$ y $\psi_1(U_1)$. Pista: usa el \cref{lema:s08-funtorialidad-pegado} con $\nu_0:U_0\to U_1$, $\nu_1:U_1\to U_0$ y, como segundos datos de pegado, los mismos con los índices $0$ y $1$ intercambiados.
\end{ejer}

\begin{ejer}[Puntos de $\mathbb P_{\mathbb F_2}^1$ y $\mathbb P_{\mathbb F_3}^1$]
\label{ejer:s08-puntos-f2-f3}
Lista los puntos cerrados de $\mathbb P_{\mathbb F_q}^1$ de grado uno y dos para $q=2$ y $q=3$. Verifica $M_2=(q^2-q)/2$ y encuentra $M_3$ para $q=2$.
\end{ejer}

\begin{ejer}[Los puntos cerrados de $\mathbb P_{\mathbb Z}^1$ sobre $2$]
\label{ejer:s08-puntos-sobre-dos}
Describe la fibra de $\mathbb P_{\mathbb Z}^1\to\Spec\mathbb Z$ sobre $(2)$ y lista sus puntos cerrados con cuerpo residual $\mathbb F_2$ y $\mathbb F_4$.
\end{ejer}

\begin{ejer}[$\mathbb P_A^n$]
\label{ejer:s08-proyectiva-n}
Muestra que $\Gamma(\mathbb P_A^n,\mathcal O)=A$ y que $\mathbb P_A^n$ no es afín si $n\geq1$ y $A\ne0$. Pista: una sección global se restringe de forma inyectiva a $U_0$, pues $A_i\to A_i[x_{0/i}^{-1}]$ es inyectivo. Basta entonces intersecar los anillos de $U_0$ y $U_1$ dentro de $L_0$, como en la \cref{prop:s08-cociclo-proyectiva-n}. \emph{[Vakil, ejercicio 4.4.E.]}
\end{ejer}

